\documentclass[11pt]{article}
\usepackage[margin=1in]{geometry}
\usepackage[utf8]{inputenc}
\usepackage[T1]{fontenc}
\usepackage{lmodern}
\usepackage{amsmath,amssymb,amsthm,mathtools}
\usepackage{booktabs,tabularx,longtable,array}
\usepackage{xcolor}
\newcolumntype{L}[1]{>{\raggedright\arraybackslash}p{#1}}
\newcolumntype{Y}{>{\raggedright\arraybackslash}X}
\usepackage[colorlinks=true,linkcolor=blue!50!black,citecolor=green!35!black,urlcolor=blue!50!black]{hyperref}


% labelled enumerations without enumitem: \begin{enumlab}[prefix]{\counterstyle}
\newenvironment{enumlab}[2][]{\begin{enumerate}\renewcommand{\theenumi}{#1#2{enumi}}%
  \renewcommand{\labelenumi}{(\theenumi)}}{\end{enumerate}}

% ---------------------------------------------------------------- theorems
\theoremstyle{plain}
\newtheorem{theorem}{Theorem}[section]
\newtheorem{lemma}[theorem]{Lemma}
\newtheorem{proposition}[theorem]{Proposition}
\newtheorem{corollary}[theorem]{Corollary}
\newtheorem{claim}{Claim}
\theoremstyle{definition}
\newtheorem{definition}[theorem]{Definition}
\newtheorem{question}[theorem]{Question}
\theoremstyle{remark}
\newtheorem{remark}[theorem]{Remark}

\newtheoremstyle{named}{\topsep}{\topsep}{\itshape}{}{\bfseries}{.}{.5em}{\thmnote{#3}}
\theoremstyle{named}
\newtheorem*{named}{}
\newtheoremstyle{namedrm}{\topsep}{\topsep}{\normalfont}{}{\bfseries}{.}{.5em}{\thmnote{#3}}
\theoremstyle{namedrm}
\newtheorem*{namedrm}{}

\newenvironment{proofsketch}[1][]{\def\psarg{#1}%
  \ifx\psarg\empty\begin{proof}[Proof sketch]\else\begin{proof}[Proof sketch (#1)]\fi}{\end{proof}}

\newcommand{\status}[1]{\par\smallskip\noindent
  \fbox{\parbox{\dimexpr\linewidth-2\fboxsep-2\fboxrule\relax}{\small\textbf{Status.} #1}}\par\smallskip}

% ---------------------------------------------------------------- notation
\newcommand{\Sym}{\operatorname{Sym}}
\newcommand{\Alt}{\operatorname{Alt}}
\newcommand{\Aut}{\operatorname{Aut}}
\newcommand{\Iso}{\operatorname{Iso}}
\newcommand{\Aff}{\operatorname{Aff}}
\newcommand{\supp}{\operatorname{supp}}
\newcommand{\cw}{\operatorname{cw}}
\newcommand{\tw}{\operatorname{tw}}
\newcommand{\Lab}{\operatorname{Lab}}
\newcommand{\polylog}{\operatorname{polylog}}
\newcommand{\Gd}{\widehat{\Gamma}_d}
\newcommand{\Gq}[1]{\widehat{\Gamma}_{#1}}
\newcommand{\fB}{\mathfrak{B}}
\newcommand{\fD}{\mathfrak{D}}
\newcommand{\fA}{\mathfrak{A}}
\newcommand{\fF}{\mathfrak{F}}
\newcommand{\fX}{\mathfrak{X}}
\newcommand{\fP}{\mathfrak{P}}
\newcommand{\cC}{\mathcal{C}}
\newcommand{\cO}{\mathcal{O}}
\newcommand{\cS}{\mathcal{S}}
\newcommand{\cT}{\mathcal{T}}
\newcommand{\cY}{\mathcal{Y}}
\newcommand{\cZ}{\mathcal{Z}}
\newcommand{\Cp}{\mathcal{C}'}
\newcommand{\vT}{\vec{T}}
\newcommand{\tup}[2]{#1^{\langle #2\rangle}}
\newcommand{\win}[2]{#1^{[#2]}}
\newcommand{\lay}[2]{#1^{\langle\!\langle #2\rangle\!\rangle}}
\newcommand{\src}[2]{[#1, src.\,l.\,#2]}

\title{Canonical labelings for graphs of maximum degree $d$\\ in time $n^{\polylog(d)}$\\[4pt]
\large A draft for expert checking}
\author{[Author names to be added]\\[2pt]
\small Unrefereed draft, prepared with substantial AI assistance. Version of 5 November 2026.}
\date{}

\begin{document}
\maketitle

\begin{abstract}
Grohe, Neuen and Schweitzer (GNS) showed that isomorphism of graphs of maximum degree $d$ can be
decided in time $n^{O((\log d)^c)}$. Whether canonical forms can be computed within the same bound was
posed as open by Schweitzer and Wiebking (STOC 2019) and by Wiebking (ICALP 2020), and it is
Question~9.6 of Wiebking's dissertation (2021). We give an argument that the answer is positive. For
every group $G\le\Sym(\Omega)$ whose composition factors all embed in $S_d$, we compute a canonical
labeling of strings $x\colon\Omega\to\Sigma$ under $G$, or under a coset of $G$, in time
$n^{O((\log d)^c)}$ with $n=|\Omega|$ (Theorem~A). A canonical version of Luks's reduction, written out
following Babai and Luks (1983), then gives canonical forms of $n$-vertex graphs of maximum degree $d$
in time $n^{O((\log d)^c)}$ (Corollary~B). The only genuinely pairwise step of GNS is replaced by a
canonical-subcoset device: an ordered test tuple is moved to a reference tuple chosen from the group
alone, and canonization takes place under the stabilizer of the reference tuple. The use of a canonical
test set over an ordered group is due to Wiebking (dissertation, Lemma~5.34).
We also give a canonical version of Neuen's hypergraph isomorphism test for such groups
(Theorem~H; correctness proved, modulo Neuen's constructions as listed in
Section~\ref{sec:assumptions}, which include a lemma of his dissertation that we have not seen; running
time only sketched), and a proof sketch that graphs of treewidth $k$ admit canonical forms in time $n^{(\log k)^{O(1)}}$, which needs
Theorem~A but not Theorem~H (Theorem~TW). While writing out the running time we found a step in the written potential
argument of GNS (proof of Lemma~7.2) that we could not justify for the algorithm as written. We describe
a choice of block system under which the step holds, and ask whether this is the intended reading; the
GNS theorem itself is not affected.

This is an unrefereed draft, prepared with substantial AI assistance. Every result is conditional on the
published results listed in Section~\ref{sec:assumptions}, among them Babai's canonical form for strings
(STOC 2019), whose published account is a preliminary report that sketches the proof (Theorem~A needs it
only for $S_k$, $k\le d$, on relational structures of arity $O(\log d)$), and on the requirements
\ref{D1}--\ref{D4} of Section~\ref{sec:P1} that the group-theoretic subroutines we use be run
deterministically on canonical generating sets. Each statement carries its status. The draft is offered
for expert checking.
\end{abstract}

\tableofcontents

% =====================================================================
\section{Introduction}\label{sec:intro}

\subsection{Canonical forms}

A \emph{canonical form} for a class of graphs assigns to each graph $X$ a graph $\mathrm{CF}(X)$ on
$\{1,\dots,n\}$ such that $\mathrm{CF}(X)\cong X$, and such that $X\cong Y$ implies
$\mathrm{CF}(X)=\mathrm{CF}(Y)$. A \emph{canonical labeling} is a bijection $\lambda$ with
$X^\lambda=\mathrm{CF}(X)$. Canonical forms solve isomorphism. They also allow hashing and database
lookup, and they can be used recursively inside other canonization algorithms. That last use is why the
decomposition framework of Schweitzer and Wiebking~\cite{SW19,Wie21} needs canonical forms and not
just isomorphism tests.

\subsection{History}

\begin{itemize}
\item Luks~\cite{Luks82} decided isomorphism of graphs of maximum degree $d$ in time $n^{O(d)}$, by
  reducing it to string isomorphism under groups in $\Gd$ (the groups all of whose composition factors
  embed in $S_d$). Babai, Kantor and Luks~\cite{BKL83} improved this to $n^{O(d/\log d)}$.
\item Babai and Luks~\cite{BL83} and Fürer, Schnyder and Specker~\cite{FSS83} computed canonical forms
  for graphs of bounded degree by making Luks's algorithm canonical (\cite[Thm.~4.2]{BL83}: time
  $n^{O(d)}$; the remark after its proof says this can be improved to $n^{O(d/\log d)}$ with the
  techniques of Luks).
\item Babai~\cite{Bab16,Bab15} decided isomorphism in quasipolynomial time, and later computed
  canonical forms in quasipolynomial time~\cite{Bab19} (``BC'' below). BC explains that the core of
  the 2016 algorithm, the $t$-ary structure built from local certificates, is only \emph{pairwise}
  canonical, and its main new contribution is to repair this \cite[Thm.~12.1]{Bab19}.
\item Grohe, Neuen and Schweitzer~\cite{GNS} (``GNS'') decided isomorphism of graphs of maximum degree
  $d$ in time $n^{O((\log d)^c)}$. GNS reuse Babai's aggregation of local certificates, so the same
  pairwise step appears there.
\end{itemize}

\subsection{The open problems}\label{sec:open}

The problem is stated as open in several places.
\begin{itemize}
\item Schweitzer and Wiebking~\cite[Outlook]{SW19}: ``Naturally, the bounded degree isomorphism
  algorithm of [GNS] should then be amenable to canonization. However, this remains as future work.''
\item Wiebking~\cite[Outlook]{Wie20}: ``With Babai's result, it is a natural question whether the
  bounded-degree isomorphism algorithm of [GNS] extends to canonization as well. The present
  isomorphism algorithm for graphs parameterized by treewidth should then be amenable to
  canonization as well.'' The introduction of~\cite{Wie20} says that only its final algorithm (treewidth
  $n^{\polylog k}$) depends on GNS, ``for which no adequate canonization version is known''.
\item Wiebking's dissertation~\cite{Wie21} (printed page numbers):
  \begin{itemize}
  \item p.~7: there are isomorphism tests with no known canonization algorithm of the same
    asymptotic running time, ``[f]or example, \dots\ the recent bounded-degree graph isomorphism test
    given by Grohe, Neuen and Schweitzer [GNS18] and also \dots\ results building on it [Neu20a]'';
  \item pp.~9--10: the $2^{k\polylog k}n^{O(1)}$ treewidth isomorphism test of~\cite{GNSW} ``is not
    able to compute canonical forms as it involves subroutines [GNS18] that do not extend to
    canonization so far'';
  \item p.~179: the $n^{\polylog h}$ isomorphism test for graphs excluding a minor (Thm.~6.29) is
    only an isomorphism test ``since the bounded-degree techniques [GNS18, Neu20a] are not available
    for canonization so far''; the treewidth result Cor.~6.30 (p.~180) is likewise stated for
    isomorphism only;
  \item p.~181: ``A natural question is whether the isomorphism test for bounded-degree graphs and
    finally our isomorphism test for graphs excluding a minor can be turned into a canonization
    algorithm'';
  \item p.~242: ``The only isomorphism test we use for which no canonization algorithm is known so far
    is the isomorphism test for bounded-degree graphs given by Grohe, Neuen and Schweitzer [GNS18] and
    its extension to hypergraphs [Neu20a].'' This is followed by
    \textbf{Question 9.6}: ``Can canonical labelings for graphs with $n$ vertices be computed in time
    $n^{\polylog(k)}$ where $k$ is the maximum degree of $G$?'', and by \textbf{Question 9.7}, the same
    question with $k$ the Hadwiger number.
  \end{itemize}
\end{itemize}

\subsection{Results and their status}\label{sec:results}

We write $\Gd$ for the class of groups whose composition factors are all isomorphic to subgroups of
$S_d$. Precise definitions are in Section~\ref{sec:prelim}. Table~\ref{tab:status} summarizes the status
of each statement. ``Proof modulo~\S\ref{sec:assumptions}'' means that we give a complete proof, using
the published results listed in Section~\ref{sec:assumptions} without re-proving them. None of the
statements has been refereed.

\begin{named}[Theorem A (strings under $\Gd$-groups)]
There is an algorithm that, given $G\le\Sym(\Omega)$ with $G\in\Gd$ (by generators), $\sigma\in\Sym(\Omega)$
and a string $x\colon\Omega\to\Sigma$, outputs a canonical form $F_{\sigma G}(x)$ and the canonical
placement coset $\Lambda_{\sigma G}(x)=\{g\in\sigma G: x^g=F_{\sigma G}(x)\}$ (Definition~\ref{def:canonizer}
and the coset rule of \S\ref{sec:canonizers}).
The output depends only on the coset $\sigma G$ as a set and on $x$, not on the generators supplied.
The algorithm runs in time $n^{O((\log d)^c)}$, where $n=|\Omega|$ and $c$ is an absolute constant.
\end{named}

\begin{named}[Corollary B (graphs of maximum degree $d$)]
Canonical forms and canonical labelings of $n$-vertex graphs of maximum degree $d$ can be computed in
time $n^{O((\log d)^c)}$, with $c$ as in Theorem~A.
\end{named}

\begin{named}[Corollary D (labeling cosets)]
Given a graph (or a string on $V^2$) and a labeling coset $\Delta\rho$ with $\Delta\in\Gd$, a canonical
labeling coset relative to $\Delta\rho$ in the sense of Schweitzer and Wiebking can be computed in time
$|V|^{O((\log d)^c)}$, with $c$ as in Theorem~A.
\end{named}

Corollary D is the canonical counterpart of the form in which Wiebking imports GNS
\cite[Thm.~31]{Wie20}.

\begin{named}[Theorem H (canonical Neuen 2020; Section~\ref{sec:H}; running time only a proof sketch)]
Given $\Gamma\in\Gd$ on $\Omega$ ($|\Omega|=n$), a coset $\sigma\Gamma$, a partition $\fP$ with
$\fP^\sigma$ $\Gamma$-invariant, and a set $\fX$ of $m$ $\fP$-strings (in particular, a hypergraph with
$m$ edges), a canonical form of $\fX$ under $\sigma\Gamma$ and its canonical placement coset can be
computed in time $(n+m)^{O((\log d)^{c'})}$, where $c'$ is an absolute constant.
\end{named}

\begin{named}[Theorem TW (treewidth; Section~\ref{sec:TW}; proof sketch)]
Canonical forms and canonical labelings of $n$-vertex graphs of treewidth $k$ can be computed in time
$n^{O((\log k)^{c''})}$, where $c''$ is an absolute constant.
\end{named}

\begin{table}[ht]
\centering\small
\begin{tabularx}{\linewidth}{@{}L{4.6cm}Y@{}}
\toprule
statement & status\\
\midrule
Theorem A & proof modulo \S\ref{sec:assumptions} (Sections~\ref{sec:framework}--\ref{sec:time}), including the running time. Depends on Babai's canonical form \cite[Thm.~2.2]{Bab19}, whose published account is a preliminary report that sketches the proof, and on requirements \ref{D1}--\ref{D4} (deterministic subroutines on canonical generating sets). Babai's result is needed only for $S_k$, $k\le d$, on structures of arity $O(\log d)$ (Remark~\ref{rem:BC})\\
Corollary B & proof modulo Theorem A (Section~\ref{sec:luks})\\
Corollary D & proof modulo Theorem A (Section~\ref{sec:D})\\
Theorem H & correctness: proof modulo Theorem A, the constructions of~\cite{Neu20} and the other results listed in \S\ref{sec:assumptions}, including the structure-graph lemma of \cite{Neu19}, which we have not seen (algorithm CanHyp and Proposition~\ref{prop:canhyp}, \S\ref{sec:canhyp}); running time: \textbf{proof sketch}\\
Theorem TW & \textbf{proof sketch}, modulo Theorem A (via Corollary D) and canonical tools of Wiebking; it incorporates three repairs TW-1--TW-3 found in an internal review\\
H1 (coset-labeled hypergraphs) & proof modulo Theorem H (so its running time is only as good as the sketch for Theorem H)\\
H2 ($t$-CR-bounded graphs), H3 ($K_{3,h}$-minor-free, genus) & \textbf{proof sketches}\\
H4 ($K_h$-minor-free, $K_h$-topological-subgraph-free) & \textbf{conditional} on Theorem H, H1, H2, with the dynamic programming only \textbf{sketched}\\
Remark~\ref{rem:gap1} (GNS Lemma 7.2) & a question about the written proof; the proposed choice is proved to make the step valid\\
\bottomrule
\end{tabularx}
\caption{Statements and their status.}\label{tab:status}
\end{table}

\subsection{The new step and its provenance}\label{sec:device}

GNS has exactly one step that is pairwise in an essential way. It occurs in the proof of
\cite[Lemma~6.6]{GNS} (Appendix~A there), in the case where the group generated by the fullness
certificates has small support. GNS define a category whose objects are test sets of \emph{both} input
strings and whose morphisms are produced by comparing local certificates \cite[Lemma~6.3]{GNS}. The
equivalence classes on ordered $t$-tuples of $D_1\sqcup D_2$ then give the relations of the two $t$-ary
structures. Since the classes are named jointly, the structure built for $x_1$ depends on $x_2$.

We instead colour each ordered $t$-tuple $\vT$ of the non-full region $D(x)$ by the canonical form,
under the pointwise stabilizer $J_0$ of a fixed reference tuple $\vT_0$, of the certificate string of
$(T,x)$ transported by any group element that moves $\vT$ to $\vT_0$. The reference tuple is a function
of the current group alone: the first $t$ labels of the ideal domain. Lemma~T shows that the
colour is well defined and equivariant, and that it separates orbits of marked pairs.
Lemma~\ref{lem:Tprime} shows that the resulting structure is strongly canonical and has relative
symmetry defect at least $3/4$. Lemma~\ref{lem:Tpp} shows that all colours together cost at most
$k^t\cdot n$ in recursive sizes. No relation between $t$ and $n$ is needed.

\paragraph{Credit.}
\begin{itemize}
\item The device ``send a tuple to a reference tuple chosen from the group alone, then canonize under
  its stabilizer'' is a special case of the canonical-subcoset step of the Babai--Luks framework
  (Lemma~S; \cite[\S2]{BL83}, \cite[\S5]{Bab19}).
\item Its use with a \emph{canonical test set} over an ordered group is due to Wiebking
  \cite[Lemma~5.34, algorithm \textsc{ProdCert}, pp.~101--103]{Wie21}: ``In our framework, we define
  the giant representation \dots\ for a group over a linearly ordered set $V^{\mathrm{can}}$. This allows
  us to choose one single (canonical) test set $T^{\mathrm{can}}\subseteq W^{\mathrm{can}}$ only.''
  There the device produces fullness certificates and partition families for multiple labeling
  cosets, at cost $|V|^{\polylog|V|}$. It is not used there to build a non-fullness structure in the
  regime of GNS.
\item Babai, in \cite[Remark~12.5]{Bab19}, notes that a strongly canonical choice
  ``is obvious: we just pick an element (or a linear order) for one representative of each
  $G$-isomorphism class \dots\ and use $G$ to translate it to all other members of the class \dots\
  Making this selection efficient is our problem.'' BC's ``right descent'' (\S16 there) transports test
  sets in the same way.
\end{itemize}
What we believe to be new is the observation that, at a fixed reference \emph{ordered} tuple, this
device turns the non-fullness comparison classes of GNS into canonical colours at test-set size
$\Theta(\log d)$, and that the rest of GNS can be made canonical with the Babai--Luks toolkit. Together
these settle Question~9.6, if the argument is correct.

\paragraph{Why BC's repair does not transplant.} BC's replacement for the pairwise structure
\cite[Thm.~12.1]{Bab19} is trivial on asymmetric test sets \cite[Cor.~12.3]{Bab19}, so Babai needs
\cite[Claim~12.4]{Bab19}, which uses \cite[Cor.~13.3, Prop.~14.2]{Bab19}. These need test sets of size
$t\ge 2\log_2 n$. In GNS, $t=\Theta(\log d)$ and $t<k/10\le d/10$. When $d\le20\log_2 n$, test sets of size
$2\log_2 n$ violate the constraint $t<k/10$ of the GNS aggregation (as $k\le d$), and when $d<2\log_2 n$
they do not even fit into the ideal domain.

\subsection{A question about the GNS analysis}

While writing out the running time, we found one step in the potential argument of GNS
(proof of \cite[Lemma~7.2]{GNS}) that we could not justify for the algorithm as written. On an explicit
input, a choice of block system that the text permits gives a next giant degree $k'=k$, whereas the
proof asserts $k'\le k-k/(3\log k)$. A one-line choice, which is itself a run permitted by the algorithm,
makes the step valid. Remark~\ref{rem:gap1} gives the details and phrases the matter as a question.
The GNS theorem itself is not in doubt. We have not seen the SIAM journal version of GNS, which may
already treat this point.

\subsection{Organisation, conventions and caveats}

Section~\ref{sec:prelim} fixes conventions and recalls GNS. Section~\ref{sec:framework} states the
framework lemmas. Section~\ref{sec:subcoset} contains the canonical subcoset lemma and the tuple
colours. Section~\ref{sec:agg} makes the aggregation canonical. Section~\ref{sec:algo} gives the
algorithm and its correctness, and Section~\ref{sec:time} the running time. Section~\ref{sec:luks}
proves Corollary~B and Section~\ref{sec:D} Corollary~D. Section~\ref{sec:H} treats Theorem~H,
Section~\ref{sec:TW} Theorem~TW, and Section~\ref{sec:further} the remaining corollaries.
Section~\ref{sec:assumptions} lists what is assumed from the literature, Section~\ref{sec:related}
discusses related work, and Section~\ref{sec:statusopen} lists the status and the open points.

References of the form \src{GNS}{2126} point to a line of the arXiv \LaTeX{} source of the cited paper
(for GNS, the file \texttt{bounded-degree-polylog-d.tex} of arXiv:1802.04659v3). Theorem numbers for
arXiv papers were computed from their \LaTeX{} sources; we list the versions in
Section~\ref{sec:assumptions}.

\paragraph{Caveats.} This draft has not been refereed. It was prepared with substantial assistance
from an AI system. Three internal adversarial reviews were carried out, and their corrections are
incorporated; apart from the corrections listed in Section~\ref{sec:statusopen}, they found no gap in
the main argument, but they are not a substitute for refereeing.
No formal verification was attempted. Small computer checks were run only for the canonical Luks
reduction (Section~\ref{sec:luks}) on small graphs, for the example of Remark~\ref{rem:gap1}, and for the
arithmetic in Lemma~\ref{lem:potential}.

% =====================================================================
\section{Preliminaries}\label{sec:prelim}

\subsection{Conventions}

All domains $\Omega$ are finite and \emph{linearly ordered}. $\Sym(\Omega)$ acts on the right:
$\alpha^{gh}=(\alpha^g)^h$. A string is a map $x\colon\Omega\to\Sigma$, where $\Sigma$ is a finite, linearly
ordered alphabet. Groups act on strings by $x^g(\alpha)=x(\alpha^{g^{-1}})$, so $(x^g)^h=x^{gh}$. Strings on
the same domain are compared lexicographically along the order of $\Omega$. We write
$\Aut_H(x)=\{h\in H:x^h=x\}$ and $\Iso_H(x,y)=\{h\in H:x^h=y\}$. For $W\subseteq\Omega$ and
$K\le\Sym(\Omega)$ we put $\Iso^W_K(x,y)=\{g\in K: x(\alpha)=y(\alpha^g)\text{ for all }\alpha\in W\}$, as
in \cite[\S2]{GNS}, and $\Aut^W_K(x)=\Iso^W_K(x,x)$. Logarithms are to base~2.

\subsection{Canonizers}\label{sec:canonizers}

\begin{definition}\label{def:canonizer}
Let $H\le\Sym(\Omega)$ and let $\mathcal X$ be an $H$-invariant set of strings. An
\emph{$H$-canonizer on $\mathcal X$} assigns to each $x\in\mathcal X$ a string $F_H(x)$ and a set
$\Lambda_H(x)\subseteq H$ such that
\begin{enumlab}[C]{\arabic}
\item $\Lambda_H(x)=\{h\in H:x^h=F_H(x)\}\neq\emptyset$;
\item $F_H(x^h)=F_H(x)$ for all $h\in H$.
\end{enumlab}
Here $H$ is regarded as a set: a canonizer may not depend on a generating set of $H$.
\end{definition}

\begin{lemma}\label{lem:canon-basic}
For every $H$-canonizer:
\begin{enumerate}
\item $\Lambda_H(x^h)=h^{-1}\Lambda_H(x)$;
\item $\Lambda_H(x)=\Aut_H(x)\lambda=\lambda\Aut_H(F_H(x))$ for every $\lambda\in\Lambda_H(x)$;
\item $\Iso_H(x,y)=\Lambda_H(x)\Lambda_H(y)^{-1}$ if $F_H(x)=F_H(y)$, and $\Iso_H(x,y)=\emptyset$ otherwise.
\end{enumerate}
\end{lemma}
\begin{proof}
Direct from (C1) and (C2).
\end{proof}

We call $\Aut_H(F_H(x))$ the \emph{right subgroup} of the canonizer at $x$. A canonizer is represented by
one element of $\Lambda_H(x)$ together with generators of the right subgroup. Note that
$\Aut_H(x)=\lambda\Aut_H(F_H(x))\lambda^{-1}$ for $\lambda\in\Lambda_H(x)$; the right subgroup is
$\Aut_H(x)$ only when $x=F_H(x)$.

\paragraph{Cosets.} For $C=\sigma H$ put $F_C(x):=F_H(x^\sigma)$ and $\Lambda_C(x):=\sigma\Lambda_H(x^\sigma)$.
By (C2) this does not depend on $\sigma$, and $\Lambda_C(x)=\{g\in C:x^g=F_C(x)\}$. We call this the
\emph{coset rule}. This is
\cite[\S2.1]{BL83} (``$\mathrm{CF}(X,\sigma G):=\mathrm{CF}(X^\sigma,G)$'') and \cite[Def.~5.1(i)]{Bab19}.
We always use \emph{left} cosets $\sigma H$. With right cosets the group visible at a child node would be
conjugated by the input-dependent element.

\begin{remark}[Labeling cosets]\label{rem:labcosets}
Schweitzer and Wiebking work with labeling cosets $\Delta\rho\subseteq\Lab(V)$, where $\Lab(V)$ is the set
of bijections $V\to\{1,\dots,|V|\}$, $\Delta\le\Sym(V)$ and $\rho\in\Lab(V)$. Put
$H:=\rho^{-1}\Delta\rho\le\Sym([n])$; it depends only on $\Delta\rho$. Then
$\mathrm{CL}(x,\Delta\rho):=\rho\Lambda_H(x^\rho)$ is a canonical labeling relative to $\Delta\rho$ in their
sense, and conversely a canonical labeling function in their sense yields canonizers. The canonical
labelings of~\cite{BL83} are the same notion. All three formalisms are interchangeable, and we use
Definition~\ref{def:canonizer}.
\end{remark}

\subsection{The class $\Gd$}

$\Gd$ is the class of groups all of whose composition factors are isomorphic to subgroups of $S_d$
\cite[Def.~2.3]{GNS}. It is closed under subgroups and homomorphic images \cite[Lemma~2.4]{GNS}, and,
by the Jordan--Hölder theorem, under extensions. A group of degree $n$ lies in $\Gq{n}$, so we may assume
$d\le n$.

\subsection{Choices that are functions of the group (principle P1)}\label{sec:P1}

The output $F_H(x)$ must depend on $H$ only as a \emph{set}. All auxiliary choices are therefore made by
\emph{deterministic} procedures whose input is a \emph{canonical generating set} of the current group,
together with the ordered domain. A canonical generating set depends only on $H\le\Sym(\Omega)$ for
ordered $\Omega$, and is computable in polynomial time \cite[Lemma~6.2]{AGMMM18}; see also
\cite[Lemma~21]{GNSW} (arXiv numbering) and the corresponding lemma of~\cite{SW19}
(\texttt{lem:canGen}, Lemma~15 in the arXiv numbering). Concretely, one may take the lexicographically
least coset representatives along the point-stabilizer chain of the ordered domain.

We write \textbf{(P1)} for a choice made this way. Choices of this kind include block systems, Babai's
giant representation \cite[Lemma~7.1]{GNS} and the labeling of its ideal domain $[k]$, reference tuples,
lexicographically least elements of cosets, and the linear order of every derived domain
($\Omega^*$, $[k]^t$, slot domains). Groups computed along different routes, for example as right
subgroups of canonizer calls on different inputs, are re-normalized to canonical generators before any
(P1) choice.

\paragraph{Determinism requirement.} This requires that every subroutine feeding a (P1) choice be
deterministic and be run on the canonical generating set. Concretely, we require deterministic
polynomial-time versions of:
\begin{enumlab}[D]{\arabic}
\item\label{D1} the permutation-group toolkit of \cite[Thm.~2.1]{GNS} (orders, membership, orbits,
  pointwise stabilizers, block systems, kernels, preimages). Seress's book~\cite{Ser03} contains
  randomized algorithms as well as deterministic ones; the deterministic polynomial-time versions are
  the ones meant here;
\item\label{D2} the algorithm of \cite[Lemma~7.1]{GNS} (= \cite[Thm.~3.2.1]{Bab15});
\item\label{D3} the constructions in the proofs of \cite[Thms.~4.3, 4.6]{GNS}, and, for Theorem~H, the
  constructions of \cite[Thms.~4.8, 4.9, 4.17]{Neu20} and the structure-graph lemma they rely on
  \cite[Lemma~6.1.17]{Neu19};
\item\label{D4} Babai's canonical form of strings \cite[Thm.~2.2]{Bab19}. BC's canonical form takes
  ``(a set of generators of) the group $G$'' as input, and its invariance property is stated for a fixed
  input group. Theorem~A calls it only for the group $S_k$ ($k\le d$) in its action on
  $[k]\sqcup[k]^2\sqcup[k]^t$, on the fixed generating set determined by $(k,t)$
  (Proposition~\ref{prop:opt1} and Remark~\ref{rem:BC}), so its output is a function of $(k,t)$ and the
  string, whatever BC does on other generating sets. We also require the computation to be deterministic.
\end{enumlab}

\subsection{Windowed and layered strings}\label{sec:tags}

For $W\subseteq\Omega$ the \emph{windowed string} $\win{x}{W}$ carries an explicit membership tag:
$\win{x}{W}(\alpha)=(1,x(\alpha))$ for $\alpha\in W$, and $\win{x}{W}(\alpha)=(0,*)$ otherwise. Thus
$\win{x}{W}$ determines $W$, whatever $\Sigma$ is. For a chain $\emptyset=U_0\subset U_1\subset\dots\subset U_r$
the \emph{layered string} $\lay{x}{U_r}$ records the layer index: $\lay{x}{U_r}(\alpha)=(i,x(\alpha))$ for
$\alpha\in U_i\setminus U_{i-1}$, and $(0,*)$ for $\alpha\notin U_r$. All tags live in one fixed totally
ordered set $\{(0,*)\}\cup(\mathbb N\times\Sigma)$.

The tags matter for three reasons.
\begin{itemize}
\item $\Aut_H(\win{x}{W})=\{h\in H:W^h=W,\ x(\alpha^h)=x(\alpha)\text{ for }\alpha\in W\}$. With an untagged blank
  symbol that also occurs in $\Sigma$, $h$ need not stabilize $W$.
\item The tags are needed for the identity $G(T,x)=\Aut_{H_T}(\lay{x}{W(T,x)})$ used in
  Lemma~\ref{lem:Tprime}(2).
\item They are needed so that Lemma~T(c), applied to windowed strings, means ``the same
  window and the same values on it''.
\end{itemize}
Recursive calls never receive windowed strings. They receive restrictions $x|_V$ to invariant sets
(Lemma~F1), so alphabets do not grow along the recursion. Tagged strings occur only as the
colours of Section~\ref{sec:colours} and as the encodings passed to BC.

\subsection{The GNS framework: the facts we use}\label{sec:gnsrecap}

Result numbers refer to arXiv:1802.04659v3.

\paragraph{Almost $d$-ary sequences \cite[Def.~4.1, Obs.~4.2]{GNS}.} A $G$-invariant sequence of
partitions $\{\Omega\}=\fB_0\succ\fB_1\succ\dots\succ\fB_m$ (discrete) is \emph{almost $d$-ary} if for every
$i$ and every $B\in\fB_{i-1}$, either $G_B$ acts semi-regularly on $\fB_i[B]$, or $|\fB_i[B]|\le d$. Such a
sequence is inherited by subgroups and by restrictions to invariant subsets. Inserting a $G$-invariant
partition $\fP$ between $\fB_{i-1}$ and $\fB_i$ keeps the sequence almost $d$-ary provided each class of
$\fB_{i-1}$ contains at most $d$ classes of $\fP$ and each class of $\fP$ contains at most $d$ classes of
$\fB_i$ (or the corresponding actions are semi-regular); this is immediate from the definition.

\paragraph{Change of action \cite[Thms.~4.3, 4.6]{GNS}.} For transitive $G\in\Gd$ there are a set
$\Omega^*$ with $|\Omega^*|\le n^{O(\log^2 d)}$, a homomorphism $\iota\colon G\to\Sym(\Omega^*)$, $g\mapsto g^*$,
whose image carries an almost $d$-ary sequence, and a map $x\mapsto x^*$ on strings. The proofs give
\src{GNS}{1131--1178, 1207--1213, 1452--1530}:
\begin{itemize}
\item $(x^g)^*=(x^*)^{g^*}$, because $x^*((N h,\alpha))=x(\alpha)$, respectively $x^*((\alpha,\bar v))=x(\alpha)$;
\item $x\mapsto x^*$ is injective, also after restriction to one $G^*$-orbit $A^*$, because every
  $\alpha\in\Omega$ occurs as a second coordinate ($G$ is transitive).
\end{itemize}

\paragraph{Giant representations and affected points \cite[Def.~5.6]{GNS}.} A homomorphism
$\varphi\colon H\to S_k$ is a \emph{giant representation} if $\varphi(H)\ge A_k$. A point $\alpha$ is
\emph{affected} if $\varphi(H_\alpha)\not\ge A_k$; $\Aff(H,\varphi)$ denotes the set of affected points. For
$B\le H$ with $\varphi(B)\ge A_k$, $\Aff(B,\varphi)$ is $B$-invariant, and $\Aff(B^h,\varphi^h)=\Aff(B,\varphi)^h$, where $\varphi^h\colon B^h\to S_k$, $c\mapsto\varphi(hch^{-1})$. (For $h\in H$, $\varphi^h$ and $\varphi|_{B^h}$ differ by conjugation with $\varphi(h)$, so also $\Aff(B^h,\varphi)=\Aff(B,\varphi)^h$; the same holds for $\psi_T$ and $\psi_{T^{\varphi(h)}}$ below.)
$\Aff$ is antimonotone: if $B'\le B$ and both images are giants, then $\Aff(B,\varphi)\subseteq\Aff(B',\varphi)$.

\paragraph{Unaffected stabilizers \cite[Thm.~5.8]{GNS}.} If $H$ has an almost $d$-ary sequence,
$k>\max\{8,2+\log d\}$, and $D$ is the set of unaffected points, then $\varphi(H_{(D)})\ge A_k$.

\paragraph{Affected orbits \cite[Thm.~5.9]{GNS} (= \cite[Thm.~6(b)]{Bab16}).} For $k\ge5$, each orbit of
$\ker\varphi$ inside an affected $H$-orbit $\Delta$ has length at most $|\Delta|/k$.

\paragraph{Local certificates \cite[Def.~6.1, Alg.~2, Lemma~6.2]{GNS}.}
Fix a node group $H$ with a giant representation $\varphi$ and a $t$-set $T\subseteq[k]$. Put
$H_T:=\varphi^{-1}(\varphi(H)_T)$ and let $\psi_T\colon H_T\to\Sym(T)$ be the induced action. The
\emph{window chain} is: $G_0:=H_T$, $W_0:=\emptyset$; while $\psi_T(G_i)\ge\Alt(T)$ and $W_i\neq\Aff(G_i,\psi_T)$,
put $W_{i+1}:=\Aff(G_i,\psi_T)$ and $G_{i+1}:=\Aut^{W_{i+1}\setminus W_i}_{G_i}(x)$. Let $i^*$ be the final
index, $W(T,x):=W_{i^*}$ and $G(T,x):=G_{i^*}$. We use:
\begin{enumlab}[L]{\arabic}
\item\label{L1} $G_i=\Aut^{W_i}_{H_T}(x)$ (proof of \cite[Lemma~6.2]{GNS}).
\item\label{L2} $G_i$ stabilizes $W_1,\dots,W_{i+1}$ setwise, because $W_{j+1}=\Aff(G_j,\psi_T)$ is
  $G_j$-invariant and $G_i\le G_j$. Hence $G_i=\Aut_{H_T}(\lay{x}{W_i})=\Aut_{H_T}(\win{x}{W_i})$, using the
  tags of \S\ref{sec:tags}.
\item\label{L3} $T$ is \emph{full} if $\psi_T(\Aut_{H_T}(x))\ge\Alt(T)$. If $T$ is full, every $G_i$ contains
  $\Aut_{H_T}(x)$, the loop ends with $W_{i^*}=\Aff(G_{i^*},\psi_T)$, and
  $K(T,x):=(G_{i^*})_{(\Omega\setminus W)}\le\Aut_{H_T}(x)$ is a fullness certificate (by
  \cite[Thm.~5.8]{GNS}). If $T$ is not full, the window grows strictly until $\psi_T(G_i)$ is not a giant,
  since otherwise \cite[Thm.~5.8]{GNS} would make $T$ full \src{GNS}{1931--1934}. So \emph{$T$ is non-full
  if and only if $\psi_T(G(T,x))$ is not a giant.}
\item\label{L4} (Equivariance.) For $h\in H$: $W_i(T^{\varphi(h)},x^h)=W_i(T,x)^h$ and
  $G_i(T^{\varphi(h)},x^h)=G_i(T,x)^h$, by induction on $i$, using $H_{T^{\varphi(h)}}=(H_T)^h$ and equivariance
  of $\Aff$. Hence $W$, $G$ and $K$ are equivariant, and fullness is invariant.
\item\label{L5} (Cost, \cite[Lemma~6.2]{GNS}.) For one $T$: recursive calls on domains of size $\le n/2$
  whose sizes sum to $\le t!\cdot n$, plus $O(t!\,n^{c})$ further work. The large-window case uses
  \cite[Thm.~5.9]{GNS}.
\end{enumlab}

\paragraph{Symmetry defect \cite[\S6.3]{GNS}.} The symmetry defect of $A\le\Sym(\Gamma)$ is
$|\Gamma|-\max\{|\Delta|:\Alt(\Delta)\le A\}$; the relative defect divides this by $|\Gamma|$.
\cite[Cor.~6.5]{GNS} (from \cite[Thm.~5.2A,B]{DM96}): if $A_k\le S\le S_k$, $k\ge24$, $A\le S$, and the
relative symmetry defect of $A$ is at least $1/4$, then $|S:A|\ge(4/3)^k$.

\paragraph{Giant representation of a large primitive group \cite[Lemma~7.1]{GNS}
(= \cite[Thm.~3.2.1]{Bab15}).} For primitive $P\le S_b$ with $|P|\ge b^{1+\log b}$ ($b$ larger than an
absolute constant), one computes in polynomial time a normal subgroup $N\trianglelefteq P$ with
$|P:N|\le b$, an $N$-invariant equipartition $\cC$ of $[b]$, and a giant representation
$\psi\colon N\to S_k$ with $k\ge\log b$ and $\ker\psi=N_{(\cC)}$. We use two consequences:
\begin{itemize}
\item $|\cC|\ge2$, since $\ker\psi\neq N$ (as $\psi$ is a giant with $k\ge3$); as $\cC$ is an equipartition,
  each class of its lift to $\Omega$ has size $n/|\cC|\le n/2$;
\item $k\le d$ when $b=|\fB|\le d$, since $N\le P\le\Sym(\fB)$ \src{GNS}{2092}; we use this for the height
  of the recursion and for the size of the call tree.
\end{itemize}

\paragraph{The main recursion \cite[Lemma~7.2, Alg.~4]{GNS}.} For transitive $G$ with almost $d$-ary
sequence:
\begin{itemize}
\item if $G^{\fB_1}$ is semi-regular, Luks reduction;
\item otherwise $|\fB_1|\le d$: coarsen $\fB_1$ to a block system $\fB$ on which $P=G^{\fB}$ is primitive;
  if $|P|$ is small, Luks reduction; otherwise apply \cite[Lemma~7.1]{GNS}, put
  $G':=\{g:g^{\fB}\in N\}$ (index $\le d$) and $\varphi:=\psi\circ(g\mapsto g^{\fB})$, so that
  $\ker\varphi=G'_{(\Cp)}$, where $\Cp$ is the lift of $\cC$ to $\Omega$;
\item put $t:=\max\{9,3+\log d\}$; if $k\le10t$, Luks reduction to $\ker\varphi$;
\item otherwise aggregate local certificates and recurse on at most $k^6$ groups $H\le G'$ with
  $|\varphi(G'):\varphi(H)|\ge(4/3)^k$.
\end{itemize}

% =====================================================================
\section{Framework lemmas}\label{sec:framework}

All of these are standard (\cite[\S2]{BL83}; \cite[\S\S5, 8]{Bab19}). We restate and prove them in the
conventions of Section~\ref{sec:prelim}, because the left/right conventions matter.

\begin{named}[Lemma F1 (restriction)]
Let $U\subseteq\Omega$ be $H$-invariant and let $\rho_U\colon H\to H^U$ be restriction. Let $(F',\Lambda')$ be
an $H^U$-canonizer for $x|_U$. Define $F(\alpha):=(1,F'(\alpha))$ for $\alpha\in U$, $F(\alpha):=(0,*)$ otherwise,
and $\Lambda:=\rho_U^{-1}(\Lambda')$. Then $(F,\Lambda)$ is an $H$-canonizer for $\win{x}{U}$. Its right subgroup is
$\Aut_H(F)=\rho_U^{-1}(\Aut_{H^U}(F'))$, and $\Aut_H(\win{x}{U})=\lambda\Aut_H(F)\lambda^{-1}$ for every
$\lambda\in\Lambda$. Everything is computable in polynomial time from $(F',\Lambda')$ (preimage of a coset,
\cite[Thm.~2.1]{GNS}).
\end{named}
\begin{proof}
Since $U$ is $H$-invariant, $h\in H$ maps $\win{x}{U}$ to $F$ iff $\rho_U(h)$ maps $x|_U$ to $F'$. (C2) holds
because $F'$ is constant on $H^U$-orbits. The last statement is Lemma~\ref{lem:canon-basic}(2).
\end{proof}

\begin{named}[Lemma F2 (sequential chain rule with input-dependent windows)]
Suppose windows $V_{i+1}$ are given as functions of the layered string $\lay{x}{U_i}$, with
$V_{i+1}\subseteq\Omega\setminus U_i$, that are $H$-equivariant ($V_{i+1}(y^h)=V_{i+1}(y)^h$), and put
$U_{i+1}:=U_i\sqcup V_{i+1}$; suppose the number of layers $r(x)$ is $H$-invariant. Given an $H$-canonizer
$(F_i,\Lambda_i)$ for $\lay{x}{U_i}$:
\begin{itemize}
\item pick $\tau\in\Lambda_i$ and put $w:=x^\tau$;
\item put $A_i:=\Aut_H(F_i)$ and $V':=V_{i+1}(F_i)$, which is $A_i$-invariant;
\item let $(F',\Lambda')$ be an $A_i$-canonizer of $\win{w}{V'}$;
\item let $F_{i+1}$ be $F_i$ overwritten on $V'$ by $F'$ with layer tag $i+1$, and $\Lambda_{i+1}:=\tau\Lambda'$.
\end{itemize}
Then $(F_{i+1},\Lambda_{i+1})$ is an $H$-canonizer for $\lay{x}{U_{i+1}}$.
\end{named}
\begin{proof}
\emph{Independence of $\tau$.} Replacing $\tau$ by $\tau a$ ($a\in A_i$) replaces $w$ by $w^a$. Since
$V'^a=V'$, $F'$ is unchanged and $\Lambda'$ becomes $a^{-1}\Lambda'$, so $\Lambda_{i+1}$ is unchanged.
\emph{(C2).} $\Lambda_i(x^h)=h^{-1}\Lambda_i(x)$, so for $x^h$ we may use $\tau'=h^{-1}\tau$, which gives the
same $w$.
\emph{(C1), $\subseteq$.} Let $a\in\Lambda'$. Then $(\lay{x}{U_i})^{\tau a}=F_i^a=F_i$ and
$V_{i+1}(\lay{x}{U_i})^\tau=V_{i+1}(F_i)=V'$, so $(\win{x}{V_{i+1}})^{\tau a}=(\win{w}{V'})^a=F'$.
\emph{(C1), $\supseteq$.} If $h$ maps $\lay{x}{U_{i+1}}$ to $F_{i+1}$, the layer tags force
$(\lay{x}{U_i})^h=F_i$, so $h=\tau a$ with $a\in A_i$, and $(\win{w}{V'})^a=F'$, so $a\in\Lambda'$.
\end{proof}

\begin{definition}[The chain rule used in this paper]\label{def:chain}
The \emph{chain rule} is the special case of Lemma~F2 in which the windows are orbits of the current
right subgroup. At stage $i$, let $V^*_{i+1}$ be the first $A_i$-orbit, in the order (size, least point),
among the positions where $F_i$ carries the tag $(0,*)$ (that is, in $\Omega\setminus U_i^{\tau}$ for
$\tau\in\Lambda_i$), and for $y$ in the $H$-orbit of $\lay{x}{U_i}$ put
$V_{i+1}(y):=(V^*_{i+1})^{\lambda^{-1}}$ for any $\lambda\in\Lambda_i(y)$. This is well defined because $V^*_{i+1}$
is $A_i$-invariant and $\Lambda_i(y)=\lambda A_i$. It is $H$-equivariant by
Lemma~\ref{lem:canon-basic}(1), and $V_{i+1}(y)$ is an orbit of $\Aut_H(y)$. The canonizer of
$\win{w}{V'}$ under $A_i$ is obtained by Lemma~F1 from one recursive call on the \emph{transitive}
instance $(V',A_i^{V'},w|_{V'})$.
\end{definition}

Applied to an intransitive $H$, the chain rule processes ``orbit by orbit'' as in \cite[Alg.~1]{GNS}.
Its windows are orbits of the shrinking right subgroups, so they are finer than the fixed orbits used by
GNS and never cost more. We use only this chain rule. In particular, every recursive call it makes is on
a transitive instance, which is used in Sections~\ref{sec:potential} and~\ref{sec:timebound}. (Note
that ``least point'' is not $H$-invariant on the input side; this is why the window is defined in the
canonical frame and pulled back.)

\begin{named}[Lemma F3 (descent)]
Let $K\le H$ with $H=\bigsqcup_{\rho\in Q}\rho K$. Put $m:=\min_\rho F_K(x^\rho)$ and
$\Lambda_H(x):=\bigcup_{\rho:F_K(x^\rho)=m}\rho\Lambda_K(x^\rho)$. This is an $H$-canonizer with $F_H(x)=m$,
and $\Lambda_H(x)=\{h\in H:x^h=m\}$.
\end{named}
\begin{proof}
\emph{(C2).} For $x^h$ the branch values are $F_K(x^{h\rho})$. The cosets $h\rho K$ run over all left
cosets of $K$, and $F_K$ is constant on $K$-orbits, so the set of branch values is unchanged.
\emph{(C1).} Every $\rho\lambda$ in the union maps $x$ to $m$. Conversely, let $x^{\rho k}=m$. Then $m$ lies in
the $K$-orbit of $x^\rho$, and $m=F_K(x^{\rho_0})$ for a minimizer $\rho_0$, so $F_K(m)=m$. Hence
$F_K(x^\rho)=m$ and $k\in\Lambda_K(x^\rho)$.
\end{proof}
The output depends only on the branch values, so branches whose values are known to coincide need not
be evaluated separately (Lemma~\ref{lem:K} uses this). The right subgroup is generated by $\Aut_K(m)$ and
one element $\lambda_1^{-1}\lambda_j$ per further minimizing branch, where $\lambda_j\in\rho_j\Lambda_K(x^{\rho_j})$.

\begin{named}[Lemma F4 (equivariant branch families)]
Let $x\mapsto\fF(x)$ assign to each $x$ a finite set of objects with $\fF(x^h)=\fF(x)^h$, and suppose an
$H$-canonizer for pairs $(x,b)$, $b\in\fF(x)$, is given. Let $F:=\min_{b\in\fF(x)}F_H(x,b)$ (pairs compared
lexicographically), let $F^x$ be its first component, and let $\Lambda_H(x)$ be the union of
$\Lambda_H(x,b)$ over the minimizing $b$. Then $\Lambda_H(x)=\{h:x^h=F^x\}$, so this is an $H$-canonizer for
$x$ with $F_H(x)=F^x$.
\end{named}
\begin{proof}
(C2) is clear. Every element of the union maps $x$ to $F^x$. Conversely let $x^h=F^x$. Write
$F=(F^x,F^b)=(x,b_0)^{\lambda_0}$ for a minimizing $b_0$. Then $F^b\in\fF(x)^{\lambda_0}=\fF(F^x)=\fF(x)^h$, so
$b:=(F^b)^{h^{-1}}\in\fF(x)$ and $(x,b)^h=F$. Since $F_H(F)=F$, we get $F_H(x,b)=F$, so $b$ minimizes and
$h\in\Lambda_H(x,b)$.
\end{proof}

\begin{named}[Lemma R (change of action)]
Let $\iota\colon G\to\Sym(\Omega^*)$, $g\mapsto g^*$, be a homomorphism and $x\mapsto x^*$ an injective map
with $(x^g)^*=(x^*)^{g^*}$. Given a $G^*$-canonizer $(F_{G^*},\Lambda_{G^*})$, $G^*:=\iota(G)$, on the
$G^*$-invariant set $\{x^*\}$ of encoded strings, put
$\Lambda_G(x):=\iota^{-1}(\Lambda_{G^*}(x^*))$, and let $F_G(x)$ be the unique $y$ with $y^*=F_{G^*}(x^*)$. This
is a $G$-canonizer.
\end{named}
\begin{proof}
$y$ exists: $F_{G^*}(x^*)=(x^*)^{g^*}=(x^g)^*$ for some $g\in G$. (C1): $g\in\Lambda_G(x)$ iff
$(x^g)^*=F_{G^*}(x^*)=(F_G(x))^*$ iff $x^g=F_G(x)$, by injectivity. (C2) follows from (C2) for $G^*$ and
equivariance.
\end{proof}
In the application, GNS's $\Omega^*$, $G^*$, the almost $d$-ary sequence and the order on $\Omega^*$ are
computed from $G$ alone (P1; rows 1--5 of Table~\ref{tab:choices}).

% =====================================================================
\section{The canonical subcoset lemma and tuple colours}\label{sec:subcoset}

\begin{named}[Lemma S (canonical subcoset)]
Let $H\le\Sym(\Omega)$ and let $\varphi\colon H\to\Sym(\Gamma)$ be a homomorphism. Let $\cS$ be a
$\varphi(H)$-invariant set of structures on $\Gamma$ with a $\varphi(H)$-canonizer $(F',\Lambda')$. Let
$x\mapsto S(x)\in\cS$ satisfy $S(x^h)=S(x)^{\varphi(h)}$, and write $S_{\mathrm{can}}=S_{\mathrm{can}}(x):=F'(S(x))$
and $R=R(x):=\varphi^{-1}(\Aut_{\varphi(H)}(S_{\mathrm{can}}))$; both depend only on the $H$-orbit of $x$.
Suppose that for each value of $S_{\mathrm{can}}$ an $R$-canonizer is available on the $R$-invariant set of
strings $x'$ with $S(x')=S_{\mathrm{can}}$. Define $E(x):=\varphi^{-1}(\Lambda'(S(x)))$; this is a left coset
$\sigma R$, and we put $\Lambda_H(x):=\Lambda_{\sigma R}(x)$, i.e.\ $F_H(x):=F_R(x^\sigma)$. Then this is an
$H$-canonizer, with right subgroup $\Aut_R(F_H(x))=\Aut_H(F_H(x))$.
\end{named}
\begin{proof}
\emph{Shape of $E(x)$.} $\Lambda'(S)=\pi\Aut_{\varphi(H)}(S_{\mathrm{can}})$, so $E(x)=\sigma R$.
\emph{Equivariance.} $\Lambda'(S^\pi)=\pi^{-1}\Lambda'(S)$ gives $E(x^h)=h^{-1}E(x)$.
\emph{(C2).} For $x^h$, the representative $h^{-1}\sigma$ gives $F_R(x^{hh^{-1}\sigma})=F_R(x^\sigma)$.
\emph{(C1), $\subseteq$.} Every $e\in\Lambda_{E(x)}(x)$ satisfies $x^e=F_H(x)$.
\emph{(C1), $\supseteq$.} Suppose $x^h=F:=F_H(x)$. We have $F=x^e$ for some $e\in E(x)$, so
$S(F)=S(x)^{\varphi(e)}=S_{\mathrm{can}}$. Then $S(x)^{\varphi(h)}=S(F)=S_{\mathrm{can}}$, so $\varphi(h)\in\Lambda'(S(x))$
and $h\in E(x)$. Hence $h\in\Lambda_{E(x)}(x)$.
\end{proof}

Isomorphic inputs $x$ and $x^h$ lead to the children $\sigma R$ and $h^{-1}\sigma R$. The coset rule maps both
to the same group node $(R,x^\sigma)$: the same group as a set, on the same ordered domain, with inputs
in one $R$-orbit. Every (P1) choice made at that node is therefore the same for both inputs.

\begin{named}[Lemma T (tuple colours)]
Let $H$ and $\varphi\colon H\to\Sym(\Gamma)$ be as above, let $\vT_0$ be an injective $t$-tuple of $\Gamma$
chosen by (P1), let $\cO:=\vT_0^{\varphi(H)}$, and let $J_0:=\{h\in H:\vT_0^{\varphi(h)}=\vT_0\}$. Let $\mathcal P$
be an $H$-invariant set of pairs $(y,\vT)$ with $\vT\in\cO$, let $\cY_0:=\{y:(y,\vT_0)\in\mathcal P\}$ (a
$J_0$-invariant set), and let $F_{J_0}$ be a $J_0$-canonizer on $\cY_0$. For $(y,\vT)\in\mathcal P$ put
\[ c(y,\vT):=F_{J_0}(y^\sigma)\quad\text{for any }\sigma\in H\text{ with }\vT^{\varphi(\sigma)}=\vT_0. \]
Then:
\begin{enumlab}{\alph}
\item $c$ is independent of $\sigma$;
\item $c(y^h,\vT^{\varphi(h)})=c(y,\vT)$ for all $h\in H$;
\item $c(y,\vT)=c(y',\vT')$ if and only if there is $h\in H$ with $y^h=y'$ and $\vT^{\varphi(h)}=\vT'$.
\end{enumlab}
\end{named}
\begin{proof}
Note $(y,\vT)^\sigma=(y^\sigma,\vT_0)\in\mathcal P$, so $y^\sigma\in\cY_0$. (a) Two valid $\sigma,\sigma'$ differ by
$\sigma^{-1}\sigma'\in J_0$, and $F_{J_0}$ is constant on $J_0$-orbits. (b) $h^{-1}\sigma$ is valid for
$(y^h,\vT^{\varphi(h)})$. (c) ``If'' is (b). For ``only if'', equal canonical forms give $j\in J_0$ with
$y^{\sigma j}=y'^{\sigma'}$. Put $h:=\sigma j\sigma'^{-1}$. Then $y^h=y'$ and
$\vT^{\varphi(h)}=\vT_0^{\varphi(j\sigma'^{-1})}=\vT'$.
\end{proof}
Lemma~T is Lemma~S with $S(y,\vT):=\vT$ and $\Lambda'(\vT):=\{\pi\in\varphi(H):\vT^\pi=\vT_0\}$, so that $R=J_0$.

\begin{remark}[Two naive variants fail]\label{rem:naive}
Take $H=S_4$ on $[4]$, $\varphi=\mathrm{id}$, $t=1$, and let $F_K(z)$ be the lexicographically least string in
the $K$-orbit of $z$, with $a<b$.
\begin{itemize}
\item \emph{No transport}, $c(x,\vT):=F_{H_{(T)}}(x)$. For $x=(a,b,b,b)$ and $T=(1)$ this gives $(a,b,b,b)$.
  For $h=(1\,2)$, $x^h=(b,a,b,b)$ and $T^h=(2)$, it gives $(b,a,b,b)$, since position~2 is pinned to $a$.
  Two isomorphic marked pairs get different colours, so (C2) fails. With transport, $\sigma=(1\,2)$ maps
  $2\mapsto1$ and $(x^h)^\sigma=x$, which gives the same colour.
\item \emph{Reference tuple read off a coset representative} (the first $t$ labels pulled back by the
  representative of the node coset). Isomorphic inputs then use conjugate groups $J_0^{\varphi(h)}$, whose
  canonical forms are different strings; the ranking of colours is permuted, and strong canonicity
  fails.
\end{itemize}
Our $\vT_0$ depends only on the group of the node (P1).
\end{remark}

% =====================================================================
\section{Canonical aggregation}\label{sec:agg}

Throughout this section the node group is $H$ (in the algorithm, $H=G'$), it carries an almost $d$-ary
sequence, $\varphi\colon H\to S_k$ satisfies $\varphi(H)\in\{A_k,S_k\}$, $t:=\max\{9,3+\lceil\log d\rceil\}$ (GNS
write $3+\log d$; we round up so that $t$ is an integer) and $k>10t$ (so $k\ge91$). The ideal domain $[k]$
is ordered (P1).

\subsection{Local certificates are canonical objects}

\begin{lemma}\label{lem:51}
For every $t$-set $T\subseteq[k]$, the chain $(W_i(T,x),G_i(T,x))_i$, the fullness of $T$, and $K(T,x)$ for
full $T$ can be computed within the recursive budget of \ref{L5}. All of these are set-theoretic
functions of $(H,\varphi,T,x)$ and are equivariant by \ref{L4}. In particular, the way the groups are
computed does not affect canonicity.
\end{lemma}
\begin{proof}
Put $W^*:=W_{i+1}\setminus W_i$; it is $G_i$-invariant by \ref{L2}, and $G_{i+1}=\Aut_{G_i}(\win{x}{W^*})$.
Compute a $G_i$-canonizer $(F,\Lambda)$ for $\win{x}{W^*}$ by Lemma~F1 and output
$G_{i+1}=\lambda\Aut_{G_i}(F)\lambda^{-1}$ for $\lambda\in\Lambda$. If $|W^*|\le n/2$, this is one recursive call on
$W^*$ under $G_i^{W^*}$; the almost $d$-ary sequence is inherited. Otherwise, $W^*$ is a union of affected
$G_i$-orbits (with respect to $\psi_T|_{G_i}$, a giant by the loop condition). Descend (Lemma~F3) from
$G_i$ to $N:=\ker(\psi_T|_{G_i})$ over at most $t!$ cosets; in each branch apply the chain rule (on $W^*$, via Lemma~F1). Its windows
lie inside $N$-orbits, which have size $\le|W^*|/t$ by \cite[Thm.~5.9]{GNS} ($t\ge9$). These are GNS's
counts.
\end{proof}

\subsection{The fullness group and the case analysis}

Put $\Phi(x):=\langle K(T,x):T\text{ full}\rangle\le\Aut_H(x)$ (the \emph{fullness group}) and
$\supp_x:=\supp(\varphi(\Phi(x)))\subseteq[k]$ (its \emph{support}). Both are equivariant: $\Phi(x^h)=\Phi(x)^h$ and
$\supp_{x^h}=(\supp_x)^{\varphi(h)}$ by \ref{L4}. We follow
\cite[App.~A]{GNS}, canonically.

\begin{center}\small
\begin{tabularx}{\linewidth}{@{}L{3.6cm}L{2.6cm}Yl@{}}
\toprule
case & GNS (App.~A) & canonical object & count\\
\midrule
(i) $k/4\le|\supp_x|\le3k/4$ & colouring by $\supp_x$ & $\fA=([k];\supp_x)$ & $r=1$\\
(ii) $|\supp_x|>3k/4$, no $\varphi(\Phi)$-orbit $\ge3k/4$ & orbit partition & $\fA=([k];\text{orbit partition})$ & $r=1$\\
(iii) $|\supp_x|>3k/4$, orbit $C$ with $|C|\ge3k/4$ and $\varphi(\Phi)|_C\ge\Alt(C)$ & ``Option 2'' & $\Delta(x):=C$, $K(x):=\Phi(x)$ & ---\\
(iv) as (iii), $\varphi(\Phi)|_C$ not a giant & individualize $I$, an orbital $R$, possibly a vertex $v$ & the family of \emph{all} structures over all choices of $(I,R,v)$, each on $[k]$ with $C'=C\setminus I$ marked & $r\le k^6$\\
(v) $|\supp_x|<k/4$ & joint comparison classes (pairwise) & $\fA(x)$ of Definition~\ref{def:Ax} & $r=1$\\
\bottomrule
\end{tabularx}
\end{center}
In case (iv), $|I|$ is the degree of transitivity of $\varphi(\Phi)|_C$ minus one, which is at most $4$ by
\cite[Thm.~A.2]{GNS} (a consequence of the classification of finite simple groups). In every structure,
the relevant domain $D\subseteq[k]$ is marked by a unary relation.

\begin{proposition}\label{prop:cases}
\begin{enumerate}
\item The case is an invariant of the $H$-orbit of $x$.
\item In cases (i), (ii), (iv), (v) the family $\fF(x)$ satisfies $\fF(x^h)=\fF(x)^{\varphi(h)}$, and every
  member, viewed as a structure on $[k]$, has relative symmetry defect at least $1/4$.
\item In case (iii), $\Delta(x^h)=\Delta(x)^{\varphi(h)}$, $K(x^h)=K(x)^h$, $K(x)\le\Aut_{H_{\Delta(x)}}(x)$ and
  $\varphi(K(x))|_{\Delta}\ge\Alt(\Delta)$.
\end{enumerate}
\end{proposition}
\begin{proof}
\emph{Equivariance.} All ingredients are equivariant: $\Phi$, $\supp_x$, the orbits, the degree of transitivity,
the orbitals, and in case (v) Lemma~\ref{lem:Tprime}(1). The family in case (iv) is the set of all
choices, so it is equivariant as a set. GNS's ``without loss of generality $R_1$ is the diagonal''
\src{GNS}{2357} is harmless since the diagonal is an invariant. GNS's ``it can be assumed that
$|D_1|=|D_2|$'' \src{GNS}{2368} is not needed, since $|D(x)|$ is an invariant.

\emph{Defects.} Let $\Alt(\Delta)\le\Aut(\fA)$ with $|\Delta|\ge3$. Since $D$ is a unary relation and $\Alt(\Delta)$
is transitive on $\Delta$, either $\Delta\subseteq D$ or $\Delta\cap D=\emptyset$.
(i) $|\Delta|\le\max(|\supp_x|,k-|\supp_x|)\le3k/4$.
(ii) $\Alt(\Delta)$ preserves the orbit partition and is primitive on $\Delta$. If $\Delta$ lies in one class,
$|\Delta|<3k/4$. Otherwise $\Delta$ meets each class in at most one point, and a class $P$ meeting $\Delta$ with
$|P|\ge2$ would contain a point fixed by $\Alt(\Delta)$, forcing $P^a=P$ for all $a\in\Alt(\Delta)$, which is
impossible; so every class meeting $\Delta$ is a singleton, i.e.\ a point of $[k]\setminus\supp_x$, and
$|\Delta|\le|[k]\setminus\supp_x|<k/4$.
(iv) If $\Delta\cap C'=\emptyset$ then $|\Delta|\le k-|C'|\le k/4+4$. If $\Delta\subseteq C'$, then either $\Alt(\Delta)$
preserves a non-trivial regular graph on $C'$, whose relative symmetry defect is at least $1/2$
\cite[Lemma~A.3]{GNS} (= \cite[Cor.~2.4.13]{Bab15}), or it preserves the colouring by in- and
out-neighbours of an individualized vertex, whose classes have size $\le|C'|/2$. So $|\Delta|\le k/2$.
(v) Lemma~\ref{lem:Tprime}(3).
This also covers a small slip in GNS: \src{GNS}{2324} asserts $|D_{i,j}|\ge3k/4$, whereas
$C'=C\setminus I$ only has at least $3k/4-4$ elements; the defect bound is unaffected.

\emph{Case (iii).} This is GNS's Option~2: $\Phi(x)\le\Aut_H(x)$, and $C$ is a $\varphi(\Phi)$-orbit, so $\Phi(x)\le H_C$.
\end{proof}

\subsection{The canonical structure in case (v)}\label{sec:colours}

Let $D(x):=[k]\setminus\supp_x$, so $|D(x)|>3k/4$. Every $t$-set $T\subseteq D(x)$ is non-full: a full $T$ has a
certificate in $\Phi(x)$ acting on $T$ as at least $\Alt(T)$, with $t\ge3$, so $T\subseteq\supp_x$. Let
$\vT_0:=(1,\dots,t)$ be the first $t$ labels of $[k]$ (P1). Since $\varphi(H)\ge A_k$ is $(k-2)$-transitive, it
is transitive on injective $t$-tuples. Put $J_0:=H_{(\vT_0)}=\ker(\psi_{T_0}\colon H_{T_0}\to\Sym(T_0))$.

\begin{definition}\label{def:Ax}
Let $\mathcal P$ be the set of all pairs $(\lay{z}{W(T',z)},\vT')$, where $z$ is a string and $\vT'$ an
injective $t$-tuple with underlying set $T'$. For $\vT\in D(x)^{\langle t\rangle}$ with underlying set $T$, put
$c_x(\vT):=c(\lay{x}{W(T,x)},\vT)$, as in Lemma~T for this $\mathcal P$ and the $J_0$-canonizer of
Lemma~\ref{lem:Tpp}. The structure is
\[ \fA(x):=\big([k];\ \text{unary }D(x);\ \text{$t$-ary colouring }\kappa_x\big),\qquad
   \kappa_x(\vT):=\text{rank of }c_x(\vT)\text{ among all values of }c_x, \]
encoded as a string on $[k]\sqcup[k]^t$, with value $0$ on tuples outside $D(x)^{\langle t\rangle}$.
\end{definition}
The set $\mathcal P$ is $H$-invariant by \ref{L4}. The set of values of $c_x$ is an invariant of the
$H$-orbit of $x$ by Lemma~T(b) and \ref{L4}, so the ranking is canonical.

\begin{lemma}\label{lem:Tprime}
\begin{enumerate}
\item (Strong canonicity.) $\fA(x^h)=\fA(x)^{\varphi(h)}$ for $h\in H$.
\item (Key identity.) For non-full $T$: $G(T,x)=\Aut_{H_T}(\lay{x}{W(T,x)})$, and $\psi_T(G(T,x))$ is not a
  giant.
\item (Symmetry defect.) If $\Alt(\Delta)\le\Aut(\fA(x))$ with $|\Delta|\ge3$, then $\Delta\subseteq[k]\setminus D(x)$ or
  $|\Delta|\le t+1$. Hence the symmetry defect of $\fA(x)$ is at least $k-\max(k/4,t+1)\ge3k/4$.
\end{enumerate}
\end{lemma}
\begin{proof}
(1) $W(T^{\varphi(h)},x^h)=W(T,x)^h$ by \ref{L4}, $D$ is equivariant, and Lemma~T(b) applies.
(2) This is \ref{L2} together with \ref{L3}.
(3) Either $\Delta\subseteq D$ or $\Delta\cap D=\emptyset$. Suppose $\Delta\subseteq D$ and $|\Delta|\ge t+2$, and pick $T\subseteq\Delta$
with $|T|=t$ and an ordering $\vT$. Every $\pi\in\Sym(T)$ is induced by an element of $\Alt(\Delta)$ (correct the
parity on $\Delta\setminus T$, which has at least two points). So all orderings of $T$ have the same colour. The
window $W(T,x)$ depends only on the set $T$, so Lemma~T(c), applied to $y=\lay{x}{W(T,x)}$, gives for every
ordering $\vT'$ of $T$ an $h\in H$ with $y^h=y$ and $\vT^{\varphi(h)}=\vT'$. Such $h$ lie in $H_T$, and $y^h=y$ means
$h\in\Aut_{H_T}(y)=G(T,x)$ by (2). So $\psi_T(G(T,x))=\Sym(T)$, contradicting (2).
\end{proof}

\begin{remark}
Running GNS's comparison \cite[Alg.~3]{GNS} for $(T_1,x_1)$ and $(T_2,x_2)$, the final coset is
$\{h:T_1^{\varphi(h)}=T_2,\ (\win{x_1}{W(T_1,x_1)})^h=\win{x_2}{W(T_2,x_2)}\}$. So the colour classes of
Definition~\ref{def:Ax}, taken over $D(x_1)^{\langle t\rangle}\sqcup D(x_2)^{\langle t\rangle}$, coincide with GNS's
comparison classes, under the reading that an empty comparison coset means ``no morphism''
(GNS's Alg.~3 does not specify what happens when the coset becomes empty). The algorithm never uses
this fact; only the symmetry defect is used. Our bound $t+1$ differs from GNS's $t-1$
\src{GNS}{2383}; both suffice.
\end{remark}

\subsection{Computing the colours}

\begin{lemma}\label{lem:Tpp}
For $\vT\in D(x)^{\langle t\rangle}$, $c_x(\vT)$ can be computed in polynomial time plus recursive canonizer
calls on pairwise disjoint subsets of $W(T,x)$, each of size at most $n/2$, with groups that are
restrictions of subgroups of $H$. Over all at most $k^t$ tuples, the recursive sizes add up to at most
$k^t\cdot n$.
\end{lemma}
\begin{proof}
\emph{Transport.} Compute $\sigma\in H$ with $\vT^{\varphi(\sigma)}=\vT_0$ (a preimage) and put $z:=x^\sigma$. By
\ref{L4}, $(\lay{x}{W(T,x)})^\sigma=\lay{z}{W(T_0,z)}$ and $G_i(T_0,z)=G_i(T,x)^\sigma$.

\emph{Definition of $F_{J_0}$.} Apply Lemma~F2 with group $J_0$ on the $J_0$-invariant set
$\cY_0:=\{\lay{z}{W(T_0,z)}\}$, with windows
$V_{i+1}(y):=\Aff(\Aut_{H_{T_0}}(y),\psi_{T_0})\setminus U_i$ for as long as
$\psi_{T_0}(\Aut_{H_{T_0}}(y))\ge\Alt(T_0)$ and this set is non-empty. These windows are $H_{T_0}$-equivariant,
hence $J_0$-equivariant, and by \ref{L2} they reproduce the chain of $(T_0,z)$. So $F_{J_0}$ is a fixed
deterministic function on $\cY_0$, which is what Lemma~T needs.

\emph{Stage $i$.} Let $F_i$ be the current canonical form and $\tau_i\in\Lambda_i$. Then
$B:=\Aut_{H_{T_0}}(F_i)=G_i(T,x)^{\sigma\tau_i}$, which is known from Lemma~\ref{lem:51}, and
$A_i=\Aut_{J_0}(F_i)=\ker(\psi_{T_0}|_B)$. The new window $V'$ is a union of $B$-orbits that are affected
with respect to the giant representation $\psi_{T_0}|_B$ of degree $t\ge9$. By \cite[Thm.~5.9]{GNS}, every
$A_i$-orbit in $V'$ has length at most $n/t$. Canonize $\win{w}{V'}$ under $A_i$ by the chain rule (after
restricting to $V'$ by Lemma~F1); its windows lie in $A_i$-orbits. (If $|V'|\le n/2$ one may instead make
one call on $V'$; whichever rule is used must be a fixed function of $|V'|$, which is $J_0$-invariant, so
that $F_{J_0}$ remains a single function on $\cY_0$.)

\emph{Sizes.} The sizes for one tuple sum to at most $|W(T,x)|\le n$. No descent over $t!$ cosets is
needed, because the right subgroup $A_i$ is already the kernel. GNS spend
$\binom kt^2t!\,n$ on comparisons; we spend at most $k^t n$.
\end{proof}

\subsection{Option 1: canonical find-structure}

\begin{proposition}\label{prop:opt1}
In cases (i), (ii), (iv), (v), an $H$-canonizer for $x$ is obtained as follows. For each $\fA\in\fF(x)$, let
$s_\fA$ be the string encoding $\fA$ on $[k]\sqcup[k]^2\sqcup[k]^t$. Let $F^{\mathrm{BC}}$ be Babai's canonical form
\cite[Thm.~2.2]{Bab19}, called only for the group $S_k$ in its induced action on this domain, on the
generating set determined by $(k,t)$ (requirement~\ref{D4}), and put
$\Lambda^{S_k}(\fA):=\Iso_{S_k}(s_\fA,F^{\mathrm{BC}}(s_\fA))$, computed by Babai's string isomorphism test
\cite[Thm.~2.1]{Bab19}. If $\varphi(H)=S_k$, put $\Lambda'(\fA):=\Lambda^{S_k}(\fA)$. If $\varphi(H)=A_k$, then
\begin{itemize}
\item if $\Aut_{S_k}(s_\fA)$ contains an odd permutation, put $\Lambda'(\fA):=\Lambda^{S_k}(\fA)\cap A_k$, with form
  $F^{\mathrm{BC}}(s_\fA)$;
\item otherwise the elements of $\Lambda^{S_k}(\fA)$ all have one parity; if it is even, put
  $\Lambda'(\fA):=\Lambda^{S_k}(\fA)$, with form $F^{\mathrm{BC}}(s_\fA)$, and if it is odd, put
  $\Lambda'(\fA):=\Lambda^{S_k}(\fA)\,(1\,2)$, with form $F^{\mathrm{BC}}(s_\fA)^{(1\,2)}$.
\end{itemize}
Then $\varphi^{-1}(\Lambda'(\fA))=\sigma_\fA R_\fA$ with
$R_\fA=\varphi^{-1}(\Aut_{\varphi(H)}(\fA_{\mathrm{can}}))$. Recursively canonize $x^{\sigma_\fA}$ under $R_\fA$, and
combine over $\fA\in\fF(x)$ by Lemma~S (for pairs $(x,\fA)$) and Lemma~F4.
\end{proposition}
\begin{proof}
$\Lambda'$ is a $\varphi(H)$-canonizer on structures: for $\varphi(H)=S_k$ because $F^{\mathrm{BC}}$, on a generating set
determined by $(k,t)$, is constant on $S_k$-orbits and takes values in them; for $\varphi(H)=A_k$ because
$\Lambda^{S_k}(s^a)=a^{-1}\Lambda^{S_k}(s)$ for $a\in A_k$, so the case distinction and the parity are constant on
$A_k$-orbits, and in each case the displayed coset is exactly the set of $g\in A_k$ mapping $s_\fA$ to the stated
form (in the odd case, $s^g=F^{(1\,2)}$ iff $g(1\,2)\in\Lambda^{S_k}(s)$). Lemma~S with $S(x,\fA)=\fA$ is equivariant
on pairs; then apply Lemma~F4.
\end{proof}
Each $R_\fA$ satisfies $|\varphi(H):\varphi(R_\fA)|\ge(4/3)^k$, by Proposition~\ref{prop:cases}(2) and
\cite[Cor.~6.5]{GNS} ($k\ge24$). The number of branches is $r\le k^6$, as in \cite[Lemma~6.7]{GNS}, where
$\fA_{1,1}$ is paired with each $\fA_{2,j}$. The calls of BC are on domains of size at most $2k^t=d^{O(\log d)}$,
so each costs $\exp(O((\log d)^{2c_{\mathrm{BC}}}))$, where $\exp(O((\log n)^{c_{\mathrm{BC}}}))$ is the running time
of BC.

\begin{remark}[What Theorem~A needs from BC]\label{rem:BC}
BC is called only in Proposition~\ref{prop:opt1}, only at Option~1 nodes, where $\varphi(H)\in\{A_k,S_k\}$ and
$k\le d$ (\S\ref{sec:gnsrecap}), and the case $A_k$ is obtained from $S_k$ by the parity argument above. So
Theorem~A needs exactly: a canonical form, under $S_k$ with $k\le d$, of relational structures on $[k]$ of
arity at most $t=O(\log d)$ (arity at most $2$ in cases (i), (ii), (iv), and $t$ in case (v)), together with Babai's
string isomorphism test \cite[Thm.~2.1]{Bab19} (the main result of \cite{Bab15}) for $\Lambda^{S_k}$. Through
the standard encoding of a $t$-ary structure by a graph on $k^{O(t)}$ vertices (cf.\ \cite[proof of
Lemma~4.33]{Neu20}), this is a special case of \cite[Cor.~2.3]{Bab19} on graphs with $k^{O(t)}=d^{O(\log d)}$
vertices. It is \emph{not} a reduction to graphs on at most $d$ vertices: for $t\ge3$ there are more $t$-ary
structures on $[k]$ than graphs on $[k]$, so case (v) cannot be encoded on $k$ vertices. Since
\cite[Cor.~2.3]{Bab19} is deduced from \cite[Thm.~2.2]{Bab19}, the dependency is also no lighter in
status: the published account of both is the preliminary report \cite{Bab19}, which says of its
main theorem, ``we sketch the proof''. This is why the dependency is flagged in the abstract and in
the status of Theorem~A.
\end{remark}

\subsection{Option 2: canonical find-symmetry}

\begin{lemma}[canonical find-symmetry]\label{lem:K}
In case (iii), an $H$-canonizer is obtained with two recursive canonizer calls under
$L:=H_{(\Delta_0)}:=\varphi^{-1}(\varphi(H)_{(\Delta_0)})$, where $\Delta_0:=\{1,\dots,|\Delta|\}$ (P1). Moreover
$|\varphi(H):\varphi(L)|\ge|\Alt(\Delta_0)|\ge(4/3)^k$.
\end{lemma}
\begin{proof}
\emph{Move $\Delta$ to $\Delta_0$.} Apply Lemma~S with $S(x)=\Delta(x)$ and
$\Lambda'(\Delta)=\{\pi\in\varphi(H):\Delta^\pi=\Delta_0\}$; this is well defined because $A_k$ is transitive on
$|\Delta|$-sets. We get $E(x)=\sigma H_{\Delta_0}$. Put $z:=x^\sigma$, so $\Delta(z)=\Delta_0$ and
$K(z)\le\Aut(z)\cap H_{\Delta_0}$.

\emph{Collapse to two cosets.} $L\trianglelefteq H_{\Delta_0}$, and the image of $K(z)$ on $\Delta_0$ contains
$\Alt(\Delta_0)$, so $K(z)L$ has index at most $2$ in $H_{\Delta_0}$ and is normal. Hence
$H_{\Delta_0}=K(z)L\cup K(z)L\tau$, where, if $K(z)L\ne H_{\Delta_0}$, $\tau$ is any element of
$H_{\Delta_0}\setminus K(z)L$ (one inducing an odd permutation on $\Delta_0$); if $K(z)L=H_{\Delta_0}$, omit the
second part (this also covers the case $\varphi(H)=A_k$, which GNS do not discuss).

\emph{Branch values.} In the descent (Lemma~F3) from $H_{\Delta_0}$ to $L$, a coset $\rho L$ with
$\rho=\kappa l\in K(z)L$ has branch value $F_L(z^{\kappa l})=F_L(z)$. A coset in $K(z)L\tau$ has a representative
$\kappa\tau$, and branch value $F_L(z^{\kappa\tau})=F_L(z^\tau)$.

\emph{Output.} Two $L$-calls suffice. With $\lambda_1\in\Lambda_L(z)$ and $\lambda_\tau\in\Lambda_L(z^\tau)$:
$\Aut(z)\cap K(z)L=K(z)\Aut_L(z)$; $\Aut(z)$ meets $K(z)L\tau$ iff $F_L(z^\tau)=F_L(z)$, and then it contains
$g_\tau:=\tau\lambda_\tau\lambda_1^{-1}$; so $\Aut_{H_{\Delta_0}}(z)=\langle K(z),\Aut_L(z)[,g_\tau]\rangle$. Finally
$\Lambda=\Aut(z)\lambda_1$ if $F_L(z)\le F_L(z^\tau)$, and $\Lambda=\Aut(z)\tau\lambda_\tau$ otherwise. By Lemma~F3 this is
exactly $\{g:z^g=m\}$.
\end{proof}
This matches \cite[Lemma~6.9]{GNS}: $r\in\{1,2\}$ and $H=G_{(\Delta_1)}$.

% =====================================================================
\section{The algorithm and its correctness}\label{sec:algo}

\subsection{Instances}

An \emph{instance} is a tuple $(\Omega,H,\vec\fB,x)$, where $\Omega$ is ordered, $H\le\Sym(\Omega)$ is given by its
canonical generating set, $\vec\fB$ is an $H$-invariant almost $d$-ary sequence, and $d$ and $t$ are fixed
globally. The sequence $\vec\fB$ is auxiliary data. It is always computed from groups higher in the
recursion and never from $x$, so isomorphic strings lead to instances that differ only in the string.

\subsection{Algorithm CanAD}\label{sec:canad}

\begin{description}
\item[Step 1.] If $|\Omega|=1$, return $(x,H)$.
\item[Step 2.] \textbf{Intransitive $H$.} Apply the chain rule (Definition~\ref{def:chain}). Each window $V'$ is
  handled by CanAD on the transitive instance $(V',A_i^{V'},\vec\fB[V'],w|_{V'})$.
\item[Step 3.] \textbf{Transitive $H$, $H^{\fB_1}$ semi-regular.} Then $H^{\fB_1}$ is regular, with $b:=|\fB_1|\ge2$
  elements. Descend (Lemma~F3) to $H_{(\fB_1)}$ over $b$ cosets. In each branch apply the chain rule. Its
  windows are orbits of subgroups of $H_{(\fB_1)}$, so each lies in a block of $\fB_1$ and has size at most
  $n/b$; per branch the window sizes sum to $n$ (there may be more than $b$ windows).
\item[Step 4.] \textbf{Transitive $H$, $|\fB_1|\le d$.} Choose (P1) a block system $\fB\succeq\fB_1$ of $H$ on which
  $P:=H^{\fB}$ is primitive, and put $b:=|\fB|\le d$.
  \begin{enumlab}{\alph}
  \item If $|P|\le b^{1+\log b}$, or $b$ is below the constant of \cite[Lemma~7.1]{GNS}: descend to
    $H_{(\fB)}$ over $|P|$ cosets and apply the chain rule; its windows have size $\le n/b\le n/2$.
  \item Otherwise apply \cite[Lemma~7.1]{GNS} to $P$ (P1), obtaining $N$, $\cC$, $\psi$ and $k$. Put
    $G':=\{g\in H:g^{\fB}\in N\}$, $\varphi:=\psi\circ(g\mapsto g^{\fB})$, and let $\Cp$ be the lift of $\cC$ to
    $\Omega$, so $\ker\varphi=G'_{(\Cp)}$. Descend (Lemma~F3) from $H$ to $G'$ over at most $b\le d$ cosets, and
    in each branch $\rho$ call $\mathrm{Agg}(G',\varphi,\Cp,\vec\fB,x^\rho)$.
  \end{enumlab}
\item[Step 5.] \textbf{$\mathrm{Agg}(G',\varphi,\Cp,\vec\fB,y)$.}
  \begin{itemize}
  \item If $k\le10t$: descend to $\ker\varphi=G'_{(\Cp)}$ over at most $k!$ cosets, then apply the chain rule.
    The windows lie inside $\Cp$-classes, of size $\le n/2$.
  \item Otherwise compute the local certificates (Lemma~\ref{lem:51}), $\Phi$, $\supp_x$, and the case
    (Proposition~\ref{prop:cases}). In cases (i), (ii), (iv), (v) apply Proposition~\ref{prop:opt1}
    (for (v) the colours come from Lemma~\ref{lem:Tpp}). In case (iii) apply Lemma~\ref{lem:K}.
  \item Every recursive call \emph{on the same domain}, under some $R_\fA$ or $L$, is made with the
    sequence $\vec\fB\oplus\Cp$ of Step~6.
  \end{itemize}
\item[Step 6.] \textbf{The top partition of same-domain children.} Put
  $\vec\fB\oplus\Cp:=(\{\Omega\}\succ\Cp\succ\fB_1\succ\fB_2\succ\cdots)$, with duplicates removed. Here
  $\Cp\succeq\fB\succeq\fB_1$. The sequence is invariant under $R_\fA,L\le G'$, because $\Cp$ is
  $G'$-invariant. It is almost $d$-ary: $|\Cp|=|\cC|\le b\le d$, and each $\Cp$-class contains at most
  $|\fB_1|\le d$ blocks of $\fB_1$; the lower levels are inherited.
\end{description}
Step~6 is the choice discussed in Remark~\ref{rem:gap1}. Since $\Cp\succeq\fB\succeq\fB_1$, whenever Step~4
applies at a same-domain child, the algorithm picks a block system above $\Cp$ restricted to the relevant
orbit, and such a block system is also one above $\fB_1$, so it is a choice that the GNS algorithm also
permits there. (CanAD tests semi-regularity on $\Cp[\Omega']$ rather than on $\fB_1[\Omega']$, so it is a
variant of the GNS algorithm rather than literally one of its runs; Section~\ref{sec:time} analyses CanAD
directly.)

\subsection{Correctness}

\begin{proposition}\label{prop:correct}
CanAD returns an $H$-canonizer for $x$.
\end{proposition}
\begin{proof}
\emph{Termination.} Induction on $(|\Omega|,|H|)$, ordered lexicographically. Every same-domain recursive
call is on a proper subgroup, since $R_\fA$ and $L$ have index at least $(4/3)^k>1$ in $G'\le H$. All other
calls are on smaller domains.
\emph{Each step preserves canonizers.} Each step is one of Lemmas F1--F4, S, \ref{lem:K}, Proposition~\ref{prop:opt1},
or a (P1) choice, and each of these maps canonizers of the sub-instances to a canonizer of the
instance.
\emph{Canonicity of auxiliary data.} Every auxiliary object is either a function of the node's group (P1)
or an $H$-equivariant function of $x$ (Lemma~\ref{lem:51}, Proposition~\ref{prop:cases},
Lemma~\ref{lem:Tprime}). The colours use one fixed $J_0$-canonizer (Lemma~\ref{lem:Tpp}).
\end{proof}

\subsection{The choice points of GNS}

Table~\ref{tab:choices} lists every place where GNS make an arbitrary choice or compare two inputs,
and how CanAD treats it. Line numbers refer to the arXiv~v3 source.

\begin{table}[!htbp]
\centering\small
\begin{tabularx}{\linewidth}{@{}rL{3.1cm}L{4.1cm}Y@{}}
\toprule
\# & GNS place (src.\ line) & choice & canonical treatment\\
\midrule
1 & Thm.~4.3 proof; 1103 & maximal block of $G_B^B$, ``$B$ arbitrary'' & P1 (block of the least point)\\
2 & Thm.~4.3; 1120--1121 & numbering $B_1,\dots,B_s$; transporters $\sigma_{1\to i}$ & P1 (blocks by least point; lex-least transporters)\\
3 & Thm.~4.3; 1207, 1254 & one $G^*$-orbit $A^*$ & P1 (orbit of the least point of $\Omega^*$)\\
4 & Thm.~4.3; 1216 & $\mathcal P$-sequence of Thm.~3.13 for $N^{B_1}$ & P1 (deterministic computation)\\
5 & Thm.~4.6; 1315 & Johnson identifications $\rho_{i,B}$ (\cite{BLS87}; Lemma~4.4) & P1\\
6 & Lemmas 6.2/6.3; 1802, 1908 & windows $\Aff(G_i,\psi)$ & equivariant \ref{L4}; no choice\\
7 & Lemmas 6.2/6.3; 1810, 1917 & coset representatives of $N$ & descent over all cosets (F3)\\
8 & Lemma 6.3; 1902 & $\sigma_0$ with $T_1\mapsto T_2$ & disappears (Lemma~T(a))\\
9 & Lemma 7.2; 2087 & coarsening to primitive $P$ & P1\\
10 & Lemma 7.2; 2090--2093 & $N,\cC,\psi$ (Lemma~7.1); $G'$; labels of $[k]$ & P1, then descent over $\le d$ cosets\\
11 & Lemma 7.2; 2101 & $k\le10t$: Luks reduction to $\ker\varphi$ & descent over $\le k!$ cosets\\
12 & Lemma 7.2, Alg.~1; 2121--2122 & order of orbits & chain rule (Definition~\ref{def:chain})\\
13 & App.~A; 2337--2349 & $\Phi$, $\supp_x$, case split & equivariant (Proposition~\ref{prop:cases})\\
14 & App.~A; 2352, 2359, 2364 & $I$, orbital $R_i$, vertex & equivariant family and F4\\
15 & App.~A; 2357--2358 & ``WLOG $R_1$ is the diagonal'' & the diagonal is an invariant\\
16 & App.~A; 2368 & ``WLOG $|D_1|=|D_2|$'' & pairwise only; $|D(x)|$ is an invariant\\
17 & App.~A; 2370--2383 & joint classes on $D_1'\sqcup D_2'$ & \textbf{Lemma T: the only essential change}\\
18 & Lemma 7.2; 2118--2126 (implicit) & block system of the next round & $\Cp$ as top partition (Step~6; Remark~\ref{rem:gap1})\\
19 & Lemma 6.7; 2017--2019 & pair $\fA_{1,1}$ with each $\fA_{2,j}$; isomorphism test of~\cite{Bab16} on structures & Lemma~S with BC per member of the family, then F4 (Proposition~\ref{prop:opt1})\\
20 & Lemma 6.9; 2049--2051 & $g$ with $\Delta_1\mapsto\Delta_2$; a transposition $\tau$ on $\Delta_1$ & Lemma~S to $\Delta_0$ (P1), then Lemma~\ref{lem:K}\\
\bottomrule
\end{tabularx}
\caption{Choice points and pairwise steps of GNS.}\label{tab:choices}
\end{table}

No other pairwise step exists in GNS. The places where two strings occur are \src{GNS}{1866--1940},
\src{GNS}{1980--1990}, \src{GNS}{1999--2057} and \src{GNS}{2370--2385}. The phrase ``for $i=1,2$'' in
\cite[Lemma~6.6]{GNS} is presentational everywhere except row~17.

\subsection{Theorem A from CanAD}\label{sec:thmA}

\begin{proof}[Proof of Theorem A (correctness)]
Cosets are reduced to groups by the coset rule. The first step computes canonical generators, so the
output depends only on the coset as a set. For intransitive $G$, apply the chain rule over orbits; each
constituent is a recursive call of the $\Gd$-canonizer, and each restricted group lies in $\Gd$. For
transitive $G$, compute $\Omega^*$, $\iota$, $x\mapsto x^*$ and the almost $d$-ary sequence by
\cite[Thms.~4.3, 4.6]{GNS} (P1; rows 1--5 of Table~\ref{tab:choices}), run CanAD on
$(\Omega^*,G^*,\vec\fB^*,x^*)$, and pull back by Lemma~R. The almost $d$-ary recursion never re-applies the
change of action \cite[Obs.~4.2]{GNS}. Correctness follows from Proposition~\ref{prop:correct} and
Lemma~R. The time bound is proved in Section~\ref{sec:timebound}.
\end{proof}
\status{Proof modulo the results listed in Section~\ref{sec:assumptions}, in particular Babai's canonical form
\cite[Thm.~2.2]{Bab19}, whose published account is a preliminary report that sketches the proof (used only for
$S_k$, $k\le d$, on structures of arity $O(\log d)$; Remark~\ref{rem:BC}), the internal constructions of
\cite[Thms.~4.3, 4.6]{GNS}, and requirements \ref{D1}--\ref{D4} (deterministic subroutines on canonical
generating sets); see Questions~\ref{q:changeofaction}, \ref{q:BC}.}

% =====================================================================
\section{Running time}\label{sec:time}

\subsection{A recursion lemma}

\begin{named}[Lemma REC (a recursion lemma)]
Let $\tau\colon\mathbb N\to\mathbb N$ with $\tau(1)=1$ and let $K\ge1$. Suppose every $n\ge2$ satisfies
$\tau(n)\le\sum_{i=1}^\ell\tau(n_i)$ with $1\le n_i<n$, where either
\begin{enumlab}{\Alph}
\item $n_i\le n/2$ for all $i$ and $\sum_i(n_i/n)^{K+1}\le1$, or
\item $\sum_in_i\le n$.
\end{enumlab}
Then $\tau(n)\le n^{K+1}$.
\end{named}
\begin{proof}
By induction: $\sum n_i^{K+1}=n^{K+1}\sum(n_i/n)^{K+1}\le n^{K+1}$ in case (A), and
$\sum n_i^{K+1}\le(\sum n_i)^{K+1}\le n^{K+1}$ in case (B).
\end{proof}
We apply the lemma to $\tau(n):=$ the maximum, over all instances with domain of size $n$, of the
number of leaves of the recursion tree of the algorithm on that instance. The hypothesis is a statement
about one instance at a time, and each instance of size $n$ has children of sizes $n_i$ as in (A) or (B).

The hypothesis of (A) holds in two situations.
\begin{itemize}
\item $n_i\le n/2$ and $\sum n_i\le2^Kn$ \cite[Lemma~2.6]{GNS}.
\item The semi-regular step (Step~3): $b\ge2$ branches, each with pieces of size $\le n/b$ that sum to $n$.
  Per branch $\sum(n_i/n)^{K+1}\le b^{-K}\sum n_i/n=b^{-K}$, so over all branches the sum is at most
  $b^{1-K}\le1$.
\end{itemize}

\begin{remark}\label{rem:gap2}
\cite[Lemma~2.8]{GNS} requires $\sum n_i\le2^Kn$. The semi-regular step of \cite[Alg.~4]{GNS} has
recursive sizes summing to $b\cdot n$ with $b$ unbounded (e.g.\ $C_n$ acting regularly), so the appeal to
Lemma~2.8 in the proof of \cite[Thm.~7.3]{GNS} does not literally cover it. Lemma~REC
covers it. This is a point of presentation only.
\end{remark}

\subsection{A question about the written potential argument of GNS Lemma 7.2}

\begin{remark}\label{rem:gap1}
In the proof of \cite[Lemma~7.2]{GNS}, the height of the recursion is bounded by showing that the
degree of the giant representation drops from $k$ to $k'\le k-k/(3\log k)$ from one round to the next.
The proof reads \src{GNS}{2126}: ``Observe that $|(H')^{\mathfrak C'}|\ge\frac{(k')!}{2}$ because
$(H')^{\varphi'}\ge A_{k'}$ and $H'_{(\mathfrak C')}\le\ker(\varphi')$.'' Combined with
$|H^{\Cp}|\le k!/(4/3)^k$, this gives the inequality and a recursion of height $O(\log^2d)$.

As far as we can see, the algorithm as written does not ensure $H'_{(\Cp)}\le\ker\varphi'$. In the next round
the algorithm starts again from $\fB_1$, restricted to the current orbit, and replaces it by ``a block system
of smaller size'' on which the group acts primitively \src{GNS}{2087}. Neuen's version of the same proof
says ``Let $\mathfrak B\succ\mathfrak B_1$ be a minimal block system'' \src{Neu20}{2683}. Nothing ties this
block system to $\Cp$. The partition $\Cp$ is used only in the analysis. The remark ``Recall that
$\mathfrak C'$ is $G'$-invariant and thus, it is also $H$-invariant'' \src{GNS}{2120} is needed in any case,
so that $H^{\Cp}$ is defined.

\emph{An example.} Let $m$ be large, $d:=m^2$, $q\ge1$, $\Omega:=[m]^2\times[q]$, and
$G:=(S_m\wr S_2)\times C_q$, with the product action on $[m]^2$ and $C_q$ regular on $[q]$. Take the
sequence $\{\Omega\}\succ\fB_1:=\{\{(a,b)\}\times[q]\}\succ$ discrete, which is almost $d$-ary. Let $U\subset[m]$
with $|U|\approx0.3m$, and $x((a,b),c):=[a\in U]$.
\begin{itemize}
\item \emph{First round.} $P=G^{\fB_1}=S_m\wr S_2$ is primitive. Every output of \cite[Lemma~7.1]{GNS} has
  $N\le S_m\times S_m$, $k=m$, and $\cC$ the rows (up to swapping coordinates), so $G'=S_m\times S_m\times C_q$ and
  $\ker\varphi=1\times S_m\times C_q$.
\item \emph{Aggregation.} A test set $T\subseteq[k]$ (a set of rows) is full iff $T\subseteq U$ or
  $T\subseteq U^c$. So $\varphi(\Phi(x))$ has orbits $U$ and $U^c$, $\supp_x=[k]$, and no orbit has $\ge3k/4$ points:
  case (ii), with structure $\{U,U^c\}$ and $H=\varphi^{-1}(\Sym(U)\times\Sym(U^c))\cap G'$.
\item \emph{Next round.} The large $H$-orbit is $\Omega'=(U^c\times[m])\times[q]$, of size about $0.7n$. The
  group $H'$ induced on the cells of $\Omega'$ is $\Sym(U^c)\times S_m$ in product action, which has exactly two
  minimal block systems: rows and columns. If the columns are chosen, $P'\cong S_m$ in its natural action,
  so $k'=m=k$, while $H'_{(\Cp)}$ (which fixes all rows) acts as $S_m$ on the columns and is not contained in
  $\ker\varphi'$.
\end{itemize}
For all sufficiently large $m$ the numerical hypotheses of \cite[Lemma~7.2]{GNS} hold (for instance
$m\ge200$ with $t=\max\{9,3+\lceil\log d\rceil\}$, apart from the unspecified constant of
\cite[Lemma~7.1]{GNS}). At $m=200$ the bound $k-k/(3\log k)$ is about $191.3$, the choice of columns gives
$k'=200$, and the choice of rows gives $k'=140$. A small computation confirmed the block systems and
groups above for $q=2$, $|U|=2$ and $m=5,\dots,8$.

\emph{The ``fewest blocks'' reading.} If ``a block system of smaller size'' is read as ``a block system with
the fewest blocks'', this example picks the rows ($0.7m$ blocks rather than $m$) and is avoided. We do not
know whether that reading suffices in general. The text does not state it, and Neuen's wording does not
suggest it; in any case the proof would need an argument for it.

\emph{A choice that makes the step valid.} Choose the next block system above $\Cp[\Omega']$. This is
always possible: $\Cp$ is $G'$-invariant, hence invariant under every same-domain descendant, and it has at
least two classes on the large orbit. Since $\Cp\succeq\fB\succeq\fB_1$, it is a run permitted by the algorithm
as written. With this choice the inequality $k'\le k-k/(3\log k)$ holds (Section~\ref{sec:potential}), and
the rest of the analysis is unchanged. Neuen's proof \cite[Lemma~4.36]{Neu20} contains the same step
\src{Neu20}{2716}, in the notation $\Lambda'_{(\mathfrak C')}\le\ker(\varphi')$, and the same choice applies there.

\emph{What is and is not claimed.} The theorem of GNS is not in doubt: the proposed choice is a legal run
of their algorithm, and with it their own inequality holds. We have not shown that the running time
fails under other choices; in the example the chain ends one round later. The same passage appears in all
arXiv versions of GNS and of~\cite{Neu20}. We found no erratum, but we could not consult the SIAM journal
version of GNS.
\end{remark}

\begin{question}\label{q:gap1}
Is ``choose the next block system above $\Cp$'' the intended reading of the proof of
\cite[Lemma~7.2]{GNS} (and of \cite[Lemma~4.36]{Neu20})? Is the point addressed in the SIAM journal
version? Does any other part of the GNS analysis use the original $\fB_1$ in a way this choice disturbs?
(We checked that local certificates need only some almost $d$-ary sequence, and that
\cite[Thm.~5.9, Cor.~6.5]{GNS} and the case analysis do not use the sequence at all.)
\end{question}

\subsection{The potential argument with $\Cp$ as top partition}\label{sec:potential}

Consider a same-domain chain: an aggregation node $(G',\varphi,k)$; a child $(R,\cdot)$ with $R\in\{R_\fA,L\}$
and $|\varphi(G'):\varphi(R)|\ge(4/3)^k$; the chain rule on $R$; and the window $\Omega'$ of size $>n/2$, if any,
processed under the current right subgroup $R^\circ\le R$. By the choice of chain rule, the next node
$H':=(R^\circ)^{\Omega'}$ is transitive, with top partition $\Cp[\Omega']$. Note $|\Cp[\Omega']|\ge2$, since
$\Cp$-classes have size $\le n/2<|\Omega'|$.

At $H'$, Step~3 or Step~4 is applied with $\fB_1=\Cp[\Omega']$. Let $\fD\succeq\Cp[\Omega']$ be the block system
chosen there. If Step~4(b) applies, it produces $G''\le H'$ (index $\le d$) and $\varphi'=\psi'\circ(\text{action
on }\fD)$ of degree $k'$.

\begin{lemma}\label{lem:potential}
In this situation, $k'!/2\le k!/(4/3)^k$, and hence $k'\le k-k/(3\log k)$.
\end{lemma}
\begin{proof}
\begin{enumerate}
\item An element of $G''$ fixing every class of $\Cp[\Omega']$ fixes every class of the coarser partition
  $\fD$, so it lies in $\ker\varphi'$. Hence $A_{k'}\le\varphi'(G'')$ is a quotient of $(G'')^{\Cp[\Omega']}$.
\item The $\Cp$-classes meeting $\Omega'$ form an $R^\circ$-invariant subset of $\Cp$, and $C\mapsto C\cap\Omega'$ is a
  bijection onto $\Cp[\Omega']$. So $(H')^{\Cp[\Omega']}$ is a quotient of $(R^\circ)^{\Cp}$, and
  $|(G'')^{\Cp[\Omega']}|\le|(R^\circ)^{\Cp}|\le|R^{\Cp}|$.
\item $R^{\Cp}\cong R/R_{(\Cp)}=R/(R\cap\ker\varphi)\cong\varphi(R)$, using $\ker\varphi=G'_{(\Cp)}$, and
  $|\varphi(R)|\le|\varphi(G')|/(4/3)^k\le k!/(4/3)^k$.
\item So $k!/k'!\ge(4/3)^k/2$. Since $k!/k'!\le k^{k-k'}$, we get $k\log(4/3)\le1+(k-k')\log k$, hence
  $k-k'\ge(0.415k-1)/\log k\ge k/(3\log k)$ for $k\ge13$.
\end{enumerate}
\end{proof}
(The arithmetic in step 4 was also checked numerically, by computing the largest $k'$ with
$k'!/2\le k!/(4/3)^k$ for every $13\le k\le5000$.)

\paragraph{Height.} Every giant degree is at most $d$, and aggregation stops once $k\le10t$. By
Lemma~\ref{lem:potential}, $k_{j+1}\le k_j(1-1/(3\log d))$ along a root-to-leaf path. So a path contains at
most $h:=3\log d\cdot\ln(d/(10t))+1=O(\log^2d)$ aggregation nodes.

\subsection{One macro step and its branch counts}

A \emph{macro step} is the computation for one transitive instance on $n$ points to which Step~4
applies. It explores a tree $\cT$ of same-domain nodes as in Section~\ref{sec:potential}: descents to
$G'$, aggregation nodes, their children $R_\fA$ or $L$, the chain rule on these children, and the transitive
node $H'$ on the window of size $>n/2$. Calls on domains of size $\le n/2$ are the \emph{exits} of the macro
step. If Step~3 (semi-regular) applies at an inner node $H'$, that node is a leaf of $\cT$; its children
have size $\le|\Omega'|/b\le n/2$ and total size $\le b|\Omega'|\le d\,n$, since $b=|\Cp[\Omega']|\le|\cC|\le d$, and
they are exits. (A semi-regular step at the \emph{root} of an instance, with $b$ unbounded, is a separate
(A)-node of Lemma~REC.)

\begin{center}\small
\begin{tabularx}{\linewidth}{@{}L{4.1cm}L{3.2cm}YL{2.3cm}@{}}
\toprule
quantity & GNS & canonical & where\\
\midrule
descent $H\to G'$ & $\le d$ & $\le d$ & F3\\
Luks reduction, small primitive $P$ & $|P|\le d^{1+\log d}$ & same & F3\\
$k\le10t$, to $\ker\varphi$ & $\le k!$ & same & F3\\
semi-regular inner leaf & $b\le d$ cosets & same & F3, chain rule\\
local certificates & $\binom kt t!\,n$ & same & Lemma~\ref{lem:51}\\
comparisons / colours & $\binom kt^2t!\,n$ & $\le k^tn$ & Lemma~\ref{lem:Tpp}\\
Option~1 branches & $r\le k^6$ & $r\le k^6$ & Prop.~\ref{prop:opt1}\\
Option~2 branches & $r\le2$ & 2 calls under $L$ & Lemma~\ref{lem:K}\\
index of children & $\ge(4/3)^k$ & $\ge(4/3)^k$ & \cite[Cor.~6.5]{GNS}\\
black box on structures & $\le k^6$ isomorphism tests & $\le k^6$ BC and isomorphism calls & Prop.~\ref{prop:opt1}\\
extra non-recursive work & --- & canonical generators, lex-least transporters: poly$(n)$ per node & \S\ref{sec:P1}\\
\bottomrule
\end{tabularx}
\end{center}

\emph{Canonization explores no branch that the isomorphism test does not.} GNS never prune by ``one
isomorphism suffices'': Luks reduction, the descent to $G'$ and the descent to $\ker\varphi$ all enumerate
every coset \src{GNS}{459--470, 2097, 2101}. Their only pruning by automorphisms is find-symmetry, which
Lemma~\ref{lem:K} reproduces. GNS pair $\fA_{1,1}$ with each of the $r$ structures of $x_2$, which is also $r$
branches. Early exits of GNS (the two strings fall into different cases) correspond to computing one side
only, which costs the same.

\paragraph{Bounds for one macro step.}
\begin{itemize}
\item \emph{Nodes of $\cT$.} Each level branches by at most $d\cdot k^6\le d^7$, so $\cT$ has at most
  $d^{7(h+2)}=2^{O(\log^3d)}$ nodes.
\item \emph{Exits per node.} At most $2k^tn$ (certificates and colours), plus $k^6n$ (small windows of the
  children), plus, at leaves, $\max(d^{1+\log d},k!,d)\cdot n$. In total $2^{O(\log^2d)}n$.
\item \emph{Total exits.} $\sum n_i\le2^{K_0\log^3d}\,n$ with every $n_i\le n/2$, for an absolute constant $K_0$.
\item \emph{Non-recursive work.} $2^{O(\log^3d)}\cdot\big(n^{O(1)}+\exp(O((\log d)^{2c_{\mathrm{BC}}}))\big)
  =d^{O((\log d)^{c_1})}n^{O(1)}$.
\end{itemize}

\subsection{Proof of the time bound}\label{sec:timebound}

\begin{proof}[Proof of Theorem A (running time)]
\emph{Almost $d$-ary instances} with domain size $N$ (and $d\le N$). A macro step satisfies
Lemma~REC(A) with $K:=K_0\log^3d$, by \cite[Lemma~2.6]{GNS}. The semi-regular step satisfies (A).
The intransitive step satisfies (B), with $\ell\ge2$. So the recursion tree has at most $N^{K+1}$ leaves. Its
depth is at most $2\log N+1$: by the choice of chain rule, every child of a (B)-node is transitive, and a
transitive node is an (A)-node, so sizes halve at least every two levels. Hence there are at most
$(2\log N+2)N^{K+1}$ macro nodes, each costing $d^{O((\log d)^{c_1})}N^{O(1)}$. The total is
$N^{O((\log d)^{c_2})}$, since $d\le N$.

\emph{$\Gd$ instances.} A transitive constituent of size $n_i$ becomes an almost $d$-ary instance on
$N_i:=|\Omega_i^*|\le n_i^{O(\log^2d)}$ points \cite[Thm.~4.6]{GNS}, costing $n_i^{O((\log d)^{c_2+2})}$. The
constituents are disjoint and the bound is superadditive, so the total is $n^{O((\log d)^c)}$ with
$c:=c_2+2$. The coset rule and Lemma~R add polynomial overhead (preimages, \cite[Thm.~2.1]{GNS}).
\end{proof}

The only changes to GNS's analysis are the choice of Remark~\ref{rem:gap1} and the recursion lemma
(Lemma~REC).

\begin{remark}[The exponent]
We have not tried to optimize $c$, and we make no claim about its value here. A separate analysis of
the exponent in GNS and in the canonical version is work in progress.
\end{remark}

% =====================================================================
\section{Corollary B: the canonical Luks reduction}\label{sec:luks}

We follow \cite[\S4.3, Thm.~4.2]{BL83} and the ``fourth trick, due to Babai'' of \cite[\S4.2]{BL83}, a
half-edge encoding that keeps the string domain of polynomial size independently of $d$. We deviate in
two cosmetic ways: we pad every vertex to $d-1$ slots instead of first canonizing the valence string, and
we fold the edges inside the new layer into the same string.

\paragraph{Setting.} Let $X$ be a graph on $V$, $|V|=n$, of maximum degree $d\ge3$ (for $d\le2$ everything
is trivial), and put $q:=d-1$.

\paragraph{Step 0 (components).} Canonize each connected component $C$, obtaining $\mathrm{CF}(C)$ on
$[|C|]$ and its coset $\Lambda(C)$. $\mathrm{CF}(X)$ places the component forms on consecutive intervals of
$[n]$, sorted by (size, form). $\Lambda(X)$ consists of all bijections that map each component, by a member of
its own $\Lambda(C)$, onto an interval carrying an equal form. It is represented by one element together
with generators of $\prod\Aut(\mathrm{CF}(C))$ and the swaps of intervals with equal forms. This is the first
paragraph of \cite[\S4.3]{BL83}.

\paragraph{Step 1 (edge rooting).} For connected $X$ with at least two vertices and each \emph{ordered}
edge $(u,v)$ (at most $nd$ of them), compute $\mathrm{CF}(X,u,v)$ on $[n]$ with $u\mapsto1$, $v\mapsto2$, together
with $\Lambda(X,u,v)=\{\mu:X^\mu=\mathrm{CF}(X,u,v),u^\mu=1,v^\mu=2\}$. Then put $\mathrm{CF}(X):=\min\mathrm{CF}(X,u,v)$
over all ordered edges, and let $\Lambda(X)$ be the union of the $\Lambda(X,u,v)$ attaining the minimum. This is
Lemma~F4 with the equivariant family of ordered edges. The roots can be forgotten because they always
go to $1$ and $2$.

\paragraph{Layers.} Put $V_0:=\{u,v\}$ and $V_r:=\{w:\mathrm{dist}(w,\{u,v\})=r\}$ for $r\ge1$; let
$X_r:=X[V_0\cup\dots\cup V_r]$, $N_r:=|V_{\le r}|$, and $I_r:=(N_{r-1},N_r]$ (the labels of layer $r$). Every
$y\in V_r$ has at most $q$ neighbours in $V_{r+1}$ (its \emph{down-degree} $\deg_\downarrow y$): for $r\ge1$
because $y$ has a neighbour in $V_{r-1}$, and for $r=0$ because $u$ and $v$ are adjacent. This is the
``outvalence $\le d-1$'' of \cite[Lemma~4.4]{BL83}.

\paragraph{Invariant $(\mathrm I_r)$.} $\mathrm{CF}_r$ is a graph on $[N_r]$ that depends only on the isomorphism
type of $(X,u,v)$, and $\Lambda_r:=\{\lambda\colon V_{\le r}\to[N_r]:u\mapsto1,v\mapsto2,X_r^\lambda=\mathrm{CF}_r\}$ is
non-empty. Every such $\lambda$ maps $V_i$ onto $I_i$, because isomorphisms fixing the roots preserve the
distance to the roots, and distances in $X_r$ equal those in $X$. The group
$\hat A_r:=\Aut(\mathrm{CF}_r;1,2)$ (the \emph{doubly rooted} automorphism group) satisfies
$\Lambda_r=\lambda\hat A_r$. The base case is $\mathrm{CF}_0=$ the edge $\{1,2\}$ on $[2]$, $\Lambda_0=\{u\mapsto1,v\mapsto2\}$.

\paragraph{Step $r\to r+1$: the group.} Let $L:=I_r$ and let $\hat H:=S_q\wr_L\hat A_r$ act on
$\Xi_r:=[N_r]\sqcup(L\times[q])\sqcup(L\times[q])^2$: an element $(\hat a,s)$ acts by $\hat a$ on $[N_r]$, by
$(\ell,i)\mapsto(\ell^{\hat a},i^{s_\ell})$ on slots, and coordinatewise on pairs of slots. $\hat H$ is generated
by the generators of $\hat A_r$ and two generators of $S_q$ per $\ell\in L$, and it acts faithfully (this needs
the summand $[N_r]$, since $\hat A_r$ need not act faithfully on $L$). It lies in $\Gq{q}$ if $\hat A_r$ does,
being an extension of $S_q^L$ by $\hat A_r$. The domain has $|\Xi_r|\le n+nq+(nq)^2\le3n^4$ points and is ordered
with $[N_r]$ first, then slots lexicographically, then pairs lexicographically. $\hat H$ and the order of
$\Xi_r$ are functions of $\mathrm{CF}_r$ (as a set) and of $(N_r,L,q)$.

\paragraph{The string.} Given $\lambda\in\Lambda_r$ and a slot assignment $\beta$ (bijections
$\beta_y\colon N_\downarrow(y)\to[\deg_\downarrow y]$ for $y\in V_r$), put $\mathrm{slot}(y,w):=(y^\lambda,\beta_y(w))$ for each
down-edge $yw$, and define $x_{\lambda,\beta}$ on $\Xi_r$ by: $*$ on $[N_r]$; on a slot $(\ell,i)$, $1$ if it is
occupied ($i\le\deg_\downarrow(\ell^{\lambda^{-1}})$) and $0$ otherwise; on a pair of occupied slots with endpoints
$w,w'\in V_{r+1}$, the value ``$=$'' if $w=w'$, ``E'' if $ww'\in E(X)$, and ``$\ne$'' otherwise; $0$ if either
slot is unoccupied.

\begin{claim}[one orbit]\label{cl:B1}
All $x_{\lambda,\beta}$, for $\lambda\in\Lambda_r$ and all $\beta$, lie in one $\hat H$-orbit.
\end{claim}
\begin{proof}
Let $\lambda'=\lambda\hat a$ and let $\beta'$ be another slot assignment. For $\ell\in L$ put $y:=\ell^{\lambda^{-1}}$
and let $s_\ell$ map $\beta_y(w)$ to $\beta'_y(w)$, extended arbitrarily on unoccupied slots. Then $(\hat a,s)$
maps $\mathrm{slot}_{\lambda,\beta}(y,w)=(y^\lambda,\beta_y(w))$ to $(y^{\lambda\hat a},\beta'_y(w))=\mathrm{slot}_{\lambda',\beta'}(y,w)$, so
it preserves occupancy and the pair values.
\end{proof}

\begin{claim}[invariance]\label{cl:B2}
Let $\pi\colon(X_{r+1},u,v)\to(X'_{r+1},u',v')$ be an isomorphism of rooted graphs (here $X_{r+1}$ is the
subgraph induced on the first $r+1$ layers). Then $x_{\pi^{-1}\lambda,\beta\circ\pi^{-1}}(X')=x_{\lambda,\beta}(X)$, and
$\hat H$ depends only on $\mathrm{CF}_r$. In particular the orbit $\cO_r$ is an isomorphism invariant of
$(X,u,v)$; in particular the first statement applies to automorphisms of $(X_{r+1},u,v)$.
\end{claim}
\begin{proof}
$x_{\lambda,\beta}$ depends only on $(X_{r+1},u,v)$, $\lambda$ and $\beta$: the layer $V_{r+1}$, the down-edges and the
edges inside $V_{r+1}$ are read off $X_{r+1}$. Transporting everything along $\pi$ gives the claim.
\end{proof}

\paragraph{Canonization.} Let $\hat F:=F_{\hat H}(x_{\lambda,\beta})$ and $\hat\Lambda:=\Lambda_{\hat H}(x_{\lambda,\beta})$,
computed by \emph{one call of Theorem~A} with group $\hat H\in\Gq{d-1}$ on at most $3n^4$ points. $\hat F$
depends only on $\cO_r$.

\paragraph{Reading off $\mathrm{CF}_{r+1}$.} In $\hat F$, ``$=$'' is an equivalence relation on occupied slots (it
is one in $x_{\lambda,\beta}$, and $\hat F=x_{\lambda,\beta}^{\hat h}$). Order its classes by least slot and label them
$N_r+1,\dots,N_{r+1}$. Let $\mathrm{CF}_{r+1}$ be $\mathrm{CF}_r$ together with an edge $\{\ell,N_r+j\}$ for each occupied
slot $(\ell,i)$ in class $j$, and an edge $\{N_r+j,N_r+j'\}$ whenever $\hat F$ has ``E'' on a pair of slots from
classes $j$ and $j'$ (this value does not depend on the choice of slots, since it depends only on the
endpoints). So $\mathrm{CF}_{r+1}$ is a function of $(\mathrm{CF}_r,\hat F)$, hence an invariant.

\paragraph{The labelings.} For $\hat h=(\hat a,s)\in\hat\Lambda$ define $\mu=\mu(\lambda,\beta,\hat h)$ by
$y\mapsto y^{\lambda\hat a}$ on $V_{\le r}$, and on $V_{r+1}$ by $w\mapsto N_r+j$, where $j$ is the class of $\hat F$
containing $\hat h(\mathrm{slot}(y,w))$ for any up-neighbour $y$ of $w$. This is well defined, because the slots of
$w$ are ``$=$''-related and every $w$ has an up-neighbour, and it is a bijection onto $I_{r+1}$.

\begin{claim}\label{cl:B3}
$X_{r+1}^\mu=\mathrm{CF}_{r+1}$.
\end{claim}
\begin{proof}
$X_r^{\lambda\hat a}=\mathrm{CF}_r^{\hat a}=\mathrm{CF}_r$. For $y\in V_r$, $w\in V_{r+1}$: $yw\in E$ iff $\mathrm{slot}(y,w)$ is
occupied with endpoint $w$, iff its $\hat h$-image (a slot of $y^\mu$) lies in class $w^\mu$. This uses that
$x_{\lambda,\beta}$ and $\hat F$ agree along $\hat h$, and that the slots of a fixed $y$ have distinct endpoints
($\beta_y$ is a bijection). No edge joins $V_{r+1}$ to $V_{\le r-1}$. Edges inside $V_{r+1}$ are read off the ``E''
values in the same way.
\end{proof}

\begin{claim}\label{cl:B4}
$\{\mu(\lambda,\beta,\hat h)\}=\Lambda_{r+1}:=\{\mu\text{ rooted}:X_{r+1}^\mu=\mathrm{CF}_{r+1}\}$.
\end{claim}
\begin{proof}
``$\subseteq$'' is Claim~\ref{cl:B3}. For ``$\supseteq$'': by Claim~\ref{cl:B2} applied to an automorphism $g$ of
$(X_{r+1},u,v)$, the choice $(g^{-1}\lambda,\beta\circ g^{-1},\hat h)$ is valid, satisfies
$x_{g^{-1}\lambda,\beta\circ g^{-1}}=x_{\lambda,\beta}$, and yields $g^{-1}\mu$. So the set of outputs is closed under
$\Aut(X_{r+1},u,v)$; $\Lambda_{r+1}$ is a single coset of that group, and it contains an output.
\end{proof}

\begin{claim}[generators]\label{cl:B5}
Let $\theta\colon\Aut_{\hat H}(\hat F)\to\Sym([N_{r+1}])$ act by $\hat a$ on $[N_r]$ and by the induced permutation
of ``$=$''-classes on the new labels. Then $\theta(\Aut_{\hat H}(\hat F))=\hat A_{r+1}$.
\end{claim}
\begin{proof}
``$\subseteq$'' is clear. For ``$\supseteq$'', let $b\in\hat A_{r+1}$. It fixes $1,2$ and preserves layers, so it preserves
$[N_r]$; put $\hat a:=b|_{[N_r]}\in\hat A_r$. For $\ell\in L$, the classes adjacent to $\ell$ correspond bijectively
to the occupied slots of $\ell$, and $b$ maps them to the classes adjacent to $\ell^b$. Let $s_\ell$ follow this
bijection, extended arbitrarily on unoccupied slots ($b$ preserves down-degrees). Then $(\hat a,s)$ fixes
$\hat F$ and $\theta(\hat a,s)=b$.
\end{proof}

So $\Lambda_{r+1}=\mu_0\hat A_{r+1}$ is obtained in polynomial time from $\hat\Lambda$, and by induction
$(\mathrm I_r)$ holds for all $r$; at $r=R$, the eccentricity of $\{u,v\}$, we have $X_R=X$. By induction
$\hat A_r\in\Gq{d-1}$: $\hat A_0=1$, and $\hat A_{r+1}$ is a quotient of a subgroup of $\hat H\in\Gq{d-1}$. This reproves
\cite[Lemma~4.3]{BL83} (due to Luks) inside the construction.

\begin{proof}[Proof of Corollary B]
There are at most $n^2$ rooted edges and at most $n$ layers, so at most $n^3$ calls of Theorem~A on
domains of size at most $3n^4$ with groups in $\Gq{d-1}$, each costing $(3n^4)^{O((\log d)^c)}$, plus
polynomial overhead. The total is $n^{O((\log d)^c)}$. As in \cite[Thm.~8.1]{GNS}, the dependence on $d$
comes only from Theorem~A.
\end{proof}
\status{Proof modulo Theorem A. In addition, an independent implementation of this construction, with
brute-force canonization in place of Theorem~A, was checked on 47 small connected graphs ($n\le8$,
$d\le4$, every ordered edge and every layer: Claims~\ref{cl:B1}, \ref{cl:B3}, \ref{cl:B4}, \ref{cl:B5} and the
minimum over roots) and on six disconnected graphs (component rule). This is a sanity check only.}

\paragraph{Comparison with \cite{BL83}.} \cite[\S4.3]{BL83} alternates two steps: canonize $Y_r$ (the edges
inside $V_r$) relative to the current coset, then extend through the bipartite graph $Z_r$ between
$V(X_r)$ and $V_{r+1}$ by the algorithm of \cite[\S4.2]{BL83}. Trick~4 there canonizes the valence string,
extends the coset to the edges ``allowing arbitrary permutations of edges with common origin'',
canonizes $y(e_1,e_2)=[e_1,e_2\text{ end at the same vertex}]$ on $E\times E$, and orders the new vertices by
their least incident edge. Our padded slots replace the valence step (the group $S_q^L$ is larger than
$\prod S_{\deg_\downarrow}$, still in $\Gq{d-1}$, and Claim~\ref{cl:B1} shows the extra freedom is harmless), and our
pair values ``E'' replace $Y_{r+1}$.

% =====================================================================
\section{Corollary D and relational structures}\label{sec:D}

\begin{proof}[Proof of Corollary D]
Let $x$ be a string on $V^2$ (e.g.\ the adjacency string of a graph on $V$) and $\Delta\rho$ a labeling coset with
$\Delta\in\Gd$. By Remark~\ref{rem:labcosets}, $\mathrm{CL}(x,\Delta\rho):=\rho\Lambda_H(x^\rho)$ with
$H:=\rho^{-1}\Delta\rho\le\Sym([n])$. The induced action of $H$ on $[n]^2$ is faithful and isomorphic to $H$, so it
lies in $\Gd$. Theorem~A on $n^2$ points gives the bound $|V|^{O((\log d)^c)}$.
\end{proof}
\status{Proof modulo Theorem A.}

Corollary D is the canonical counterpart of \cite[Thm.~31]{Wie20}, which computes
$\Iso(G_1;G_2)\cap\Delta\varphi$ in time $|V|^{\polylog(\cw\Delta)}$. Theorem~A with $\Sym(D)$ acting on $D^t$ also
gives canonical forms of $t$-ary structures and hypergraphs in time $n^{O(t(\log n)^c)}$, the canonical
counterparts of \cite[Thm.~8.3, Cor.~8.4]{GNS}. These are weaker than known canonization results
\cite[Thms.~5.2, 5.4]{Wie21}, so we claim nothing new there.

% =====================================================================
\section{Theorem H: canonical Neuen 2020}\label{sec:H}

Throughout this section, ``Neu20'' refers to arXiv:2002.06997v2, and result numbers were computed from its
source. Neuen's results are used as listed in Section~\ref{sec:assumptions}.

\subsection{Neuen's algorithm}

Neuen solves \emph{generalized string isomorphism}: a group $\Gamma$, a $\Gamma$-invariant partition $\fP$, and
sets $\fX,\mathfrak Y$ of $\fP$-strings, i.e.\ pairs $(P,x)$ with $P\in\fP$ and $x\colon P\to\Sigma$. Hypergraphs and
sets of strings reduce to it equivariantly \src{Neu20}{987--1005}. Progress is measured by the
\emph{$d$-virtual size} $s=\sum_{P\in\fP}|P|\cdot m_{\fX}(P)^{f(d)+1}$ with $f(d)=O(\log d)$, which is at most
$(n+m)^{O(\log d)}$ \src{Neu20}{1570--1578}. An instance is normalized once by
\cite[Thm.~4.8]{Neu20}, which produces an equivalent instance with an almost $d$-ary sequence, built from
structure graphs. The main algorithm \cite[Lemma~4.36, Thm.~4.37]{Neu20} proceeds as follows:
intransitive groups are processed orbit by orbit and combined by \cite[Lemma~4.15]{Neu20}
(CombineWindows); a semi-regular top is handled by Luks reduction with BalanceOrbits
\cite[Lemma~4.3]{Neu20}; otherwise a giant representation \cite[Lemma~4.35]{Neu20}, the local
certificates routine \cite[Thm.~4.24]{Neu20}, aggregation \cite[Lemma~4.29]{Neu20}, find-structure and
find-symmetry \cite[Lemmas~4.33, 4.34]{Neu20}, and recursion on the potential $k$.

The difference from GNS lies inside the local certificates loop. There, the window automorphisms
$\Gamma_{(\Omega\setminus W)}$ are automorphisms of a set of strings only if $\fX[W]$ is \emph{simple}
\cite[Lemma~4.16]{Neu20}. So after each window step the instance is simplified \cite[Thm.~4.17]{Neu20}
(copies of each block, one per restriction to $W$) and renormalized \cite[Thm.~4.9]{Neu20} (implanted
structure graphs). Neuen notes that renormalization ``is not isomorphism-invariant''
\src{Neu20}{1930}, and therefore runs the local certificates routine on all pairs of test objects
simultaneously \src{Neu20}{270--275}.

\paragraph{Sets of strings as objects.} For $U\subseteq\Omega$ put
$\fP[U]:=\{P\cap U:P\in\fP,\ P\cap U\ne\emptyset\}$, and for a set $Z$ of $\fP$-strings let
$Z[U]:=\{(P\cap U,x[P\cap U]):(P,x)\in Z,\ P\cap U\ne\emptyset\}$, a set of $\fP[U]$-strings. A permutation acts by
$Z^g:=\{(P^g,x^g):(P,x)\in Z\}$. We call $Z$ \emph{completely occupied} if $m_Z(P)\ge1$ for all blocks $P$
of its partition; this is Neuen's standing assumption \src{Neu20}{1036}, it is inherited by every
$Z[U]$, and we need it for \cite[Obs.~4.12]{Neu20}: if $Z$ is completely occupied and $U\subsetneq U'$, then the
$d$-virtual size of $Z[U]$ is strictly smaller than that of $Z[U']$. Sets of strings on a fixed $\Omega$ are
compared in a fixed total order (the sorted list of the strings, compared along the order of $\Omega$). The
notions of canonizer (Definition~\ref{def:canonizer}), Lemma~\ref{lem:canon-basic}, the coset rule, and the
statements and proofs of Lemmas F3, F4, R, S, Lemma~\ref{lem:K} and Proposition~\ref{prop:opt1} use only a
right group action and this total order on the objects, so they apply verbatim to sets of strings in place
of strings.

\subsection{Equivariant coset maps}

\begin{definition}\label{def:ECM}
Let $H\le\Sym(\Omega)$ and let $\cZ$ be an $H$-invariant set of objects. An \emph{$H$-equivariant coset map}
(ECM) assigns to each $Z\in\cZ$ a left coset $C(Z)=\pi R(Z)\subseteq H$ of a subgroup $R(Z)\le H$ such that
$C(Z^h)=h^{-1}C(Z)$ for all $h\in H$. We call $R(Z)$ the \emph{right group} of $C$ at $Z$ (the analogue of
the right subgroup of a canonizer).
\end{definition}
The $\Lambda_H$ of a canonizer is an ECM with $R=\Aut_H(F_H(x))$. ECMs are the weak notion of canonical
labeling used by Schweitzer and Wiebking.

\begin{named}[Lemma E1]
For an $H$-ECM, $R(Z^h)=R(Z)$, and $\Aut_H(Z)\le C(Z)C(Z)^{-1}=\pi R(Z)\pi^{-1}$ for any $\pi\in C(Z)$.
\end{named}
\begin{proof}
$h^{-1}C(Z)$ has right group $R(Z)$. If $a\in\Aut_H(Z)$ then $C(Z)=C(Z^a)=a^{-1}C(Z)$, so $a\in C(Z)C(Z)^{-1}$.
\end{proof}

\begin{named}[Lemma E2 (composition at a prefix)]
Let $C_1$ be an $H$-ECM on $\cZ$ with right group $R_1(Z)$. A \emph{prefix} is a value $\bar p(Z)$ that is
computed from $Z^\pi$ and is the same for every $\pi\in C_1(Z)$, satisfies $\bar p(Z^h)=\bar p(Z)$, and
determines $R_1(Z)$. For each prefix value $p$ with right group $R$, the set
$\cZ_p:=\{Z^\pi:\pi\in C_1(Z),\ \bar p(Z)=p\}$ is $R$-invariant; let $C_2^p$ be an $R$-ECM on $\cZ_p$. Put
$C(Z):=\pi\cdot C_2^{\bar p(Z)}(Z^\pi)$ for any $\pi\in C_1(Z)$. Then $C$ is a well-defined $H$-ECM with
$C(Z)\subseteq C_1(Z)$.
\end{named}
\begin{proof}
For $\pi'=\pi r$ ($r\in R$): $\pi r\cdot C_2(Z^{\pi r})=\pi r\cdot r^{-1}C_2(Z^\pi)$. For $Z^h$ use
$\pi'=h^{-1}\pi\in C_1(Z^h)$; then $Z^{h\pi'}=Z^\pi$ and $\bar p$ is unchanged, so $C(Z^h)=h^{-1}C(Z)$.
\end{proof}
When $C_1$ and $C_2$ come from canonizers, this is Lemma~F2. The point of E2 is that \emph{anything computed
from the prefix by a deterministic procedure, for example a renormalization, is identical for every object
in the $H$-orbit.} This replaces Neuen's synchronization of pairs.

\begin{named}[Lemma E3 (change of action at a prefix)]
In the setting of E2, suppose the prefix $p$ determines an injective homomorphism
$\iota_p\colon R\to\Sym(\Omega_p)$ and an injective map $Z'\mapsto Z'_*$ on $\cZ_p$ with
$(Z'^r)_*=(Z'_*)^{\iota_p(r)}$. If $C^*$ is an $\iota_p(R)$-ECM on $\{Z'_*:Z'\in\cZ_p\}$, then
$C_2(Z'):=\iota_p^{-1}(C^*(Z'_*))$ is an $R$-ECM on $\cZ_p$.
\end{named}
\begin{proof}
$\iota_p^{-1}(C^*((Z'^r)_*))=\iota_p^{-1}(\iota_p(r)^{-1}C^*(Z'_*))=r^{-1}C_2(Z')$.
\end{proof}

\begin{named}[Lemma E4 (from an ECM to a canonizer)]
Let $C$ be an $H$-ECM on an $H$-invariant set $\cZ$ with right group $R(Z)$. Suppose that for every value $R$ of
the right group an $R$-canonizer $(F_R,\Lambda_R)$ is given on the $R$-invariant set
$\{Z^\pi:\pi\in C(Z),\ R(Z)=R\}$. Then
\[ F_H(Z):=F_{R(Z)}(Z^\pi),\qquad \Lambda_H(Z):=\pi\,\Lambda_{R(Z)}(Z^\pi)\qquad(\pi\in C(Z)) \]
defines an $H$-canonizer on $\cZ$. Its right subgroup is $\Aut_{R(Z)}(F_H(Z))$, and $\Lambda_H(Z)\subseteq C(Z)$.
\end{named}
\begin{proof}
\emph{Independence of $\pi$.} Let $\pi'=\pi r$ with $r\in R$. Then $F_R(Z^{\pi r})=F_R(Z^\pi)$, and
$\pi r\,\Lambda_R(Z^{\pi r})=\pi r\,r^{-1}\Lambda_R(Z^\pi)$.
\emph{(C2).} For $Z^h$ use $h^{-1}\pi\in C(Z^h)$; then $Z^{hh^{-1}\pi}=Z^\pi$, and $R(Z^h)=R(Z)$ by Lemma~E1.
\emph{(C1).} Every element of $\Lambda_H(Z)$ maps $Z$ to $F:=F_H(Z)$. Conversely, let $Z^h=F$. Write $F=Z^{\pi\lambda}$
with $\lambda\in\Lambda_R(Z^\pi)$. Then $C(F)=(\pi\lambda)^{-1}C(Z)=\lambda^{-1}R=R$, so $1\in C(F)$; on the other hand
$C(F)=C(Z^h)=h^{-1}C(Z)$, so $h\in C(Z)$, say $h=\pi r$ with $r\in R$. Then $(Z^\pi)^r=F$, so $r\in\Lambda_R(Z^\pi)$.
\emph{Right subgroup.} If $h\in\Aut_H(F)$, the same computation gives $h\in C(F)=R$.
\end{proof}
This is Lemma~S in the language of ECMs: an ECM combined with canonizers at the right group gives a
canonizer.

\subsection{Ordered colours from a setwise reference test set}

This is Lemma~T moved from the pointwise stabilizer of a reference \emph{tuple} to the setwise stabilizer
of a reference \emph{set}, which is the form of Wiebking's canonical test set \cite[pp.~102--103]{Wie21}
(``$T^{\mathrm{can}}:=\{1,\dots,m\}\subseteq W^{\mathrm{can}}$''; the group ``stabilize[s] the test set
$T^{\mathrm{can}}$ setwise'').

\begin{named}[Lemma T$^{\text{set}}$]
Let $H\le\Sym(\Omega)$ and $\varphi\colon H\to\Sym([k])$ with $\varphi(H)\ge A_k$ and $k\ge t+2$. Let
$T_0:=\{1,\dots,t\}$ (P1), $H_0:=\varphi^{-1}(\mathrm{Stab}(T_0))$, and $\psi\colon H_0\to\Sym(T_0)$ the induced action.
Let $C$ be an $H_0$-ECM with right group $R$ on an $H$-invariant set $\cZ$. For $Z\in\cZ$ and an ordered tuple
$\vT$ of distinct points with underlying set $T$, choose $\sigma\in H$ with $T^{\varphi(\sigma)}=T_0$ and put
$E(Z,T):=\sigma C(Z^\sigma)$, $R(Z,T):=R(Z^\sigma)$, and
\[ c(Z,\vT):=\big(R(Z^\sigma),\ \text{the }\psi(R(Z^\sigma))\text{-orbit of }\vT^{\varphi(e)}\big)\quad\text{for any }e\in E(Z,T). \]
Then:
\begin{enumlab}{\alph}
\item $E$, $R$ and $c$ are well defined;
\item $c(Z^h,\vT^{\varphi(h)})=c(Z,\vT)$ for $h\in H$;
\item if $c(Z,\vT^\pi)=c(Z,\vT)$ for all $\pi\in\Sym(T)$, then $\psi(R(Z^\sigma))=\Sym(T_0)$;
\item $\psi_T(\Aut_{H_T}(Z))$ is conjugate to a subgroup of $\psi(R(Z^\sigma))$, where $H_T:=\varphi^{-1}(\mathrm{Stab}(T))$.
\end{enumlab}
\end{named}
\begin{proof}
(a) If $\sigma'=\sigma g$ with $g\in H_0$, then $\sigma g\cdot C(Z^{\sigma g})=\sigma g\cdot g^{-1}C(Z^\sigma)$, and $R$ is
unchanged by Lemma~E1. Two elements of $E$ differ by $r\in R$ on the right, and
$\vT^{\varphi(er)}=(\vT^{\varphi(e)})^{\psi(r)}$. (b) Use $h^{-1}\sigma$ for $(Z^h,T^{\varphi(h)})$: the transported object is
again $Z^\sigma$, $E$ becomes $h^{-1}E$, and $(\vT^{\varphi(h)})^{\varphi(h^{-1}e)}=\vT^{\varphi(e)}$. (c) With $\sigma$ and $e$ fixed
(they depend on the set $T$ only), $(\vT^\pi)^{\varphi(e)}$ runs through all $t!$ orderings of $T_0$, which then lie
in one $\psi(R)$-orbit. (d) Lemma~E1 for $H_0$ and $Z^\sigma$, conjugated by $\sigma$.
\end{proof}
Colours are compared through canonical generating sets of $R$ and the lexicographically least ordering in
the orbit. Only (b) and (c) are used below. In particular, the colours need \emph{not} reproduce Neuen's
comparison classes.

\subsection{The pairwise steps of Neu20}

An internal reading of Neu20 identified 21 choice points or pairwise steps. Four are pairwise in an
essential way:
\begin{enumlab}{\alph}
\item the local certificates routine, run ``on all pairs in parallel'' with an arbitrary test object and a
  realignment \src{Neu20}{2201--2258};
\item the representative group $\Gamma''_1$ in Simplify \src{Neu20}{2053--2061};
\item the joint category on $D'_1\sqcup D'_2$ in the aggregation \src{Neu20}{2589--2596};
\item find-structure and find-symmetry \src{Neu20}{2603--2653}.
\end{enumlab}
(c) and (d) are as in GNS and are handled by Lemma~T$^{\text{set}}$, Lemma~S with BC, and Lemma~\ref{lem:K}. (a) and (b)
are new. The others are handled as in Section~\ref{sec:algo}: (P1) choices (structure graphs and their
transporters, the choice of $v_P$, the minimal block system, the giant representation), descents
(F3), equivariant branch families (F4), equivariant encodings (Lemma~R), and the two routine pairwise
subroutines CombineWindows and BalanceOrbits (Lemmas~W and~B below). All group-only constructions are
subject to requirement~\ref{D3}: they must be deterministic functions of the group and the ordered domain,
computed from canonical generating sets.

\subsection{Lemmas W, F2\texorpdfstring{$^{\text{set}}$}{-set} and B}

Theorem~A is a canonizer for strings under a group. Below it is also applied to strings on a set $\Gamma$ on
which the current group $A$ acts through a homomorphism $\theta\colon A\to\Sym(\Gamma)$. This is Lemma~S with
$S(x)$ the string and $\Lambda'$ the Theorem~A canonizer under $\theta(A)\in\Gd$ (quotients stay in $\Gd$).

\begin{named}[Lemma W (canonical CombineWindows)]
Let $A\le\Sym(\Omega)$ be in $\Gd$, let $\fP$ be $A$-invariant, let $W_1,W_2\subseteq\Omega$ be disjoint and $A$-invariant, let
$F_j$ be a set of $\fP[W_j]$-strings with $A[W_j]\le\Aut(F_j)$ ($j=1,2$), and put $W:=W_1\cup W_2$. Let $\cZ$ be the
$A$-invariant set of sets $Z$ of $\fP$-strings with $Z[W_j]=F_j$ for $j=1,2$. Let $V:=F_1\sqcup F_2$, let $\theta\colon
A\to\Sym(V)$ be the induced action, and for $Z\in\cZ$ let the \emph{linking graph} $G(Z)$ be the graph on $V$ with
edges $\{(P\cap W_1,x[W_1]),(P\cap W_2,x[W_2])\}$ for $(P,x)\in Z$ with $P\cap W_1\ne\emptyset\ne P\cap W_2$
\src{Neu20}{1648--1651}. Let $(F_G,\Lambda_G)$ be the $\theta(A)$-canonizer on these graphs given by Corollary~D
under $\theta(A)\in\Gd$. Then there is a set $F_W(Z)$ of $\fP[W]$-strings, depending only on $V$ and
$F_G(G(Z))$, such that for $Z\in\cZ$ and $E(Z):=\theta^{-1}(\Lambda_G(G(Z)))$, a left coset $\sigma R$ with
$R=\theta^{-1}(\Aut_{\theta(A)}(F_G(G(Z))))$,
\[ E(Z)=\{h\in A: Z^h[W]=F_W(Z)\},\qquad F_W(Z^a)=F_W(Z)\quad\text{for }a\in A. \]
Thus $(F_W,E)$ is an $A$-canonizer of $Z\mapsto Z[W]$ on $\cZ$. Cost: one call of Theorem~A on a graph with at most
$2|Z[W]|$ vertices, i.e.\ on a string on at most $(2|Z[W]|)^2$ points (through $\theta$).
\end{named}
\begin{proof}
For $h\in A$ we have $F_j^h=F_j$, so $\theta(h)$ permutes $V$, $Z^h\in\cZ$, and $G(Z^h)=G(Z)^{\theta(h)}$. A set $Z\in\cZ$
is recovered on $W$ from its linking graph: $Z[W]$ consists of the strings $(Q_1\cup Q_2,z_1\cup z_2)$ over the edges
$\{(Q_1,z_1),(Q_2,z_2)\}$, together with the strings of $F_1\cup F_2$ on blocks that meet only one of $W_1,W_2$;
and $G(Z)$ is determined by $Z[W]$ and $V$. So $Z[W]\leftrightarrow G(Z)$ is a bijection; let $F_W(Z)$ be the set of
strings corresponding to the graph $F_G(G(Z))$. For $h\in A$ this gives $Z^h[W]=F_W(Z)$ iff $G(Z)^{\theta(h)}=F_G(G(Z))$
iff $\theta(h)\in\Lambda_G(G(Z))$ iff $h\in E(Z)$. The remaining statements follow from (C1), (C2) for
$(F_G,\Lambda_G)$. (This is the forward and backward direction of the proof of \cite[Lemma~4.15]{Neu20}.)
\end{proof}

\begin{named}[Lemma F2$^{\text{set}}$ (one step of the chain rule for sets of strings)]
Let $A\le\Sym(\Omega)$ be in $\Gd$, let $\fP$ be $A$-invariant, and let $\cZ$ be an $A$-invariant set of sets of $\fP$-strings. Let
$C$ be an $A$-ECM on $\cZ$ with right group $R$, and suppose that there is a prefix $(U,F,R)$, constant on $A$-orbits,
where $U$ is $R$-invariant and $Z^\pi[U]=F$ for all $\pi\in C(Z)$ (so $R[U]\le\Aut(F)$). Let $V\subseteq\Omega\setminus U$ be $R$-invariant,
given by a function of the prefix, and let $(F_V,\Lambda_V)$ be an $R$-canonizer on the $R$-invariant set
$\{Z^\pi[V]:\pi\in C(Z)\}$. For $\pi\in C(Z)$ put $w:=Z^\pi$, take $\lambda\in\Lambda_V(w[V])$, put $y:=w^\lambda$ and
$R_2:=\Aut_R(F_V(w[V]))$, and apply Lemma~W to $R_2$, $W_1:=U$, $W_2:=V$, $F_1:=F$, $F_2:=F_V(w[V])$ and the
object $y$, obtaining $E(y)=\sigma R_3$ and $F_+:=F_W(y)$. Then $C_+(Z):=\pi\lambda E(y)$ is an $A$-ECM on $\cZ$ with right group
$R_3$ and $C_+(Z)\subseteq C(Z)$; and $(U\cup V,F_+,R_3)$ is a prefix of the same kind: it is constant on $A$-orbits,
and $Z^{\pi'}[U\cup V]=F_+$ for all $\pi'\in C_+(Z)$.
\end{named}
\begin{proof}
Since $V\cap U=\emptyset$, we have $y[U]=F$ and $y[V]=F_V(w[V])$, and $R_2$ preserves both $F$ and $F_V(w[V])$, so Lemma~W
applies. By Lemma~E2 with $C_1:=C$ (prefix $(U,F,R)$) it suffices that $w\mapsto\lambda E(w^\lambda)$ is an $R$-ECM on
$\{Z^\pi\}$. This is again Lemma~E2, now with $C_1:=(w\mapsto\Lambda_V(w[V]))$, an $R$-ECM by
Lemma~\ref{lem:canon-basic}(1) and $(w^r)[V]=(w[V])^r$, with prefix $(F_V(w[V]),R_2)$, and with the
$R_2$-ECM $y\mapsto E(y)$ of Lemma~W. For $\pi'=\pi\lambda e$ with $e\in E(y)$ we get $Z^{\pi'}[U\cup V]=y^e[U\cup V]=F_+$ by Lemma~W. Replacing
$(\pi,\lambda)$ by another choice replaces $y$ by an element of $y^{R_2}$, and $Z$ by $Z^h$ leaves $y$ unchanged, so
$F_+=F_W(y)$ and $R_3$ are constant on $A$-orbits by Lemma~W.
\end{proof}
This is the analogue of Lemma~F2 for sets of strings; Lemma~W replaces the concatenation of canonical
forms. The chain rule itself is the iteration of this step in the subroutine Chain of \S\ref{sec:canhyp}.

\begin{named}[Lemma B (canonical BalanceOrbits)]
Let $A\le\Sym(\Omega)$ be in $\Gd$, let $\fP$ be $A$-invariant, and let $\cZ$ be an $A$-invariant set of sets of $\fP$-strings. There is
an $A$-ECM $C_{\mathrm B}$ on $\cZ$, with right group $R$, such that for every $Z\in\cZ$, every $\pi\in C_{\mathrm B}(Z)$ and every
$R$-orbit $O$, the set $Z^\pi[O]$ is balanced. It is computed by at most $n\log n$ calls of Theorem~A on at most
$n$ points, with $n=|\Omega|$. (This is \cite[Lemma~4.3]{Neu20} made canonical.)
\end{named}
\begin{proof}
Start with $C_0(Z):=A$, $R_0:=A$. Suppose $C_i$ is an $A$-ECM with right group $R_i$; let $\pi\in C_i(Z)$ and $z:=Z^\pi$. For
$r\in R_i$ and an $R_i$-orbit $O$ we have $z^r[O]=(z[O])^r$, so whether $z[O]$ is balanced depends only on the
$R_i$-orbit of $z$. If all $R_i$-orbits are balanced, stop. Otherwise let $O$ be the first $R_i$-orbit in the order
(size, least point) with $z[O]$ unbalanced; it depends only on $R_i$ and the $R_i$-orbit of $z$. Let
$\mu\colon\fP[O]\to\mathbb N$, $Q\mapsto m_{z[O]}(Q)$, and let $\rho\colon R_i\to R_i[\fP[O]]=:B\in\Gd$ be the natural
homomorphism. Then $B$ is transitive on $\fP[O]$ and $\mu$ is not constant. Let $(F_\mu,\Lambda_\mu)$ be the Theorem~A
canonizer of $\mu$ under $B$ and $E(z):=\rho^{-1}(\Lambda_\mu(\mu))$, a coset $\sigma R_{i+1}$ with
$R_{i+1}:=\rho^{-1}(\Aut_B(F_\mu))$. Since $\mu_{z^r}=\mu_z^{\rho(r)}$, we have $E(z^r)=r^{-1}E(z)$, so $E$ is an $R_i$-ECM and
Lemma~E2 (with prefix $R_i$) makes $C_{i+1}(Z):=\pi E(z)$ an $A$-ECM with right group $R_{i+1}$. Moreover $R_{i+1}$ is a
proper subgroup of $R_i$, since $F_\mu$ is not constant and $B$ is transitive. A subgroup chain in $\Sym(\Omega)$ has length
at most $\log_2 n!\le n\log n$, which bounds the number of rounds.
\end{proof}

\subsection{Canonical simplification and renormalization}

\begin{named}[Lemma SIMP]
Let $A\le\Sym(\Omega)$ be in $\Gd$, with an almost $d$-ary $A$-invariant sequence $\vec\fB$ having $\fP$ as a level; let $W\subseteq\Omega$ be
$A$-invariant, let $F$ be a set of $\fP[W]$-strings with $A[W]\le\Aut(F)$, and let $\cZ_F$ be the $A$-invariant set of
completely occupied sets $Z$ of $\fP$-strings with $Z[W]=F$. From these data one computes an $A$-ECM $C$ on $\cZ_F$ with
right group $A''\le A$ (by one call of Theorem~A) and, from $(A,\fP,\vec\fB,W,F)$ and the canonical string of step~3
below, a domain $\Omega^*$, a partition $\fP^*$ with a sequence $\vec\fB^*$, a window $W^*$, a monomorphism
$\iota\colon A''\to\Sym(\Omega^*)$ with image $\Lambda^*$, and, for each $\pi\in C(Z)$, a completely
occupied set $(Z^\pi)_*$ of $\fP^*$-strings, such that
\begin{enumlab}[S]{\arabic}
\item $\vec\fB^*$ is an almost $d$-ary sequence of $\Lambda^*$-invariant partitions having $\fP^*$ as a level;
\item $W^*$ is $\Lambda^*$-invariant, $\Lambda^*[W^*]\le\Aut(F^*)$, where $F^*:=(Z^\pi)_*[W^*]$ is simple and depends only on the
  computed data, not on $Z$ or $\pi$;
\item $Z^\pi\mapsto(Z^\pi)_*$ is injective and $(Z^{\pi r})_*=((Z^\pi)_*)^{\iota(r)}$ for $r\in A''$;
\item the $d$-virtual size of $(Z^\pi)_*$ is at most that of $Z$, and that of $(Z^\pi)_*[\Omega^*\setminus W^*]$ is at most
  that of $Z[\Omega\setminus W]$.
\end{enumlab}
If $F$ is simple, put $C:=A$, $\Omega^*:=\Omega$, $(Z^\pi)_*:=Z^\pi$ (nothing is done).
\end{named}
\begin{proof}
We follow the proof of \cite[Thm.~4.17]{Neu20} for a single object.

\emph{Step 1 (copy blocks).} $\Omega':=\bigcup_PP\times\{z:(P\cap W,z)\in F\}$, where a block $P$ with $P\cap W=\emptyset$ gets the
single copy $P\times\{\emptyset\}$ \src{Neu20}{1974--1985}; since $Z$ is completely occupied, $\Omega'$ depends only on $\fP$ and $F$.
Since $A[W]\le\Aut(F)$, the rule $(\alpha,z)^\gamma:=(\alpha^\gamma,z^\gamma)$ defines an action $\psi'\colon A\to\Sym(\Omega')$, faithful because
every point of $\Omega$ occurs as a first coordinate. The partition $\fP'$, the copies $Z\mapsto Z'$, and $W':=W\times(\cdot)$ are as in
\src{Neu20}{1974--2003}: $(P,x)\mapsto(P\times\{x[W\cap P]\},(\alpha,x[W\cap P])\mapsto x(\alpha))$, so $Z'$ determines $Z$ and $(Z^\gamma)'=(Z')^{\psi'(\gamma)}$
(claim \texttt{preserve-isomorphisms}, \src{Neu20}{1987--2002}). Moreover $Z'$ is completely occupied, $Z'[W']$ depends only on $F$,
$\psi'(A)[W']\le\Aut(Z'[W'])$ and $Z'[W']$ is simple \src{Neu20}{2006--2029}.

\emph{Step 2 (lift the sequence)} to $\vec\fB'$ \src{Neu20}{2031--2040}, a function of $(\vec\fB,F)$; only the level $\fP'$ can violate almost
$d$-arity \src{Neu20}{2041--2046}.

\emph{Step 3 (balance).} Apply Lemma~S to the string $x(Z'):\fP'\to\mathbb N^2$,
$P'\mapsto(m_{Z'}(P'),m_{Z'[\Omega'\setminus W']}(P'\setminus W'))$ \src{Neu20}{2048}, canonized by Theorem~A under $\psi'(A)[\fP']$. This
yields a coset $E(Z')=\sigma A'''$, with $A'''\le\psi'(A)$ the preimage of the stabilizer of the canonical string, and we put
$C(Z):=(\psi')^{-1}(E(Z'))$, $A'':=(\psi')^{-1}(A''')$. For $\pi\in C(Z)$ put $Z'':=(Z^\pi)'=Z'^{\psi'(\pi)}$. Then
$x(Z'')$ is the canonical string, so the claim \texttt{respect-block-sizes} \src{Neu20}{2069--2074} holds for $A'''$ and $Z''$.

\emph{Step 4 (renormalize).} Apply \cite[Thm.~4.9]{Neu20} to $(A''',\fP',\vec\fB')$. The construction in its proof
\src{Neu20}{1455--1554} uses these data only; the strings enter only through the definition of $\fX^*$
\src{Neu20}{1521--1525}. So $\Omega^*,\fP^*,\vec\fB^*,f,f'$ and $\iota_0\colon A'''\to\Sym(\Omega^*)$ are functions of the prefix, and
$\iota:=\iota_0\circ\psi'|_{A''}$. Define $(Z^\pi)_*:=\{(P^*,\alpha^*\mapsto x(f(\alpha^*))):(f'(P^*),x)\in Z''\}$ (property (VII)) and
$W^*:=f^{-1}(W')$.

\emph{Verification.} (S1) is (I), (VI) of \cite[Thm.~4.9]{Neu20}, whose hypothesis is met by Step~2. (S2): invariance of $W^*$ follows from (II); $\Lambda^*[W^*]\le\Aut(F^*)$ from (IX),
applied to the relation $\psi'(A)[W']\le\Aut(Z'[W'])$ of Step 1. Since $Z''[W']$ is simple and equals the object
$Z'[W']$ of Step 1 (a function of $F$), and the strings of $(Z^\pi)_*$ on a block $P^*$ are $x\circ f$ with
$(f'(P^*),x)\in Z''$, the set $F^*$ is simple and determined by the computed data.
Complete occupancy of $(Z^\pi)_*$ is immediate from (VII) and the complete occupancy of $Z''$. (S3): $f'$ is
onto $\fP'$ (otherwise (VIII) fails for two sets differing on a block outside the image); each $P^*$ is mapped
bijectively onto $f'(P^*)$ by $f$ ((III), (V)), so $x\circ f$ determines $x$; equivariance is (VIII) applied to $\fX:=Z''$,
$\mathfrak Y:=Z''^{r}$, together with Step~1. (S4) is the estimate in the proof of properties (F), (G) of
\cite[Thm.~4.17]{Neu20} \src{Neu20}{2115--2141}, which uses only (I)--(V) and \texttt{respect-block-sizes}; it holds with
$\fX_1:=Z^\pi$, $\fX_1'=\fX_1'':=Z''$ and $\Gamma'':=A'''$; (S4) for $(Z^\pi)_*[\Omega^*\setminus W^*]$ is proved in the same way.

In the language of ECMs: $C$ is an $A$-ECM because $x(Z'^{\psi'(a)})=x(Z')^{\psi'(a)}$ (action on $\fP'$) and the Theorem~A
canonizer of $x(Z')$ is equivariant; and $Z\mapsto Z'$ and $Z''\mapsto(Z^\pi)_*$ are the changes of action of Lemma~E3 (using
the injectivity of Step~1 and (S3)) that transport canonizers back along $\psi'$ and $\iota_0$.
\end{proof}
Neuen runs steps 3 and 4 once per equivalence class of pairs, with a representative group $\Gamma''_1$. Here
step 3 is a canonizer and step 4 is computed from its canonical output.

\subsection{Canonical local certificates}

\begin{named}[Lemma LC]
Assume the hypotheses of \cite[Thm.~4.24]{Neu20}: $\Gamma\in\Gd$ on $\Omega$ with an almost $d$-ary sequence
containing $\fP$, a giant representation $\varphi\colon\Gamma\to S_k$ with $\max\{8,2+\log d\}<t\le k-2$, and
$T_0$, $H_0=\Gamma_{T_0}$, $\psi$ as in Lemma~T$^{\text{set}}$. Let $\cZ$ be the set of completely occupied sets of
$\fP$-strings of $d$-virtual size at most $s$. Algorithm CanLC below computes an $H_0$-ECM $C$ on $\cZ$ with right group $R$, and
an $H_0$-invariant flag, such that:
\begin{enumlab}[LC]{\arabic}
\item $C(Z^g)=g^{-1}C(Z)$ for $g\in H_0$;
\item if the flag is \emph{nonfull}, then $\psi(R(Z))\not\ge\Alt(T_0)$;
\item if the flag is \emph{full}, CanLC also outputs $\Delta(Z)\le\Aut_{H_0}(Z)$ with $\psi(\Delta(Z))\ge\Alt(T_0)$ and
  $\Delta(Z^g)=g^{-1}\Delta(Z)g$.
\end{enumlab}
Cost: polynomial work in $s$, polynomially many calls of Theorem~A on domains of size polynomial in $s$,
and recursive calls of the main canonizer on instances of virtual size at most $s/2$, of total virtual size
at most $t!\,s$.
\end{named}

\emph{State at round $i$} (for an object $Z$, with $\pi\in C_i(Z)$). $C_i$ is an $H_0$-ECM with right group $R_i$.
Its prefix $\bar p_i(Z)$ consists of a domain $\Omega_i$, a partition $\fP_i$ and a sequence $\fB^i$; a window
$W_i$; a group $\Lambda_i\le\Sym(\Omega_i)$ with an injective $\tau_i\colon\Lambda_i\to H_0$ such that $\tau_i(\Lambda_i)=R_i$;
the object map $Z'\mapsto Z'_*$; and the window part $F_i:=(Z^\pi)_*[W_i]$, the same for all $\pi\in C_i(Z)$.
Invariants:
\begin{enumlab}[I]{\arabic}
\item $\fB^i$ is an almost $d$-ary sequence of $\Lambda_i$-invariant partitions containing $\fP_i$;
\item $W_i$ is $\Lambda_i$-invariant, $\Lambda_i[W_i]\le\Aut(F_i)$, and $F_i$ is simple;
\item for $\lambda\in\Lambda_i$: $\lambda\in\Aut((Z^\pi)_*)$ iff $\tau_i(\lambda)\in\Aut(Z^\pi)$;
\item the virtual size of $(Z^\pi)_*$ is at most $s$, and the virtual size $s'_i$ of $(Z^\pi)_*[\Omega_i\setminus W_i]$
  strictly decreases with $i$;
\item $(Z^\pi)_*$ is completely occupied.
\end{enumlab}
Start: $C_0=H_0$, $\Omega_0=\Omega$, $W_0=\emptyset$, $\Lambda_0=H_0$, $\tau_0=\mathrm{id}$.

\emph{Round $i$} (each step is Lemma~E2 applied with the named ECM):
\begin{enumlab}{\alph}
\item If $\psi(R_i)\not\ge\Alt(T_0)$, stop with \emph{nonfull}.
\item Put $W':=\Aff(\Lambda_i,\psi\tau_i)\setminus W_i$, a function of the prefix. If $W'=\emptyset$, stop with \emph{full}
  and $\Delta(Z):=\pi\,\tau_i\big((\Lambda_i)_{(\Omega_i\setminus W_i)}\big)\,\pi^{-1}$.
\item \emph{Balance}: Lemma~B on $(Z^\pi)_*$ under $\Lambda_i$, giving the group $\Lambda'$. If $\psi\tau_i(\Lambda')$ is not
  $\ge\Alt(T_0)$, stop with \emph{nonfull}.
\item \emph{Window}: canonize $(Z^\pi)_*[W']$ under $\Lambda'$, joining it to the window part $F_i$ on $W_i$, by iterating
  Lemma~F2$^{\text{set}}$ from the prefix $(W_i,F_i,\Lambda')$ over the $\Lambda'$-orbits $O$ of $W'$ in (P1) order. For an
  orbit $O$, processed under the current right group $R\le\Lambda'$, the required $R$-canonizer of the object on $O$ is
  obtained by descent (F3) from $R$ to $R\cap N'$, where $N':=\ker(\psi\tau_i|_{\Lambda'})$ (index at most $t!$ in $R$), and
  by the subroutine Chain of \S\ref{sec:canhyp} under $(R\cap N')^O$ (intransitive on $O$) on the object in each
  branch; its windows are recursive calls of the main canonizer (CanHyp, at smaller virtual size).
\item \emph{Simplify and renormalize}: Lemma~SIMP at the new prefix (window $W_i\cup W'$). Lemma~E3 yields
  $\Lambda_{i+1}$, $\tau_{i+1}$, $\Omega_{i+1}$, $W_{i+1}$.
\end{enumlab}

\begin{proof}[Proof of Lemma LC]
\emph{(LC1).} By induction on rounds. Each step composes (Lemma~E2) an ECM that comes from a canonizer
(Theorem~A, recursive calls, Lemmas~B and~W), a descent (F3), or a change of action computed from the prefix
(Lemma~E3, Lemma~SIMP). The stopping tests depend only on $R_i$ and the prefix, which are
$H_0$-invariant.

\emph{Invariants.} (I1) and (I2) are (S1), (S2) of Lemma~SIMP. (I3) follows from (S3), the claim
\texttt{preserve-isomorphisms} and property (VIII) of \cite[Thm.~4.9]{Neu20}, and the fact that balancing
and windowing only pass to subgroups containing the relevant automorphisms (Lemma~E1). (I5) holds at the start by
hypothesis and is preserved by (S3); it is needed for strict monotonicity of the virtual size. (I4) is Neuen's
argument \src{Neu20}{2400--2412}: $W'\neq\emptyset$ strictly lowers $s'$ \cite[Obs.~4.12]{Neu20} (using (I5)), and Simplify does
not raise it (S4). Hence there are at most $s$ rounds.

\emph{(LC2)} holds by the tests at (a) and (c).

\emph{(LC3).} At (b), $\Aff(\Lambda_i,\psi\tau_i)\subseteq W_i$, $\Lambda_i$ is $d$-normalized by (I1), and $\psi\tau_i(\Lambda_i)$ is a
giant by test (a). So \cite[Thm.~4.6]{Neu20} (= \cite[Thm.~5.8]{GNS}) gives
$\psi\tau_i((\Lambda_i)_{(\Omega_i\setminus W_i)})\ge\Alt(T_0)$. \cite[Lemma~4.16]{Neu20} with (I2) gives
$(\Lambda_i)_{(\Omega_i\setminus W_i)}\le\Aut((Z^\pi)_*)$, and (I3) transfers this to $\tau_i(\dots)\le\Aut_{H_0}(Z^\pi)$.
Conjugating by $\pi$ gives $\Delta(Z)\le\Aut_{H_0}(Z)$ with a giant image. $\Delta(Z)$ does not depend on $\pi$, because
$(\Lambda_i)_{(\Omega_i\setminus W_i)}\trianglelefteq\Lambda_i$ ($W_i$ is $\Lambda_i$-invariant), so its $\tau_i$-image is normalized by
$R_i$. Equivariance follows from (LC1).

\emph{Cost.} Step (d) is Neuen's recursion for one pair \src{Neu20}{2436--2452}, with one change:
\cite[Lemma~4.7]{Neu20} (= \cite[Thm.~6(b)]{Bab16}) is applied to $\Lambda'$ itself, which is a giant by test (c);
the orbits in $W'$ are affected for $\Lambda'$ because $\Aff$ is antimonotone. So each $N'$-orbit in a $\Lambda'$-orbit
$O$ has size at most $|O|/t$, and the $(R\cap N')$-orbits lie inside $N'$-orbits. Since $(Z^\pi)_*[O]$ is balanced
(Lemma~B), each piece has virtual size at most $s/2$ \cite[Obs.~4.13]{Neu20}. Over the rounds the window
virtual sizes sum to at most $s$ \src{Neu20}{2425--2434}; with the $t!$ cosets the total is at most $t!\,s$.
The calls of Theorem~A (Lemmas~B, W and SIMP) are polynomially many, on domains of polynomial size.
\end{proof}

\emph{Comparison with Neu20.} Neuen computes, for every pair of test objects, the coset of isomorphisms,
and runs all $4k^{2t}$ pairs in one synchronized loop so that renormalizations are comparable
\src{Neu20}{270--275, 1930--1933}. CanLC runs $\binom kt$ independent single-object loops, one per $t$-set:
transport the $t$-set to $T_0$, then run the loop on $Z^\sigma$. Consistency between test sets is automatic,
because every auxiliary construction is a function of the canonical prefix, and Lemma~T$^{\text{set}}$(a)
makes the result independent of the transporter. The recursive cost drops from $4k^{2t}t!\,s$ to
$k^tt!\,s$.

\emph{Why the setwise stabilizer.} In Section~\ref{sec:agg} we canonize under $J_0=\ker\psi$ and take the
windows from the GNS groups $G_i(T,x)=\Aut_{H_T}(\win{x}{W_i})$, which are intrinsic to $(T,x)$. In Neu20 the
corresponding groups live on renormalized domains, and they depend on structure-graph choices made at
earlier rounds. They are intrinsic only once those choices are made canonically, which requires
canonizing at the level $H_0$ at which the windows are defined.

\subsection{Canonical aggregation}

\begin{named}[Lemma AGG]
In the setting of \cite[Lemma~4.29]{Neu20}, for a single $\fX$ (and $\Gamma'$, $\varphi$, $t<k/10$ as there), one
computes one of:
\begin{enumlab}{\roman}
\item a family $\{\fA_j(\fX)\}_{j\le k^6}$ of $t$-ary structures on subsets $V_j\subseteq[k]$ with $|V_j|\ge3k/4$, each
  of relative symmetry defect at least $1/4$, equivariant as a family:
  $\{\fA_j(\fX^h)\}=\{\fA_j(\fX)\}^{\varphi(h)}$;
\item a pair $(M(\fX),\Delta(\fX))$ with $|M|\ge3k/4$, $\Delta\le\Aut_{\Gamma'_M}(\fX)$ and $\varphi(\Delta)[M]\ge\Alt(M)$, both
  equivariant.
\end{enumlab}
In the case $|\supp_\fX|<k/4$ the structure has relative symmetry defect at least $3/4$.
\end{named}
\begin{proof}
Run CanLC on $\fX^\sigma$ for each $t$-set $T$. This gives the canonical fullness certificates
$\Delta(T,\fX):=\sigma\Delta(\fX^\sigma)\sigma^{-1}$ for routine-full $T$ (equivariant, since
$\Delta(\fX^{\sigma g})=g^{-1}\Delta(\fX^\sigma)g$), and the colours $c(\fX,\vT)$ of Lemma~T$^{\text{set}}$ for all orderings
of routine-nonfull $T$. By (LC2), (LC3) and Lemma~E1, ``routine-full'' coincides with Neuen's ``full''. The
fullness group $\Lambda_\fX:=\langle\Delta(T,\fX):T\text{ routine-full}\rangle$ and its support $\supp_\fX$ are equivariant. The
cases $|\supp_\fX|\in[k/4,3k/4]$ and $|\supp_\fX|>3k/4$ \src{Neu20}{2554--2582} produce GNS's structures from $\Lambda_\fX$ alone, and
each individualization is a member of an equivariant branch family of size at most $k^6$
(Proposition~\ref{prop:cases}). In the case $|\supp_\fX|<k/4$, put $D:=[k]\setminus\supp_\fX$ and
$\fA(\fX):=([k];\text{unary }D;\ \vT\mapsto\text{rank of }c(\fX,\vT)\text{ for }\vT\in D^{\langle t\rangle})$. Every
$T\subseteq D$ is routine-nonfull. $\fA$ is strongly canonical by Lemma~T$^{\text{set}}$(b), and the ranking is
canonical because the set of colour values is an invariant. If $\Alt(M)\le\Aut(\fA)$ with $|M|\ge t+2$ and
$M\subseteq D$, then for a $t$-set $T\subseteq M$ all orderings are induced by $\Alt(M)$ and have equal colour, so
Lemma~T$^{\text{set}}$(c) gives $\psi(R)=\Sym(T_0)$, contradicting (LC2). Hence every such $M$ lies in $\supp_\fX$ or has at
most $t+1$ points, and the relative defect is at least $1-\max(1/4,1/10+1/k)=3/4$ for $k\ge7$.
\end{proof}
Neuen uses the comparison classes only to build his structures \src{Neu20}{2583--2597}, and find-structure
uses only the invariance of the family and the defect bound \src{Neu20}{2610--2630}; exactness of the colour
classes is not needed.

\subsection{The algorithm CanHyp}\label{sec:canhyp}

This is the counterpart of Section~\ref{sec:algo} for sets of strings. \emph{Instances} are tuples
$(\Omega,\Gamma,\fP,\vec\fB,\fX)$, where $\Omega$ is ordered, $\Gamma\le\Sym(\Omega)$ is in $\Gd$ and given by its
canonical generating set, $\vec\fB$ is a $\Gamma$-invariant almost $d$-ary sequence with $\fP$ as a level, and $\fX$ is a
completely occupied set of $\fP$-strings. As in Section~\ref{sec:algo}, $d$ and $t:=\max\{9,3+\lceil\log d\rceil\}$ are fixed
globally, and $(\fP,\vec\fB)$ are auxiliary data: they are functions of the group data above the node, or of
canonical prefixes (Lemma~E2), and never depend on $\fX$ directly. The \emph{measure} of an instance is
\[ \mu(\Omega,\Gamma,\fP,\vec\fB,\fX):=(s,|\Gamma|)\in\mathbb N^2,\qquad s:=\text{the $d$-virtual size of }\fX, \]
ordered lexicographically; this order is a well-order. All recursive calls of CanHyp below are made on
instances of strictly smaller measure (Proposition~\ref{prop:canhyp}).

\paragraph{Chain.} For $A\le\Sym(\Omega)$ intransitive and $Z$ a completely occupied set of $\fP$-strings,
$\mathrm{Chain}(A,Z)$ (with $\Omega,\fP,\vec\fB$ understood from the context) is the iteration of Lemma~F2$^{\text{set}}$ over orbits of the right groups:
put $C_0(Z):=A$, $R_0:=A$, $U_0:=\emptyset$, $F_0:=\emptyset$. At stage $i$, if $U_i=\Omega$, stop and return $(F_i,C_i(Z))$.
Otherwise let $V$ be the first $R_i$-orbit in $\Omega\setminus U_i$ in the order (size, least point), let $\pi\in C_i(Z)$, and
call CanHyp on the transitive instance $(V,R_i^V,\fP[V],\vec\fB[V],Z^\pi[V])$ (a valid instance, by \cite[Obs.~4.2]{GNS}, and
the restriction of a completely occupied set is completely occupied). Its output, lifted through the restriction
$R_i\to R_i^V$ to an $R_i$-canonizer of $Z^\pi[V]$, is the input of Lemma~F2$^{\text{set}}$; this gives
$C_{i+1},R_{i+1},F_{i+1}$ and $U_{i+1}:=U_i\cup V$.

\paragraph{Luks.} For an $A$-invariant partition $\fD$ with at least two classes,
$\mathrm{Luks}(A,\fD,Z)$ is: apply Lemma~B to $(A,\fP,Z)$, obtaining $C_{\mathrm B}$ with right group $B$; let
$\pi\in C_{\mathrm B}(Z)$ and $y:=Z^\pi$. If $B$ is intransitive, let $(F,\Lambda):=\mathrm{Chain}(B,y)$. Otherwise descend
(Lemma~F3) from $B$ to $B_{(\fD)}$ over at most $|B^\fD|$ cosets and run $\mathrm{Chain}(B_{(\fD)},y^\rho)$ on each branch
($B_{(\fD)}$ is intransitive, as its orbits lie in classes of $\fD$), obtaining a $B$-canonizer $(F,\Lambda)$ of $y$. Return
$(F,\pi\Lambda)$, an $A$-canonizer by Lemma~E4.

\paragraph{Algorithm CanHyp.} On an instance $(\Omega,\Gamma,\fP,\vec\fB,\fX)$, with
$\vec\fB=(\{\Omega\}=\fB_0\succ\fB_1\succ\dots)$:
\begin{description}
\item[Step 0.] If $|\Omega|=1$, return $(\fX,\Gamma)$.
\item[Step 1. Intransitive $\Gamma$.] Return $\mathrm{Chain}(\Gamma,\fX)$.
\item[Step 2. Transitive $\Gamma$, $\Gamma^{\fB_1}$ semi-regular.] Return $\mathrm{Luks}(\Gamma,\fB_1,\fX)$.
\item[Step 3. Transitive $\Gamma$, $|\fB_1|\le d$.] Choose (P1) a block system $\fB\succeq\fB_1$ of $\Gamma$ with $|\fB|\ge2$ on which
  $P:=\Gamma^{\fB}$ is primitive, and put $b:=|\fB|\le d$.
  \begin{enumlab}{\alph}
  \item If $|P|\le b^{1+\log b}$, or $b$ is below the constant of \cite[Lemma~4.35]{Neu20}: return
    $\mathrm{Luks}(\Gamma,\fB,\fX)$.
  \item Otherwise apply \cite[Lemma~4.35]{Neu20} (= \cite[Lemma~7.1]{GNS}) to $P$ (P1), obtaining $N$, $\cC$, $\psi$ and $k$.
    Put $G':=\{g\in\Gamma:g^{\fB}\in N\}$, $\varphi:=\psi\circ(g\mapsto g^{\fB})$, and let $\Cp$ be the lift of $\cC$, so that
    $\ker\varphi=G'_{(\Cp)}$, $|\Cp|\ge2$ and $\Cp\succeq\fB\succeq\fB_1$. Descend (Lemma~F3) from $\Gamma$ to $G'$ over at most $b$
    cosets, and in each branch $\rho$ run $\mathrm{Agg}(G',\fX^\rho)$.
  \end{enumlab}
\item[Step 4. $\mathrm{Agg}(G',y)$] (with $\varphi,\Cp,\vec\fB$ from Step~3) returns a $G'$-canonizer of $y$.
  \begin{itemize}
  \item If $k\le10t$: return $\mathrm{Luks}(G',\Cp,y)$ (its descent is to $B_{(\Cp)}\le\ker\varphi$, of index
    at most $k!$).
  \item If $k>10t$: run CanLC for every $t$-set of $[k]$ (with recursive calls of CanHyp, Lemma~LC) and apply
    Lemma~AGG to $y$. In case (i) apply Proposition~\ref{prop:opt1} with the family of Lemma~AGG, each structure
    encoded on $[k]$ with its domain marked by a unary relation; for each member $\fA$ of the family, call
    CanHyp on $(\Omega,R_\fA,\fP,\vec\fB\oplus\Cp,y^{\sigma_\fA})$, and combine by Lemma~S and Lemma~F4. In case (ii) apply
    Lemma~\ref{lem:K} with $\Delta(y):=M(y)$ and $K(y):=\Delta(y)$ and call CanHyp on its two children
    $(\Omega,L,\fP,\vec\fB\oplus\Cp,\cdot)$.
  \end{itemize}
\item[Step 5. Top partition of same-domain children.] Put
  $\vec\fB\oplus\Cp:=(\{\Omega\}\succ\Cp\succ\fB_1\succ\fB_2\succ\cdots)$, with duplicates removed. As in Step~6 of
  \S\ref{sec:canad}, this is an almost $d$-ary sequence for every subgroup of $G'$, since $|\Cp|\le b\le d$ and each
  $\Cp$-class contains at most $|\fB_1|\le d$ blocks of $\fB_1$; and $\fP$ is still a level.
\end{description}

\begin{proposition}\label{prop:canhyp}
CanHyp terminates on every instance $(\Omega,\Gamma,\fP,\vec\fB,\fX)$, and it returns a $\Gamma$-canonizer for $\fX$
(on the completely occupied sets of $\fP$-strings): its output $(F,\Lambda)$ satisfies $\Lambda=\{g\in\Gamma:\fX^g=F\}\ne\emptyset$, and
$F$ depends only on $(\Omega,\Gamma,\fP,\vec\fB)$ (with $\Gamma$ taken as a set) and on the $\Gamma$-orbit of $\fX$.
\end{proposition}
\begin{proof}
\emph{Well-foundedness.} We show that every recursive call of CanHyp on an instance $I$, made directly or inside
Chain, Luks, Agg or CanLC, has strictly smaller measure than $I$.
\begin{enumerate}
\item \emph{Windows of Chain.} Chain is entered only for an intransitive group, so every window $V$ is a proper
  subset of the current domain (at stage $0$ because the group is intransitive, later because $U_i\ne\emptyset$ and
  $V\cap U_i=\emptyset$). Since $\fX$ is completely occupied and transported copies have the same virtual size,
  \cite[Obs.~4.12]{Neu20} gives that $s$ strictly decreases. (Chain is entered in Step~1, inside Luks, and in step (d)
  of CanLC; in the last case the objects are completely occupied by (I5).)
\item \emph{Calls inside CanLC.} They are the windows of Chain in step (d), on subsets of $W'\subseteq\Omega_i$ of
  $(Z^\pi)_*$, which has virtual size at most $s$ by (I4). Each is a proper subset of $\Omega_i$, because every
  $N'$-orbit in a $\Lambda'$-orbit $O$ has size at most $|O|/t$ (Cost paragraph of Lemma~LC). So by (I5)
  and \cite[Obs.~4.12]{Neu20} their virtual size is strictly smaller than $s$.
\item \emph{Same-domain children} in Step~4 are on proper subgroups: $R_\fA$ and $L$ have index at least
  $(4/3)^k>1$ in $G'\le\Gamma$ (Proposition~\ref{prop:opt1}, Lemma~\ref{lem:K}, \cite[Cor.~4.28]{Neu20}), and the virtual size is unchanged,
  since the object is only transported.
\end{enumerate}
Calls of Theorem~A, Corollary~D and BC are calls of other algorithms, which terminate. So CanHyp terminates, and we
may prove the rest by induction on $\mu$, assuming that the recursive calls of CanHyp return canonizers
(Lemma~LC and Lemma~AGG were stated and proved under this assumption).

\emph{Each step returns a canonizer.}
\begin{itemize}
\item Step~0 is trivial.
\item \emph{Chain.} By induction, Lemma~F2$^{\text{set}}$ applies at every stage, so each $C_i$ is an $A$-ECM with prefix
  $(U_i,F_i,R_i)$. At the end $U_r=\Omega$, so $C_r$ is an $A$-ECM with right group $R_r$ and $Z^\pi=F_r$ for all $\pi\in C_r(Z)$;
  thus $R_r\le\Aut(F_r)$. Lemma~E4, with the trivial $R_r$-canonizer on $\{F_r\}$, shows that $(F_r,C_r(Z))$ is an
  $A$-canonizer.
\item \emph{Luks} and Steps~2, 3(a): Lemma~B, Lemma~E4, Lemma~F3 and the previous item.
\item \emph{Step~3(b).} The objects $\fB,N,\psi,\Cp,G',\varphi$ are (P1) functions of $\Gamma$, and Lemma~F3 is applied to the
  $G'$-canonizers returned by Agg.
\item \emph{Step~4.} If $k\le10t$ this is Luks. If $k>10t$, Lemma~AGG gives the equivariant data needed: an equivariant
  family of structures of relative symmetry defect at least $1/4$, or an equivariant pair $(M(y),\Delta(y))$ with
  $\Delta(y)\le\Aut_{G'_{M(y)}}(y)$ and $\varphi(\Delta(y))[M(y)]\ge\Alt(M(y))$. By induction the children return canonizers, and
  Proposition~\ref{prop:opt1} (through Lemma~S and Lemma~F4), respectively Lemma~\ref{lem:K}, give a $G'$-canonizer.
  Step~5 keeps the instance valid.
\end{itemize}

\emph{Canonicity of auxiliary data.} Every auxiliary object is a (P1) function of the current group, or of a prefix
of an ECM (Lemma~E2), or an equivariant function of the current object (Lemmas~W, B, SIMP, LC, AGG). Hence, for
fixed $(\Omega,\Gamma,\fP,\vec\fB)$, the output depends only on the $\Gamma$-orbit of $\fX$.
\end{proof}

\subsection{Theorem H}

\begin{proof}[Proof of Theorem H (correctness)]
For a coset $\sigma\Gamma$, the coset rule gives $\Lambda_{\sigma\Gamma}(\fX)=\sigma\Lambda_\Gamma(\fX^\sigma)$, and $\fX^\sigma$ is a set of
$\fP^\sigma$-strings, which is why $\fP^\sigma$ must be $\Gamma$-invariant; so let $\sigma=1$. Let $\fX^+$ be $\fX$ together with
the constant string $\bot$ (a new letter) on every block of $\fP$ that carries no string of $\fX$; the map $\fX\mapsto\fX^+$
is injective and equivariant and $\fX^+$ is completely occupied. Compute (P1), from $(\Gamma,\fP)$ alone, the
normalization of \cite[Thm.~4.8]{Neu20}: a set $\Omega^*$, a monomorphism $\iota\colon\Gamma\to\Sym(\Omega^*)$, an almost $d$-ary
sequence $\vec\fB^*$ of $\iota(\Gamma)$-invariant partitions, and a level $\fP^*$ of it. The map $\fX^+\mapsto(\fX^+)^*$ of that theorem is injective,
since $f$ maps each $P^*$ onto a block of $\fP$, every block is hit, and $x^*=x\circ f$; it is equivariant by the proof
there, and $(\fX^+)^*$ is completely occupied. Run CanHyp on $(\Omega^*,\iota(\Gamma),\fP^*,\vec\fB^*,(\fX^+)^*)$; by
Proposition~\ref{prop:canhyp} it returns an $\iota(\Gamma)$-canonizer. Lemma~R (with $\iota$ and the map $\fX\mapsto(\fX^+)^*$) gives a
$\Gamma$-canonizer for $\fX$.
\end{proof}

\begin{proofsketch}[running time of Theorem H]
The recursion tree has the same node types as Neuen's, and the canonical version adds no branch type: the
descent to $G'$ ($\le d$), Luks steps ($\le k!$, or $b$ cosets for a semi-regular top), the $t!$ cosets in
LC, at most $k^6$ individualizations and at most $k^6$ structures in Option~1, and two cosets in Option~2
all occur in Neuen's test with the same counts. LC runs $\binom kt$ single-object loops instead of $4k^{2t}$
pair loops. Isomorphism calls of GNS become calls of Theorem~A on domains of size at most
$s\le(n+m)^{O(\log d)}$, costing $(n+m)^{O((\log d)^{c+1})}$, which is absorbed into the exponent. Babai's test on
$t$-ary structures over at most $k$ points becomes BC, at cost $\exp(\polylog d)$. Neuen's
orbit-by-orbit processing is sequential too \src{Neu20}{1677--1686}, so the chain rule has the same
recursion, and the potential argument \src{Neu20}{2710--2725} only needs upper bounds on $|\Lambda[\Cp]|$,
which pass to subgroups. With the choice of Remark~\ref{rem:gap1}, \cite[Lemma~4.36]{Neu20} should give
$\sum s_i\le2^{O(\log^3d)}s$ at a node with $s_i\le s/2$, and \cite[Lemma~2.6]{Neu20} then gives
$s^{O(\log^3d)}(n+m)^{\polylog d}$. With $s\le(n+m)^{O(\log d)}$ this is $(n+m)^{O((\log d)^{c'})}$. What is
missing is a written transfer of Neuen's accounting \cite[Lemma~4.36, Lemma~2.6]{Neu20} to the canonical
recursion, including the two points of Remark~\ref{rem:neu20}.
\end{proofsketch}
\status{Correctness (algorithm CanHyp, Proposition~\ref{prop:canhyp}): proof modulo Theorem~A; \cite[Thms.~4.8, 4.9, 4.17]{Neu20} (constructions and the listed
properties), \cite[Lemma~4.16]{Neu20}, the virtual-size facts \cite[Obs.~4.12, 4.13, Lemma~4.14]{Neu20};
\cite[Thm.~5.8]{GNS}; \cite[Thm.~6(b)]{Bab16}; the CFSG facts of the aggregation \cite[Thm.~4.31,
Lemma~4.32]{Neu20}; \cite[Lemma~7.1]{GNS}; BC; canonical generating sets; and requirement~\ref{D3}, which
includes the structure-graph lemma of Neuen's dissertation \cite[Lemma~6.1.17]{Neu19}, which we have not
seen. Running time: \textbf{proof sketch}.}

\begin{remark}[Two points about Neu20 as written]\label{rem:neu20}
\begin{enumerate}
\item The potential argument in the proof of \cite[Lemma~4.36]{Neu20} contains the step discussed in
  Remark~\ref{rem:gap1} \src{Neu20}{2716}, and the choice of block system \src{Neu20}{2683} is equally
  unconstrained. The same choice applies: $\Cp\succeq\fB\succeq\fB_1$, and inserting $\Cp$ at the top keeps the
  sequence almost $d$-ary for $\Gamma'$, since $|\Cp|\le|\fB|\le d$.
\item \emph{Order of the giant test (cost accounting only).} In the local certificates algorithm the
  recursive call on $W'$ \src{Neu20}{2222} precedes the test ``$(\Lambda')^\psi\not\ge\Alt(T)$''
  \src{Neu20}{2224}. The cost proof bounds kernel orbits using the kernel in the \emph{pre-balancing} group
  $\Lambda$, which is a giant \src{Neu20}{2442}, but bounds virtual size using balance, which holds for orbits of
  the \emph{post-balancing} group $\Lambda'$. If $\psi(\Lambda')$ is not a giant, kernel orbits of $\Lambda'$ need not be
  small, and the bound $s_i\le s/2$ is not justified as written. Testing $\psi(\Lambda')$ immediately after
  balancing, as in step (c) of CanLC, fixes this; correctness of Neu20 is unaffected. Whether the journal
  version handles this we do not know.
\end{enumerate}
There is also a harmless typo: ``$\ker\psi=N_{(\mathfrak B)}$'' \src{Neu20}{2687} should read $N_{(\mathfrak C)}$.
\end{remark}

% =====================================================================
\section{Treewidth: Theorem TW (proof sketch)}\label{sec:TW}

This section is a \textbf{proof sketch}. It rewrites Wiebking's $n^{\polylog k}$ isomorphism test for graphs of
treewidth $k$ \cite[Thm.~35]{Wie20} canonically, within the Schweitzer--Wiebking framework. The key
observation is that the canonical version needs only Theorem~A (through Corollary~D), and not the
canonical version of Neu20.

\subsection{What is known}

\begin{center}\small
\begin{tabularx}{\linewidth}{@{}L{4.3cm}YL{2.2cm}@{}}
\toprule
result & statement & canonical?\\
\midrule
\cite[Thm.~6.13]{Wie21} (also \cite{GNSW}) & treewidth $k$, $2^{O(k^2\log k)}n^{O(1)}$ & yes\\
\cite{LPPS17} & treewidth $k$, $2^{O(k^5\log k)}n^5$ & yes\\
\cite{ES16} & bounded treewidth, logspace for each fixed $k$ & yes\\
\cite{GNSW} (journal version) & treewidth $k$, $2^{k\polylog k}n^{O(1)}$ & no \cite[pp.~9--10]{Wie21}\\
\cite[Thm.~35]{Wie20} $=$ \cite[Cor.~6.30]{Wie21} & treewidth $k$, $n^{\polylog k}$ & no\\
\cite[Thm.~6.29]{Wie21}, \cite{GNW} & excluded $K_h$ minor, $n^{\polylog h}$ (uses \cite{Neu20}) & no \cite[p.~179]{Wie21}\\
BC \cite{Bab19} & all graphs, $n^{\polylog n}$ & yes\\
\bottomrule
\end{tabularx}
\end{center}
GNS conjectured an $n^{O((\log k)^c)}$ isomorphism test for treewidth $k$ \src{GNS}{2270--2274}, which
Wiebking proved. Wiebking's outlook predicts that canonization should follow once GNS is canonical
(\S\ref{sec:open}).

\subsection{What Wiebking's algorithm imports}

Wiebking's algorithm \cite[proof of Thm.~35]{Wie20} approximates the treewidth, $k$-improves the graph
\cite[Lemma~29]{Wie20} (from \cite{LPPS17}), computes the clique-separator decomposition
\cite[Thm.~30]{Wie20} (from \cite{Lei93,ES16}), and then runs a bottom-up procedure \textsc{isoTree}.

\begin{center}\small
\begin{tabularx}{\linewidth}{@{}L{4cm}L{3.2cm}Y@{}}
\toprule
component & non-canonical element & canonical replacement\\
\midrule
treewidth approximation \cite{Amir01} & output need not be invariant & invariant doubling search with invariant checks (fix TW-3)\\
$k$-improvement $G^k$ & --- & invariant \cite[Lemma~29]{Wie20}\\
clique-separator decomposition & --- & output is isomorphism-invariant \cite[Thm.~30]{Wie20}\\
root of the decomposition tree & arbitrary root & centre of the tree (subdivide a central edge)\\
children: isomorphisms for all pairs & pairwise & canonical labelings of children (recursion)\\
``$\rho_1$ arbitrary labeling'', ``representative $t^*$'' & choice & colour $=$ canonical form of the child, label $=$ its canonical labeling coset restricted to the adhesion set \cite[\S6.1.3]{Wie21}\\
equal adhesion sets: \texttt{canObj} on the bag & none (bag $\le k+1$) & \cite[Cor.~27]{Wie20}, time $n^{\polylog k}$\\
distinct adhesion sets: \textsc{isoBasicClique} & GNS calls, arbitrary choices & rewrite (b) below, with fix TW-1\\
distinct adhesion sets: \texttt{cosetHyper} & Neuen-type isomorphism & Lemma~CH: Corollary~D only\\
distinct adhesion sets: GNS on the bag & GNS isomorphism & Corollary~D\\
extension to subtrees & arbitrary extensions of generators & canonical extension \cite[pp.~145--152]{Wie21}\\
\textsc{isoBasic} & ``some $S_1$, $S_2$'', ``some bijection'', GNS calls & rewrite (a) below\\
parent adhesion set & not used in \textsc{isoTree} & marked in the bag object (fix TW-2)\\
\bottomrule
\end{tabularx}
\end{center}

\subsection{Coset-labeled hyperedges under a hyperedge-respecting coset}

Wiebking's \cite[Lemma~34]{Wie20} (\texttt{cosetHyper}; Wiebking cites it as Lemma~9 of the ICALP version
of~\cite{GNSW} and refers to~\cite{Neu19}; that number does not match the arXiv~v1 numbering we use for
\cite{GNSW})
assumes that the given coset ``maps the hyperedges in $H_1$ to hyperedges in $H_2$''. Its only use, in the
proof of \cite[Thm.~35]{Wie20}, comes right after the isomorphisms of the bag with its hyperedges have been
computed by \cite[Lemma~33]{Wie20}, so the hypothesis holds there. Only the coset labels remain. (An
earlier version of our notes listed \texttt{cosetHyper} as an obstacle requiring canonical Neu20; that was
wrong for this application.)

\begin{named}[Lemma CH]
Let $\Delta\rho$ be a labeling coset of $V$ with $\Delta\le\Sym(V)$, $\Delta\in\Gd$. Let $\cS\subseteq2^V$ and $\chi\colon\cS\to\mathbb N$
be preserved by $\Delta$. Let each $S\in\cS$ carry a labeling coset $\tau_S\Theta_S$ of $S$, with
$\tau_S\colon S\to[|S|]$, $\Theta_S\le\Sym([|S|])$, and $\Theta_S=\Theta_{S'}$ whenever $\chi(S)=\chi(S')$. Let
$d':=\max_S\cw(\Theta_S)\le\max_S|S|$. Then a canonical labeling coset of $X=(V,\cS,\chi,\{\tau_S\Theta_S\})$ relative
to $\Delta\rho$ can be computed in time $(|V|+\sum_S|S|)^{O((\log\max(d,d'))^c)}$, with $c$ as in Theorem~A.
\end{named}
Note the two conventions in this statement: in $\Delta\rho$ the group sits on the side of $V$ (as in
Schweitzer and Wiebking), whereas in $\tau_S\Theta_S$ it sits on the label side (as in \cite{GNSW}); both are
consistent with their sources.
\begin{proof}
This is the second case of \cite[Lemma~23]{GNSW} (arXiv~v1 numbering; ``we can assume that $\Delta$
preserves $\mathcal S$ and $\chi$''). Let $U:=\bigsqcup_S\{S\}\times S$ and let $\hat\Lambda$ be the labeling coset of
$V\cup U$ that orders $V$ by $\delta\rho\in\Delta\rho$, orders the blocks $\{S\}\times S$ by the order of $S^{\delta\rho}$, and
orders inside each block by some element of $\tau_S\Theta_S$. Since $\Delta$ preserves $\chi$, the block of $S$ moves to
the block of $S^\delta$ with inner permutation in $\Theta_S$, so $\hat\Lambda$ is a coset; its group is an extension of
$\prod_S\Theta_S$ by $\Delta$, of composition width at most $\max(\cw\Delta,d')$. Let $Y$ be the incidence graph
$v$---$(S,v)$, and return $\mathrm{CL}(Y;\hat\Lambda)|_V$. GNSW prove the required properties of this output; they
canonize $Y$ with a Luks-type algorithm, and we use Corollary~D instead.
\end{proof}
\status{Proof modulo Corollary~D and \cite[Lemma~23]{GNSW}.}

\subsection{Canonical rewrites of Wiebking's isomorphism-only steps}

Everything below takes place in the Schweitzer--Wiebking framework: an isomorphism-invariant labeling
coset is refined step by step, and $\mathrm{CL}(X_2;\mathrm{CL}(X_1;\Lambda))$ is a canonical labeling of $(X_1,X_2)$
relative to $\Lambda$. Every step uses only Corollary~D, (P1) constructions, Theorem~A through a homomorphism,
or Wiebking's quasipolynomial canonizers on at most $k+1$ points. Each item is a sketch.

\paragraph{(a) \textsc{isoBasic} \cite[Lemma~32]{Wie20}.} Input: a clique-separator-free graph $G$ (here a
bag $G^k[D]$, with separability at most $k$ and degeneracy at most $k$ by the checks of~(c)).
\begin{itemize}
\item \emph{Branching.} Branch over $S\in\cS:=\{N(v):\deg v=\min\deg(G)\}$, an invariant family, by Lemma~F4. In
  a branch, start from the invariant labeling coset $\Lab(S)$.
\item \emph{Separator sizes.} Canonize the complete directed graph on $S$, coloured by the sizes of the leftmost
  minimal separators $|S_{v,w}|$, relative to the current coset (a string on $S^2$; Corollary~D).
\item \emph{Wreath product.} Form $\hat S:=S\sqcup\bigsqcup_{v,w}\hat S_{v,w}$ and the coset $\hat\Lambda$ that orders $S$ by
  the current labeling, the blocks by $(v^\lambda,w^\lambda)$, and the inside of each block arbitrarily. This is an
  invariant construction, with composition width at most $\max(|S|,\mathrm{sep})\le k$.
\item \emph{Identification.} Canonize the identification graph relative to $\hat\Lambda$ (Corollary~D), and pass to
  the induced labeling coset of the quotient, ordering classes by least label.
\item Repeat until $S=V$ (at most $n$ rounds), and finally canonize $G$ relative to the coset (Corollary~D).
\end{itemize}
Each $\hat\Lambda$ orders $S$ first, so its group stabilizes $S$; the final group is $\Aut(G,S)$, and Lemma~F4 over the
invariant family $\cS$ gives $\Aut(G)$.

\paragraph{(b) \textsc{isoBasicClique} \cite[Lemma~33]{Wie20}, with fix TW-1.} Input: $G$ as in (a) and a
set $\mathcal H$ of cliques of $G$, possibly coloured by $\chi\colon\mathcal H\to\mathbb N$.
\begin{itemize}
\item Keep the branch $S$ of (a), so that all cosets have composition width at most $k$.
\item The clique cover $\alpha$ of \src{Wie20}{185--207} (peel the vertices of minimum degree; $\alpha(v)=N(v)\cup\{v\}$
  at the moment $v$ is peeled) is invariant, and $|\alpha(v)|\le\mathrm{degeneracy}(G)+1\le k+1$.
\item For each $v$, canonize the hypergraph $\mathcal H_v:=\{S'\in\mathcal H:S'\subseteq\alpha(v)\}$ on $\alpha(v)$, with the
  colours $\chi$, by Wiebking's hypergraph canonization \cite[Thm.~9]{Wie20} in time
  $(k+1+|\mathcal H_v|)^{\polylog(k+1)}$.
\item Colour $v$ by the canonical form of $\mathcal H_v$, and canonize relative to the current coset
  (Corollary~D).
\item Form the wreath product over $U:=V\sqcup\bigsqcup_v\alpha(v)\times\{v\}$ from the canonical cosets of the
  $\mathcal H_v$, canonize the identification graph (Corollary~D), and induce the coset on $V$.
\end{itemize}
The output group preserves $\mathcal H$: identification forces the restriction to $\alpha(v)$ into the canonical coset of
$\mathcal H_v\to\mathcal H_{v^\varphi}$, and every hyperedge lies in some $\alpha(v)$.
\emph{Fix TW-1.} Wiebking's version never needs $\chi$, because in the distinct-adhesion case the label
image encodes the class. With compressed labels it does not, and Lemma~CH requires $\chi$ to be preserved.
So $\chi$ is passed into the hypergraph canonizations as above or, equivalently, after (b) the string
$\chi\colon\mathcal H\to\mathbb N$ is canonized under the induced action of the output group on $\mathcal H$ (Theorem~A
through a homomorphism; $\Gd$ is closed under quotients). The resulting group preserves $\mathcal H$ and~$\chi$.

\paragraph{(c) \textsc{isoTree}, with fixes TW-2 and TW-3.} An instance is
$\hat X=(\hat G,T,\beta,r,S)$, where $S$ is the adhesion set to the parent ($S=\emptyset$ at the global root). Let
$D:=\beta(r)$, and let the children $t$ have adhesion sets $S_t:=D\cap\beta(t)$.
\begin{itemize}
\item \emph{Children.} Canonical labelings by recursion. A child's colour is its canonical form (with its
  adhesion set marked); its label is its canonical labeling coset restricted to $S_t$ and compressed to
  $\tau_t\Theta_t$. Equal colours give equal $\Theta_t$, as Lemma~CH needs.
\item \emph{Fix TW-2 (parent adhesion set).} The root-bag object must contain $S$; otherwise the
  canonical labeling of $\hat X$ need not map $S$ consistently and the parent's step fails. In the
  distinct-adhesion case, add $S$ as a hyperedge with a reserved colour. ($S$ is a clique in $G^k$, so the
  clique cover covers it, and by \cite[Thm.~30(3)]{Wie20} it differs from the children's adhesion sets.) In
  the equal-adhesion case $S$ equals every child's adhesion set and is marked as such.
\item \emph{Equal adhesion sets.} The bag has at most $k+1$ vertices (check (i) below). Apply Wiebking's
  canonization of objects with multiple labeling cosets \cite[Cor.~27]{Wie20} to the bag graph, the child
  labels and colours, and the marked $S$, in time $n^{\polylog k}$. (The $|D|^{O(\cw)}$ route of
  \cite[Thm.~6.11]{Wie21} costs $2^{O(k\log k)}$ and must not be used when $k\gg\log n$.) Extend to the
  subtrees as in the equal-adhesion cases of the proof of \cite[Thm.~6.11]{Wie21} (pp.~148--151), bundling
  children with identical colour and label.
\item \emph{Distinct adhesion sets.} Branch over the minimum-degree neighbourhoods of $Y:=G^k[D]$ (Lemma~F4).
  Let $\Lambda_1$ be the output of (b) on $(Y,\{S_t\}\cup\{S\})$ with colours, including the $\chi$-step of TW-1. Let
  $\Lambda_2$ be the output of Lemma~CH relative to $\Lambda_1$, with colours and labels from the children. Let $\Lambda_3$
  be the canonical labeling coset of $G[D]$ relative to $\Lambda_2$ (Corollary~D). Take the minimum over the
  branches, and extend canonically to the subtrees by the ``partition into singletons'' case of the proof of
  \cite[Thm.~6.11]{Wie21} (pp.~151--152), which accepts any canonical coset on $D$ whose group respects the
  labeled hyperedges; $\Lambda_3$ is such a coset.
\item \emph{Root.} Use the centre of the decomposition tree. A subdivided central edge gives a bag that is a
  clique in $G^k$ with two equal adhesion sets of size at most $k+1$. Components are sorted by canonical
  form.
\item \emph{Fix TW-3 (the parameter).} For $k=1,2,4,\dots$ compute $G^k$ and its clique-separator
  decomposition, and accept the first $k$ for which:
  (i) every equal-adhesion bag has at most $k+1$ vertices;
  (ii) every basic bag $G^k[D]$ has degeneracy at most $k$;
  (iii) every adhesion set has at most $k+1$ vertices.
  These checks are invariant and polynomial. Canonicity holds for every $k$; the checks only certify the
  time bound. They hold for every $k\ge\tw(G)$: by \cite[Lemma~29(3)]{Wie20}, $\tw(G^k)=\tw(G)\le k$, so
  $\mathrm{degeneracy}(G^k[D])\le\tw(G^k[D])\le k$, equal-adhesion bags have at most $\tw(G^k)+1\le k+1$
  vertices by \cite[Thm.~30(3)]{Wie20}, and adhesion sets are cliques of $G^k$. So the accepted $k$ is at most
  $\max(1,2\tw(G))$. Check (ii) uses degeneracy rather than minimum degree: the time analysis needs
  $|\alpha(v)|\le k+1$ for the clique cover in (b), which Wiebking derives from ``minimum degree $\le\tw$'' for
  every peeled subgraph, i.e.\ from degeneracy. For $k<\tw(G)$, $\tw(G^k)$ can be much larger than
  $\tw(G)$, so a minimum-degree check would not certify that bound. Separability at most $k$ needs no
  check, since local connectivity is monotone under induced subgraphs and $\mathrm{sep}(G^k)\le k$.
\end{itemize}

\subsection{Theorem TW}

\begin{proofsketch}[Theorem TW]
Assemble (a)--(c). Each node of the clique-separator decomposition makes polynomially many calls of
Corollary~D with composition width at most $k+1$ on domains of polynomial size, polynomially many calls of
Theorem~A through homomorphisms (TW-1), and calls of \cite[Thm.~9, Cor.~27]{Wie20} on at most $k+1$
points. There are polynomially many nodes and at most $n$ branches per bag. The total is
$n^{O((\log k)^{c''})}$ with $k\le\max(1,2\tw(G))$.
\end{proofsketch}
\status{\textbf{Proof sketch}, conditional on Theorem~A (via Corollary~D); \cite[Lemma~29, Thm.~30, Thm.~9,
Cor.~27]{Wie20}; the canonical-extension construction in the proof of \cite[Thm.~6.11]{Wie21}; and
\cite[Lemma~23]{GNSW}. Fixes TW-1, TW-2, TW-3 come from an internal review. The rewrites (a)--(c) have been
checked against Wiebking's text step by step, but they are not written as full proofs.}

Theorem~TW is incomparable with \cite[Thm.~6.13]{Wie21} ($2^{O(k^2\log k)}n^{O(1)}$); it is faster once
$k\ge(\log n)^{1/2+\varepsilon}$. For $k=\polylog(n)$ it gives $n^{(\log\log n)^{O(1)}}$, against BC's
$n^{\polylog n}$. A second route, through the minor result (as in \cite[Cor.~6.30]{Wie21}) and Corollary~H4
below, would need Theorem~H; the route above needs less.

% =====================================================================
\section{Further corollaries}\label{sec:further}

Every statement in this section is conditional or a sketch, as marked.

\begin{named}[Corollary H1 (coset-labeled hypergraphs)]
Given a coset-labeled hypergraph $(V,\mathcal E,\mathfrak p)$ in the sense of \cite{GNW} (labeling cosets of the
hyperedges, and colours) with $|E|\le d$ for all $E\in\mathcal E$, and a labeling coset $\Delta\rho$ of $V$ with $\Delta\in\Gd$,
a canonical labeling coset can be computed in time $(n+m)^{O((\log d)^{c'})}$, with $c'$ as in Theorem~H.
\end{named}
\begin{proof}
Refine each edge colour $c$ to $(c,\Theta)$; this is invariant, since isomorphisms in the sense of \cite{GNW}
require equal $\Theta$. First canonize the coloured hypergraph $(V,\mathcal E,\chi)$ relative to $\Delta\rho$ by Theorem~H;
its right group then preserves $\mathcal E$ and $\chi$. Then apply Lemma~CH with $d'\le d$.
\end{proof}
\status{Proof modulo Theorem~H and Corollary~D; hence its running time is only as good as the sketch for
Theorem~H. The isomorphism version in \cite{GNW} cites \cite{Neu20} and \cite[Thm.~6.6.7]{Neu19}.}

\begin{named}[Corollary H2 ($t$-CR-bounded graphs; sketch)]
Canonical labelings of $t$-CR-bounded graphs and pairs (\cite[Def.~5.1, Thm.~5.3, Cor.~5.4]{Neu20}) can be
computed in time $n^{\polylog t}$; and for the closure $D=\mathrm{cl}_t(X)$ of \cite{GNW} one can compute an
invariant labeling coset of $D$ of composition width at most $t$ containing all isomorphisms restricted to $D$.
\end{named}
\begin{proofsketch}
Neuen's proof maintains a coset on the colour classes and refines it, each refinement being an isomorphism
of quotient graphs (GNS) or of sets of strings (Neu20). Canonically, one maintains an invariant labeling
coset: small classes are split by a wreath product with $\Sym(P)$, $|P|\le t$, ordered by the parent label,
followed by Corollary~D on the quotient graph; refinement rounds canonize the set of strings relative to the
current coset (Theorem~H), induce a labeling coset on the new classes by the lexicographic order of the
transported strings, and canonize the quotient graph (Corollary~D). The composition width stays at most $t$,
and there are polynomially many steps. The closure output is a weak canonical labeling (an equivariant coset
containing the automorphisms), not a canonical one.
\end{proofsketch}
\status{\textbf{Proof sketch}, modulo Theorem~H and Corollary~D.}

\begin{named}[Corollary H3 ($K_{3,h}$-minor-free graphs, Euler genus; sketch)]
Canonical forms of $K_{3,h}$-minor-free graphs can be computed in time $n^{\polylog h}$, and of graphs of
Euler genus $g$ in time $n^{\polylog g}$.
\end{named}
\begin{proofsketch}
Canonical forms of 3-connected (vertex- and arc-coloured) components combine into canonical forms of the
graph by the standard canonical decomposition into 2- and 3-connected components (not re-derived here). For
a 3-connected component, branch over ordered triples $(v_1,v_2,v_3)$ ($n^3$ branches, Lemma~F4);
\cite[Cor.~6.3]{Neu20}, which holds for any three vertices, makes the result $(h-1)$-CR-bounded, and H2 applies.
By \cite[Thm.~6.1]{Neu20} (Ringel), graphs of Euler genus $g$ exclude $K_{3,4g+3}$. To obtain a bound in terms of
the unknown genus, run an invariant doubling search over $h$, using the invariant $(h-1)$-CR test.
\end{proofsketch}
\status{\textbf{Proof sketch}, modulo H2.}

\begin{named}[Corollary H4 (excluded minors and topological subgraphs; conditional)]
Canonical forms of graphs excluding $K_h$ as a minor, or as a topological subgraph, can be computed in time
$n^{\polylog h}$.
\end{named}
\begin{proofsketch}
The non-canonical elements of the proof of \cite[main theorem]{GNW} are: ``fix one representative'',
``an arbitrary bijection'' (twice), ``$v_1\in X_1$ arbitrary'', and calls to the $t$-CR closure theorem, the
coset-labeled hypergraph theorem, Wiebking's canonization for multiple labeling cosets (canonical already),
and Babai's test on fewer than $h$ vertices. Canonically: colours and labels come from canonical labelings of
the parts as in \cite[\S6.1.3]{Wie21}; equal separators are handled by Wiebking's multiple-labeling-coset
canonization on fewer than $h$ points; distinct separators by H1 relative to the coset from H2; branching
over $v\in X$ by Lemma~F4; and BC on fewer than $h$ vertices. The decomposition steps of \cite{GNW} are
invariant constructions and are used as black boxes. The proof of \cite{Neu22t} has the same structure; its
extra call computes a set stabilizer, which Theorem~A handles.
\end{proofsketch}
\status{\textbf{Conditional} on Theorem~H, H1 and H2; the canonical dynamic programming is only
\textbf{sketched}. To answer Question~9.7 of \cite{Wie21} (parameter the Hadwiger number of the input), one
also needs a parameter-free form: the canonical counterpart of \cite[main theorem]{GNW} must have the form
``either a $K_h$ minor is certified, or a canonical form is output'', with the failure outcome produced from
canonical data so that it is invariant, followed by an invariant doubling over $h$. We have not written this
out.}

\paragraph{Prior art for H4.} Canonization of graphs excluding a fixed minor is known: in polynomial time for
each fixed excluded minor by Grohe \cite{Gro12,Gro17}, and in time $f(H)\cdot n^{O(1)}$ by Lokshtanov,
Pilipczuk, Pilipczuk and Saurabh \cite{LPPS22}. Neither gives a bound of the form $n^{\polylog h}$, so
Question~9.7 remains open as stated, and H4, if completed, would be new only in its dependence on $h$.

\paragraph{Hypergraphs under $\Sym(V)$.} Theorem~H with $d=n$ gives $(n+m)^{\polylog n}$, which is Wiebking's
hypergraph canonization \cite[Thm.~5.4(1)]{Wie21} again; nothing new.

\paragraph{FPT treewidth canonization (speculation).} The FPT isomorphism test of \cite{GNSW} in time
$2^{k\polylog k}n^{O(1)}$ uses GNS only for strings under subgroups of automorphism groups of trees, which lie in
$\Gd$, so Theorem~A covers those calls; its other non-canonical step, coset intersection for multiple coloured
cosets on $O(k^2)$ points, has a canonical replacement by Wiebking. So the obstacle named in
\cite[p.~10]{Wie21} may be removed. We have not traced the rest of that algorithm, and we make no claim.

% =====================================================================
\section{Assumptions from the literature}\label{sec:assumptions}

\paragraph{Versions.} GNS: arXiv:1802.04659v3. Neu20: arXiv:2002.06997v2. Wie20: arXiv:1911.11257 (the
version in our possession). GNSW: arXiv:1803.06858v1. SW19: arXiv:1806.07466. BC: the STOC~2019 paper.
Wie21: the dissertation, printed page numbers. Theorem numbers for arXiv papers were computed from their
\LaTeX{} sources by counting theorem environments; please treat them as pointers and check them against the
labels quoted.

\begin{center}\small
\begin{longtable}{@{}rL{9.4cm}L{4.6cm}@{}}
\toprule
\# & result & used for\\
\midrule
\endhead
1 & \cite[Thm.~2.1]{GNS}: polynomial-time permutation group algorithms (Seress), in deterministic versions \ref{D1} & throughout\\
2 & \cite[Lemma~2.4]{GNS}: $\Gd$ closed under subgroups and quotients; closure under extensions (Jordan--Hölder) & \S\ref{sec:prelim}, \S\ref{sec:luks}\\
3 & \cite[Thm.~3.13]{GNS}: structure of large primitive $\Gd$-groups, via Thm.~4.3 & change of action\\
4 & \cite[Thms.~4.3, 4.6]{GNS}, \emph{including the constructions in their proofs}: $\iota$ a homomorphism; $x\mapsto x^*$ injective and equivariant; the bound on $|\Omega^*|$; the almost $d$-ary sequence; computability; all choices deterministic from canonical generators \ref{D3}. Also \cite[Lemma~4.4]{GNS} (from \cite{BLS87}) and \cite[Lemma~4.5]{GNS} & Lemma~R\\
5 & \cite[Obs.~4.2]{GNS} & inheritance of sequences\\
6 & \cite[Thm.~5.8]{GNS} (= \cite[Thm.~4.6]{Neu20}) & \ref{L3}, Lemma~\ref{lem:51}, Lemma~LC\\
7 & \cite[Thm.~5.9]{GNS} = \cite[Thm.~6(b)]{Bab16} (= \cite[Lemma~4.7]{Neu20}) & Lemmas~\ref{lem:51}, \ref{lem:Tpp}, LC\\
8 & \cite[Lemma~6.2]{GNS} and the termination argument of \cite[Lemma~6.3]{GNS} & \ref{L1}--\ref{L5}\\
9 & \cite[Cor.~6.5]{GNS} (from \cite[Thm.~5.2A,B]{DM96}) & index $\ge(4/3)^k$\\
10 & \cite[Thm.~A.2]{GNS} (CFSG: non-giants are at most 5-transitive); \cite[Lemma~A.3]{GNS} = \cite[Cor.~2.4.13]{Bab15} & case (iv)\\
11 & \cite[Lemma~7.1]{GNS} = \cite[Thm.~3.2.1]{Bab15}, deterministic \ref{D2}; with $|\cC|\ge2$, $\cC$ an equipartition, and $k\le d$ & Step~4, \S\ref{sec:potential}\\
12 & \cite[Lemma~2.6]{GNS} (an inequality) & \S\ref{sec:time}\\
13 & \cite[Thm.~2.2]{Bab19} (canonical form of strings under any group, quasipolynomial; a preliminary report that sketches the proof), called only for $S_k$, $k\le d$, on relational structures of arity $O(\log d)$ on $[k]$ (domains of size $\le2k^t=d^{O(\log d)}$; \ref{D4}, Remark~\ref{rem:BC}); and Babai's string isomorphism test \cite[Thm.~2.1]{Bab19} & Prop.~\ref{prop:opt1}\\
14 & Canonical generating sets: \cite[Lemma~6.2]{AGMMM18} = \cite[Lemma~21]{GNSW} = \cite[Lemma~15]{SW19} (arXiv numbering) & (P1)\\
15 & Lexicographically least element of a coset in polynomial time (standard; point-stabilizer chain) & (P1)\\
\midrule
\multicolumn{3}{@{}l}{\emph{Additionally for Theorem H}}\\
16 & \cite[Thm.~4.8]{Neu20} (normalize), \cite[Thm.~4.9]{Neu20} (renormalize), \cite[Thm.~4.17]{Neu20} (simplify on window), with their constructions and listed properties, and the claims in their proofs; deterministic \ref{D3} & Lemma~SIMP, Theorem~H\\
17 & \cite[Lemma~6.1.17]{Neu19} (structure graphs), as used by \cite{Neu20}; \textbf{not seen by us} & normalization\\
18 & \cite[Lemma~4.16]{Neu20} (extending window automorphisms) & Lemma~LC\\
19 & \cite[Obs.~4.12, 4.13, Lemma~4.14]{Neu20} (virtual size); \cite[Lemma~4.36, Lemma~2.6]{Neu20} for the time sketch & Lemma~LC, Theorem~H\\
20 & \cite[Thm.~4.31, Lemma~4.32, Cor.~4.28]{Neu20} (as items 9, 10) & Lemma~AGG\\
\midrule
\multicolumn{3}{@{}l}{\emph{Additionally for Theorem TW and the corollaries}}\\
21 & \cite[Lemma~29]{Wie20} (from \cite{LPPS17}), \cite[Thm.~30]{Wie20} (from \cite{Lei93,ES16}) & decomposition\\
22 & \cite[Thm.~9]{Wie20} (hypergraph canonization), \cite[Cor.~27]{Wie20} (multiple labeling cosets) & Theorem~TW\\
23 & \cite[Lemma~23]{GNSW} (arXiv v1), second case & Lemma~CH\\
24 & The canonical-extension construction in the proof of \cite[Thm.~6.11]{Wie21} (pp.~145--152) & Theorem~TW\\
25 & \cite[Def.~5.1, Thm.~5.3, Cor.~5.4, Thm.~6.1, Cor.~6.3]{Neu20}; the structure theorems of \cite{GNW,Neu22t} as invariant constructions & H2--H4\\
\bottomrule
\end{longtable}
\end{center}

\paragraph{Not assumed; reproved here.} The Babai--Luks framework lemmas F1--F4; the canonical bounded-degree
reduction of \cite{BL83} (Section~\ref{sec:luks}), including \cite[Lemma~4.3]{BL83}; the recursion lemma
(Lemma~REC); the potential argument (Lemma~\ref{lem:potential}); the canonical versions of \cite[Lemmas~4.3, 4.15]{Neu20}
(Lemmas~B and~W, which we prove directly).

\paragraph{New here, as far as we know.} Lemma~T and its use in Definition~\ref{def:Ax} and
Lemmas~\ref{lem:Tprime}, \ref{lem:Tpp}; Lemmas~T$^{\text{set}}$, LC, SIMP and AGG (the single-object local
certificates loop for Neu20); the observation that Wiebking's treewidth algorithm needs only canonical GNS;
Lemma~\ref{lem:K} and Lemma~R as stated (routine); the choice of Remark~\ref{rem:gap1}.

% =====================================================================
\section{Related work}\label{sec:related}

\begin{itemize}
\item \emph{Bounded degree.} Luks \cite{Luks82}; Babai, Kantor and Luks \cite{BKL83} ($n^{O(d/\log d)}$);
  canonical forms by Babai and Luks \cite{BL83} and by Fürer, Schnyder and Specker \cite{FSS83}. \cite{BL83}
  introduces canonical labeling cosets, canonical forms relative to a coset, and the half-edge trick.
\item \emph{Quasipolynomial time.} Babai's isomorphism test \cite{Bab16,Bab15} and canonical form
  \cite{Bab19}. BC makes the local-certificate structure strongly canonical \cite[Thm.~12.1]{Bab19} by
  different means (a structure built from the certificates, and a selection in asymmetric test sets,
  Claim~12.4), which need $t\ge2\log n$.
\item \emph{$n^{\polylog}$ in a parameter.} GNS \cite{GNS} (maximum degree); Neuen \cite{Neu20} (hypergraphs
  under $\Gd$, $t$-CR-bounded graphs, genus); Wiebking \cite{Wie20} (treewidth); Grohe, Neuen and Wiebking
  \cite{GNW} (excluded minors); Neuen \cite{Neu22t} (excluded topological subgraphs). All are isomorphism
  tests. Neuen's survey \cite{Neu26} states that canonization ``is possible for many of the algorithms
  discussed'', without a specific statement for these results. See also the survey \cite{GN21}.
\item \emph{Canonization frameworks.} Schweitzer and Wiebking \cite{SW19}, Wiebking \cite{Wie20,Wie21}:
  labeling cosets, hypergraphs, multiple labeling cosets, canonical FPT for treewidth \cite[Thm.~6.13]{Wie21},
  and DeepWL \cite{GSW21} (which turns polynomial-time DeepWL isomorphism tests into canonization; GNS is
  group-theoretic and not of this form). Wiebking's dissertation lists our problem as Question~9.6 and
  contains the canonical test set device \cite[Lemma~5.34]{Wie21}.
\item \emph{Treewidth.} Canonical FPT algorithms by Lokshtanov, Pilipczuk, Pilipczuk and Saurabh
  \cite{LPPS17} and by Grohe, Neuen, Schweitzer and Wiebking \cite{GNSW,Wie21}; logspace canonization for
  bounded treewidth by Elberfeld and Schweitzer \cite{ES16}.
\item \emph{Excluded minors and rank width.} Canonization for each fixed excluded minor in polynomial time
  by Grohe \cite{Gro12,Gro17}, and in time $f(H)\,n^{O(1)}$ by Lokshtanov, Pilipczuk, Pilipczuk and Saurabh
  \cite{LPPS22}. Canonization for bounded rank width by Grohe and Neuen \cite{GN19}, and in parallel by
  Levet, Rombach and Sieger \cite{LRS23}. Neuen also gave an FPT isomorphism test parameterized by genus
  \cite{Neu24genus}.
\item \emph{Practical canonizers.} Individualization--refinement tools (e.g.\ \cite{MP14}) routinely relabel
  individualized vertices to reference labels. Lemma~T is the group-theoretic analogue at a fixed ordered
  tuple.
\item \emph{Prior-art searches.} Searches of the arXiv API, Semantic Scholar, OpenAlex and the web through
  5~November 2026 found no $n^{\polylog d}$ canonization for maximum degree $d$, no canonical version of
  \cite{Neu20}, and no $n^{\polylog k}$ canonization for treewidth $k$. We did not have access to the SIAM
  journal version of GNS or to Neuen's dissertation \cite{Neu19}. Since the dissertation predates Wiebking's
  2021 statement of Question~9.6, it is unlikely to contain a solution.
\end{itemize}

% =====================================================================
\section{Status and open points}\label{sec:statusopen}

\subsection{Status of every statement}

\begin{center}\small
\begin{tabularx}{\linewidth}{@{}L{4.6cm}Y@{}}
\toprule
statement & status\\
\midrule
Lemmas F1--F4, R, S, T, \ref{lem:K}; Prop.~\ref{prop:cases}, \ref{prop:opt1}; Lemmas~\ref{lem:51}, \ref{lem:Tprime}, \ref{lem:Tpp} & proofs given\\
Proposition~\ref{prop:correct} (correctness of CanAD) & proof modulo \S\ref{sec:assumptions}\\
Lemma~REC, Lemma~\ref{lem:potential} & proofs given\\
Theorem~A & proof modulo \S\ref{sec:assumptions}\\
Corollary~B & proof modulo Theorem~A; small computer check of the construction\\
Corollary~D & proof modulo Theorem~A\\
Lemmas E1--E4, T$^{\text{set}}$ & proofs given\\
Lemmas W, F2$^{\text{set}}$, B, SIMP, LC, AGG & proofs modulo Theorem~A and the results of Neu20 listed in \S\ref{sec:assumptions}\\
Proposition~\ref{prop:canhyp} (correctness of CanHyp) & proof modulo \S\ref{sec:assumptions}, including the unseen \cite[Lemma~6.1.17]{Neu19}\\
Theorem~H & correctness: proof modulo \S\ref{sec:assumptions} (Proposition~\ref{prop:canhyp}); running time: \textbf{proof sketch}\\
Lemma~CH & proof modulo Corollary~D and \cite[Lemma~23]{GNSW}\\
Theorem~TW & \textbf{proof sketch}, conditional on Theorem~A, with fixes TW-1--TW-3\\
H1 & proof modulo Theorem~H\\
H2, H3 & \textbf{proof sketches}\\
H4 & \textbf{conditional} on Theorem~H, H1, H2; \textbf{sketch}; parameter-free form not written\\
FPT treewidth canonization & speculation; no claim\\
Remark~\ref{rem:gap1} & question about the written proof of \cite[Lemma~7.2]{GNS}; the proposed choice is proved to work (Lemma~\ref{lem:potential})\\
Remark~\ref{rem:gap2} & presentation point in \cite[Thm.~7.3]{GNS}; covered by Lemma~REC\\
Remark~\ref{rem:neu20} & two points in the cost analysis of \cite{Neu20}; correctness of Neu20 unaffected\\
\bottomrule
\end{tabularx}
\end{center}

\subsection{How this draft was checked}

Three internal adversarial reviews were carried out (all AI-assisted, none of them refereeing).
\begin{itemize}
\item The first attacked the canonical colour construction (coset versus group transport, reference tuples,
  canonical generating sets), the exactness of colour classes, the completeness of the list of choice points
  (Table~\ref{tab:choices}), the running time, the Luks reduction, prior art, and the question why the problem
  was considered open. It found no break; its corrections (explicit membership tags; credit to
  \cite[Lemma~5.34]{Wie21}; writing out the running time and Corollary~B) are incorporated.
\item The second checked the running time and the canonical Luks reduction line by line, with a computer
  check of Section~\ref{sec:luks} on small graphs, and found four undocumented assumptions (determinism and
  canonical generating sets for all subroutines feeding (P1) choices, including BC; closure of $\Gd$ under
  extensions; $|\cC|\ge2$; $k\le d$), a slip in the description of the right subgroup in Lemma~F1 and
  Lemma~\ref{lem:51}, and imprecise wording for the semi-regular step and the chain rule. All are corrected
  here.
\item The third checked Theorem~H, Theorem~TW and the corollaries. It found the repairs TW-1--TW-3, the
  descent group in step (d) of CanLC, the coset convention for $\fP$ in Theorem~H, the need for a
  parameter-free form of H4, and the missing prior art \cite{Gro12,Gro17,LPPS22}. All are incorporated.
\end{itemize}
A separate check of Remark~\ref{rem:gap1} reproduced the example by computer for small $m$.
The written-out algorithm CanHyp of \S\ref{sec:canhyp}, Lemma~E4, Lemma~F2$^{\text{set}}$ and the proofs of Lemmas~W, B and
SIMP were added after these reviews, in response to an internal referee report, and have not been independently
checked.

\subsection{Questions for experts}

\begin{question}\label{q:T}
(Lemmas~\ref{lem:Tprime}, \ref{lem:Tpp}.) Is the colour of an ordered tuple, computed under $J_0=H_{(\vT_0)}$
after transport, a single well-defined $J_0$-canonizer on the layered certificate strings? Does ``all
orderings of $T$ have the same colour'' force $\psi_T(G(T,x))=\Sym(T)$? Do the colours cost at most $k^tn$ in
recursive sizes, with no relation between $t$ and $n$?
\end{question}

\begin{question}\label{q:changeofaction}
(Lemma~R; rows 1--5 of Table~\ref{tab:choices}.) Can every auxiliary choice in the proofs of
\cite[Thms.~4.3, 4.6]{GNS} be made as a deterministic function of the group, in particular the
$\mathcal P$-sequence of \cite[Thm.~3.13]{GNS} and the Johnson identifications? Is $\Omega^*$ ordered canonically,
and is the restriction to one $G^*$-orbit still injective on strings? We checked the relevant claims in the
proofs but did not re-verify \cite[Thm.~3.13]{GNS} or the group theory behind \cite[Lemma~4.5]{GNS}.
\end{question}

\begin{question}\label{q:BC}
Is \cite[Thm.~2.2]{Bab19} accepted as proved? BC is a preliminary report that sketches its proof; Theorem~A
uses it only for $S_k$, $k\le d$, on relational structures of arity $O(\log d)$ (Remark~\ref{rem:BC}).
\end{question}

Question~\ref{q:gap1} (GNS Lemma 7.2) should be checked first.

\begin{question}
(Neuen.) Is the ``all pairs in parallel'' device of \cite{Neu20} needed for any reason other than making
renormalization consistent? Our single-object loop at the setwise stabilizer of $T_0$ claims to make it
unnecessary. Is the order of the giant test (Remark~\ref{rem:neu20}(2)) handled in the journal version?
\end{question}

\begin{question}
(Wiebking.) Is it right that the treewidth algorithm of \cite{Wie20} needs only GNS, with \texttt{cosetHyper}
applied only to hyperedge-respecting cosets, so that its canonical version follows from canonical GNS plus
\cite[Thm.~9, Cor.~27]{Wie20} on at most $k+1$ points?
\end{question}

\begin{question}
Do the experts know of unpublished work on Question~9.6 of \cite{Wie21}? The device is theirs.
\end{question}

\subsection{Known limitations}

\begin{itemize}
\item We have not seen the SIAM journal version of GNS, Neuen's dissertation \cite{Neu19}, or the journal
  versions of \cite{Neu20} and \cite{GNSW}.
\item BC is a preliminary report that sketches its proof. Theorem~A depends on it only for structures of
  arity $O(\log d)$ on at most $d$ points (Remark~\ref{rem:BC}); Wiebking's work uses it in the same way.
\item We did not re-verify the group theory in \cite[Thm.~3.13, Lemma~4.5]{GNS}, or the internals of
  \cite[Thms.~4.8, 4.9, 4.17]{Neu20} beyond the facts used.
\item The running time of Theorem~H, Theorem~TW and H2--H4 are sketches. The correctness proof of
  Theorem~H (Proposition~\ref{prop:canhyp}) is written out, but it relies on the structure-graph lemma of
  \cite{Neu19}, which we have not seen, and on the constructions of \cite{Neu20} as listed in
  \S\ref{sec:assumptions}.
\item No formal verification was attempted, and the computer checks cover only small instances of
  Section~\ref{sec:luks} and of the example in Remark~\ref{rem:gap1}.
\end{itemize}

\section*{Acknowledgements and disclosure}

[To be completed by the author(s).] This draft was prepared with substantial assistance from an AI system
(Claude, Anthropic), including the derivations, the internal reviews and the small computer checks. It has
not been refereed. The canonical test set device is due to Daniel Wiebking, and the overall approach follows
Grohe, Neuen and Schweitzer, Babai and Luks, and Babai.

\bibliographystyle{alpha}
\bibliography{refs}

\end{document}
